% Contraction : Rédiger un résumé — Français Tle Tech — Cours signé OnBuch+
% Programme officiel MINESEC. Compiler avec : tectonic N04-contraction-rediger-resume.tex
\newcommand{\DOCMATIERE}{Français}
\newcommand{\DOCNIVEAU}{Tle Tech}
\newcommand{\DOCMODULE}{Français – classe de Terminale technique}
\newcommand{\DOCLECON}{N04}
\newcommand{\DOCTITRE}{Contraction : Rédiger un résumé}
% OnBuch+ — préambule « cours signés » ENRICHI (compilé avec tectonic / XeLaTeX)
% Système visuel « Pop » d'OnBuch+ : fond crème (jamais blanc), bordures encre
% épaisses, ombres « dures » décalées, titres Archivo Black en capitales.
% Version MATHÉMATIQUES : graphes (pgfplots/tikz), boîtes pédagogiques
% (définition, propriété, méthode, attention, le savais-tu, exercice résolu,
% exercice, TP, bilan), page de garde et signature OnBuch+ ancrée en bas.
%
% Macros attendues dans main.tex AVANT \input{preamble} :
%   \DOCMATIERE  \DOCNIVEAU  \DOCMODULE  \DOCLECON (numéro)  \DOCTITRE
\documentclass[11pt]{article}

\usepackage{fontspec}
\usepackage[french]{babel}
\usepackage[a4paper,margin=2cm,top=3.4cm,bottom=2.8cm,headheight=1.4cm,footskip=1.3cm]{geometry}
\usepackage{amsmath,amssymb,amsfonts}
\usepackage{mathtools} % \xmapsto, \coloneqq... souvent utilisés par les modèles
\usepackage{cancel} % simplifications barrées (bilans de forces)
% \abs{z} (module, valeur absolue) et \norm{u} (norme), utilisés par les modèles
\providecommand{\abs}{}\let\abs\relax\DeclarePairedDelimiter{\abs}{\lvert}{\rvert}
\providecommand{\norm}{}\let\norm\relax\DeclarePairedDelimiter{\norm}{\lVert}{\rVert}
\usepackage[table]{xcolor} % [table] : fournit aussi \rowcolors (alternance de lignes)
\usepackage{pagecolor}
\usepackage{tikz}
\usetikzlibrary{arrows.meta,positioning,calc,shapes.geometric,shapes.misc,shapes.arrows,decorations.pathmorphing,decorations.pathreplacing,decorations.markings,patterns,shadows,mindmap,trees,fit,backgrounds,matrix,intersections,angles,quotes}
\usepackage{pgfplots}
\usepackage[version=4]{mhchem} % équations nucléaires \ce{^{235}_{92}U}
\usepackage[european,siunitx]{circuitikz} % schémas électriques
\pgfplotsset{compat=1.18}
\usepgfplotslibrary{fillbetween,groupplots} % aires entre courbes (calcul intégral)
\usepackage{enumitem}
\usepackage{fancyhdr}
% [most] charge toutes les bibliothèques tcolorbox (listings, minted, raster,
% documentation...) et gonfle inutilement la mémoire TeX sur un document long
% (~300 boîtes) : on ne charge que ce qui est réellement utilisé.
\usepackage[skins,breakable]{tcolorbox}
\usepackage{graphicx}
\usepackage{colortbl}
\usepackage{booktabs}
\usepackage{tabularx}
\usepackage{multirow}
\usepackage{diagbox} % case d'en-tête diagonale des tableaux à double entrée
% Colonnes extensibles centrée / gauche / droite (C, L, R), souvent utilisées par les modèles
\newcolumntype{C}{>{\centering\arraybackslash}X}
\newcolumntype{L}{>{\raggedright\arraybackslash}X}
\newcolumntype{R}{>{\raggedleft\arraybackslash}X}
\usepackage{array}
\usepackage{multicol}
\usepackage{siunitx}
% parse-numbers=false : accepte aussi \SI{2,5 \times 10^{-3}}{...}
\sisetup{locale=FR,output-decimal-marker={,},group-minimum-digits=4,parse-numbers=false}
\DeclareSIUnit\ohm{\ensuremath{\symup{\Omega}}} % U+2126 absent de la police
\DeclareSIUnit\division{div} % réglages d'oscilloscope (ms/div, V/div)
\DeclareSIUnit\tr{tr} % tours (vitesse de rotation, ex. \SI{1500}{\tr\per\minute})
\DeclareSIUnit\molar{M} % concentration molaire (ex. \SI{3}{\molar}, \SI{1}{\micro\molar})
\DeclareSIUnit\an{an} % année (le modèle confond parfois avec \year, un compteur TeX) — ex. \SI{5}{\kilo\meter\per\an}
\DeclareSIUnit\FCFA{FCFA} % franc CFA (ex. \SI{500}{\FCFA\per\kilo\gram})
\DeclareSIUnit\degreeBrix{\textdegree Brix} % teneur en sucre d'un sirop/jus (réfractométrie)
\DeclareSIUnit\month{mois} % le modèle utilise parfois \month, un compteur TeX, comme unité de durée
\DeclareSIUnit\microliter{\micro\liter} % le modèle écrit parfois \microliter en un seul token
\DeclareSIUnit\million{millions} % dénombrement (ex. \SI{15}{\million\per\mL} spermatozoïdes)
\usepackage{qrcode}
% \degree : commande fréquemment utilisée par les modèles pour un angle ou une
% température en dehors de \SI{}{} (non fournie par les paquets chargés).
\providecommand{\degree}{^{\circ}}
% \qed (fin de démonstration), utilisé par les modèles sans amsthm
\providecommand{\qed}{\hfill$\square$}
% \barre : double barre des valeurs interdites dans un tableau de variations (tkz-tab)
\providecommand{\barre}{\ensuremath{\|}}
% environnement proof (amsthm non chargé) utilisé par les modèles
\ifdefined\proof\else\newenvironment{proof}[1][Démonstration]{\par\noindent\textit{#1.}\ }{\qed\par}\fi
\usepackage{lastpage}
\usepackage{hyperref}
\hypersetup{hidelinks}

% ============================================================
% Palette « Pop » OnBuch+ — mêmes hex que la classe P (plus_ui.dart)
% ============================================================
\definecolor{poporange}{HTML}{F59321}
\definecolor{popdark}{HTML}{E07A0C}
\definecolor{popcream}{HTML}{FFF6E8}
\definecolor{popink}{HTML}{1C1714}
\definecolor{poppurple}{HTML}{7A5AE0}
\definecolor{popgreen}{HTML}{2E9E6A}
\definecolor{poppink}{HTML}{E0457B}
\definecolor{popblue}{HTML}{2F7DE1}
\definecolor{popgold}{HTML}{C98A12}
\definecolor{popmuted}{HTML}{8A7F76}
\definecolor{popline}{HTML}{EADFD2}
\definecolor{popgray}{HTML}{8A7F76}
\colorlet{popgrey}{popgray}
% Alias tolérants (le modèle nomme parfois les couleurs différemment)
\colorlet{popwhite}{white}
\colorlet{popblack}{popink}
\colorlet{popred}{poppink}
\colorlet{popyellow}{popgold}
% Teintes claires (fonds de tableaux/encadrés) — le modèle nomme parfois
% les couleurs avec un suffixe L (ex. popblueL) sans les avoir définies.
\colorlet{poporangeL}{poporange!20}
\colorlet{popdarkL}{popdark!20}
\colorlet{poppurpleL}{poppurple!20}
\colorlet{popgreenL}{popgreen!20}
\colorlet{poppinkL}{poppink!20}
\colorlet{popblueL}{popblue!20}
\colorlet{popgoldL}{popgold!20}
\colorlet{popredL}{popred!20}
\colorlet{popgrayL}{popgray!20}
% Teintes claires pour fonds de tableaux / zones
\colorlet{popblueL}{popblue!10!white}
\colorlet{poporangeL}{poporange!14!white}
\colorlet{popgreenL}{popgreen!12!white}
\colorlet{poppurpleL}{poppurple!12!white}
\colorlet{poppinkL}{poppink!12!white}
\colorlet{popgoldL}{popgold!14!white}

\pagecolor{popcream}

% ============================================================
% Typographie
% ============================================================
\defaultfontfeatures{Ligatures=TeX}
\setmainfont{PlusJakartaSans-Medium.ttf}[Path=../../../fonts/,BoldFont=PlusJakartaSans-Bold.ttf]
% Maths en sans-serif (Fira Math) avec chiffres et lettres droites de Plus
% Jakarta, pour que les formules se fondent dans le texte.
\usepackage[math-style=ISO,bold-style=ISO]{unicode-math}
\setmathfont{FiraMath-Regular.otf}[Path=../../../fonts/]
\setmathfont{PlusJakartaSans-Medium.ttf}[Path=../../../fonts/,range={up/num,up/{Latin,latin},\mathup}]
\usepackage{icomma} % virgule décimale sans espace en mode math
% Caractères mathématiques Unicode absents des polices de texte (Plus Jakarta,
% Archivo Black) : rendus par leur équivalent mathématique, en texte comme en
% formule — sinon ils s'affichent en carré vide (ex. « Arithmétique dans ℤ »).
\usepackage{newunicodechar}
\newunicodechar{ℝ}{\ensuremath{\mathbb{R}}}
\newunicodechar{ℤ}{\ensuremath{\mathbb{Z}}}
\newunicodechar{ℂ}{\ensuremath{\mathbb{C}}}
\newunicodechar{ℕ}{\ensuremath{\mathbb{N}}}
\newunicodechar{ℚ}{\ensuremath{\mathbb{Q}}}
\newunicodechar{𝔻}{\ensuremath{\mathbb{D}}}
\newunicodechar{α}{\ensuremath{\alpha}}
\newunicodechar{β}{\ensuremath{\beta}}
\newunicodechar{γ}{\ensuremath{\gamma}}
\newunicodechar{δ}{\ensuremath{\delta}}
\newunicodechar{ε}{\ensuremath{\varepsilon}}
\newunicodechar{θ}{\ensuremath{\theta}}
\newunicodechar{λ}{\ensuremath{\lambda}}
\newunicodechar{μ}{\ensuremath{\mu}}
\newunicodechar{σ}{\ensuremath{\sigma}}
\newunicodechar{τ}{\ensuremath{\tau}}
\newunicodechar{φ}{\ensuremath{\varphi}}
\newunicodechar{ψ}{\ensuremath{\psi}}
\newunicodechar{ω}{\ensuremath{\omega}}
\newunicodechar{Σ}{\ensuremath{\Sigma}}
% majuscules produites par \MakeUppercase dans les titres (α-aminés -> Α)
\newunicodechar{Α}{\ensuremath{\alpha}}
\newunicodechar{Β}{\ensuremath{\beta}}
\newunicodechar{Γ}{\ensuremath{\Gamma}}
\newunicodechar{Δ}{\ensuremath{\Delta}}
\newunicodechar{Ε}{\ensuremath{\varepsilon}}
\newunicodechar{Θ}{\ensuremath{\Theta}}
\newunicodechar{Λ}{\ensuremath{\Lambda}}
\newunicodechar{Μ}{\ensuremath{\mu}}
\newunicodechar{Τ}{\ensuremath{\tau}}
\newunicodechar{Φ}{\ensuremath{\Phi}}
\newunicodechar{Ψ}{\ensuremath{\Psi}}
\newunicodechar{Ω}{\ensuremath{\Omega}}
\newunicodechar{↦}{\ensuremath{\mapsto}}
\newunicodechar{∈}{\ensuremath{\in}}
\newunicodechar{∉}{\ensuremath{\notin}}
\newunicodechar{∧}{\ensuremath{\wedge}}
\newunicodechar{⇔}{\ensuremath{\Leftrightarrow}}
\newunicodechar{⟺}{\ensuremath{\Longleftrightarrow}}
\newunicodechar{⇒}{\ensuremath{\Rightarrow}}
\newunicodechar{⟹}{\ensuremath{\Longrightarrow}}
\newunicodechar{∘}{\ensuremath{\circ}}
\newunicodechar{∪}{\ensuremath{\cup}}
\newunicodechar{∩}{\ensuremath{\cap}}
\newunicodechar{≡}{\ensuremath{\equiv}}
\newunicodechar{⊂}{\ensuremath{\subset}}
\newunicodechar{⊥}{\ensuremath{\perp}}
\newunicodechar{∃}{\ensuremath{\exists}}
\newunicodechar{∀}{\ensuremath{\forall}}
\newunicodechar{∛}{\ensuremath{\sqrt[3]{\,}}}
\newunicodechar{∜}{\ensuremath{\sqrt[4]{\,}}}
\newunicodechar{⃗}{\ensuremath{{}^{\to}}}
\newunicodechar{ⁿ}{\ifmmode^{n}\else\textsuperscript{\ensuremath{n}}\fi}
\newunicodechar{ˣ}{\ifmmode^{x}\else\textsuperscript{\ensuremath{x}}\fi}
\newunicodechar{ᵇ}{\ifmmode^{b}\else\textsuperscript{\ensuremath{b}}\fi}
\newunicodechar{ʳ}{\ifmmode^{r}\else\textsuperscript{\ensuremath{r}}\fi}
\newunicodechar{ᵘ}{\ifmmode^{u}\else\textsuperscript{\ensuremath{u}}\fi}
\newunicodechar{ʸ}{\ifmmode^{y}\else\textsuperscript{\ensuremath{y}}\fi}
\newunicodechar{ᵃ}{\ifmmode^{a}\else\textsuperscript{\ensuremath{a}}\fi}
\newunicodechar{ᵉ}{\ifmmode^{e}\else\textsuperscript{\ensuremath{e}}\fi}
\newunicodechar{ᵏ}{\ifmmode^{k}\else\textsuperscript{\ensuremath{k}}\fi}
\newunicodechar{ᵖ}{\ifmmode^{p}\else\textsuperscript{\ensuremath{p}}\fi}
\newunicodechar{ᴺ}{\ifmmode^{N}\else\textsuperscript{\ensuremath{N}}\fi}
\newunicodechar{ᶜ}{\ifmmode^{c}\else\textsuperscript{\ensuremath{c}}\fi}
\newunicodechar{ʰ}{\ifmmode^{h}\else\textsuperscript{\ensuremath{h}}\fi}
\newunicodechar{ᵠ}{\ifmmode^{q}\else\textsuperscript{\ensuremath{q}}\fi}
\newunicodechar{ₙ}{\ifmmode_{n}\else\textsubscript{\ensuremath{n}}\fi}
\newunicodechar{ₐ}{\ifmmode_{a}\else\textsubscript{\ensuremath{a}}\fi}
\newunicodechar{ᵢ}{\ifmmode_{i}\else\textsubscript{\ensuremath{i}}\fi}
\newunicodechar{ᵣ}{\ifmmode_{r}\else\textsubscript{\ensuremath{r}}\fi}
\newunicodechar{ₖ}{\ifmmode_{k}\else\textsubscript{\ensuremath{k}}\fi}
\newunicodechar{ᵦ}{\ifmmode_{\beta}\else\textsubscript{\ensuremath{\beta}}\fi}
\newunicodechar{ₓ}{\ifmmode_{x}\else\textsubscript{\ensuremath{x}}\fi}
\newfontfamily\popdisplay{ArchivoBlack-Regular.ttf}[Path=../../../fonts/]
\newfontfamily\poplogo{SpaceGrotesk-Bold.ttf}[Path=../../../fonts/]
\newfontfamily\popsemi{PlusJakartaSans-SemiBold.ttf}[Path=../../../fonts/]
\newfontfamily\popxbold{PlusJakartaSans-ExtraBold.ttf}[Path=../../../fonts/]
\newfontfamily\popxboldls{PlusJakartaSans-ExtraBold.ttf}[Path=../../../fonts/,LetterSpace=3.0]

\setlist[enumerate]{leftmargin=1.7em,itemsep=5pt,topsep=4pt}
\setlist[itemize]{leftmargin=1.7em,itemsep=4pt,topsep=4pt,label={\color{poporange}\textbullet}}
\setlength{\parindent}{0pt}
\setlength{\parskip}{5pt}
\renewcommand{\arraystretch}{1.25}

% pgfplots : style Pop par défaut
\pgfplotsset{
  every axis/.append style={axis line style={popink,line width=1pt},tick style={popink},
    grid style={popline,line width=0.5pt},label style={font=\small},tick label style={font=\footnotesize}},
  popcurve/.style={poporange,line width=1.8pt,no markers,smooth},
}
% Même style disponible dans un \draw TikZ simple (hors pgfplots)
\tikzset{popcurve/.style={poporange,line width=1.8pt}}
\tikzset{popgrid/.style={popline,very thin}}
\colorlet{popcurve}{poporange} % le nom du style est parfois utilisé comme couleur

% ============================================================
% Pill / squiggle
% ============================================================
\newcommand{\poppill}[3]{%
  \tikz[baseline=-0.55ex]\node[draw=popink,line width=1.2pt,rounded corners=2.6pt,%
    fill=#2,inner xsep=6.5pt,inner ysep=3pt]{%
    \popxboldls\fontsize{7}{7}\selectfont\color{#3}\MakeUppercase{#1}};%
}
\newcommand{\popsquiggle}[1]{%
  \tikz[baseline=-0.4ex]{
    \draw[#1,line width=2.6pt,line cap=round]
      plot [smooth] coordinates {(0,0.09) (0.34,0.01) (0.68,0.21) (1.02,0.01) (1.36,0.10)};
  }%
}
% Mot-clé surligné (marqueur orange sous le texte)
\newcommand{\cle}[1]{{\popxbold\color{popdark}#1}}

% ============================================================
% En-tête / pied
% ============================================================
\pagestyle{fancy}
\fancyhf{}
\renewcommand{\headrulewidth}{0pt}
\renewcommand{\footrulewidth}{0pt}
\lhead{\raisebox{-0.15ex}{\poplogo\fontsize{13}{13}\selectfont\color{popink}OnBuch\color{poporange}+}}
\rhead{\poppill{\DOCMATIERE\ \textbullet\ \DOCNIVEAU\ \textbullet\ Leçon \DOCLECON}{white}{popink}}
\lfoot{\popxbold\fontsize{7.6}{7.6}\selectfont\color{popmuted}\MakeUppercase{Cours signé OnBuch+}\ \textbullet\ {\popsemi\DOCTITRE}}
\rfoot{\tikz[baseline=-0.6ex]\node[rounded corners=8.5pt,fill=popink,minimum height=17pt,minimum width=17pt,inner xsep=5pt,inner sep=0]{\popxbold\fontsize{8}{8}\selectfont\color{popcream}\thepage\,/\,\pageref*{LastPage}};}

% ============================================================
% Titres
% ============================================================
\newcommand{\chaptitre}[2]{%
  \begin{tcolorbox}[enhanced,colback=poporange,colframe=popink,boxrule=2.6pt,arc=11pt,
      drop shadow=popink,left=15pt,right=15pt,top=12pt,bottom=11pt,before skip=3pt,after skip=18pt]
    \poppill{#2}{popink}{popcream}\par\vspace{9pt}
    {\raggedright\hyphenpenalty=10000\exhyphenpenalty=10000\popdisplay\fontsize{22}{24}\selectfont\color{popcream}\MakeUppercase{#1}\par}
  \end{tcolorbox}
}
% Section numérotée : pastille orange avec le numéro + titre Archivo
\newcounter{popsec}
\newcommand{\coursec}[1]{%
  \stepcounter{popsec}\setcounter{popsub}{0}%
  \par\vspace{16pt}\noindent
  \begin{minipage}{\linewidth}
  \tikz[baseline=-0.7ex]\node[draw=popink,line width=1.6pt,fill=poporange,rounded corners=4pt,minimum size=22pt,inner sep=0]{\popdisplay\fontsize{12}{12}\selectfont\color{popcream}\thepopsec};%
  \hspace{8pt}{\popdisplay\fontsize{15}{17}\selectfont\color{popink}\MakeUppercase{#1}}\par
  \vspace{4pt}\noindent{\color{poporange}\rule{36pt}{3.6pt}}
  \end{minipage}\par\nopagebreak\vspace{8pt}
}
\newcounter{popsub}
\newcommand{\courssub}[1]{%
  \stepcounter{popsub}%
  \par\vspace{9pt}\noindent{\popxbold\fontsize{11.5}{13}\selectfont\color{popink}{\color{poporange}\thepopsec.\thepopsub}\hspace{6pt}#1}\par\nopagebreak\vspace{4pt}
}

% ============================================================
% Boîtes pédagogiques — même recette : fond blanc, bordure épaisse colorée,
% ombre dure encre, badge plein. Titre optionnel [#1].
% ============================================================
\tcbset{popbase/.style={breakable,enhanced,colback=white,boxrule=2.3pt,arc=9pt,
  shadow={3pt}{-3pt}{0pt}{popink},left=14pt,right=14pt,top=11pt,bottom=11pt,before skip=11pt,after skip=15pt,
  fontupper=\popsemi\fontsize{10.6}{14.5}\selectfont\color{popink},
  fontlower=\popsemi\fontsize{10.6}{14.5}\selectfont\color{popink}}}
% Coche dessinée (le glyphe ✓ n'existe pas dans les polices du document)
\newcommand{\popcheck}[1]{\tikz[baseline=-0.5ex]\draw[#1,line width=1.6pt,line cap=round,line join=round](-0.13,0.0)--(-0.04,-0.09)--(0.13,0.1);}
\providecommand{\coche}{\popcheck{popgreen}}
\providecommand{\resultat}[1]{\par\smallskip\noindent\fcolorbox{popgreen}{popgreenL}{\parbox{\dimexpr\linewidth-2\fboxsep-2\fboxrule\relax}{\textbf{Résultat.} #1}}\par\smallskip}
\newcommand{\popboxhead}[3]{\poppill{#1}{#2}{white}\ifx\relax#3\relax\else\hspace{6pt}{\popxbold\fontsize{10.6}{12}\selectfont\color{popink}#3}\fi\par\vspace{7pt}}

\newtcolorbox{aretenir}[1][]{popbase,colframe=popgreen,before upper={\popboxhead{À retenir}{popgreen}{#1}}}
\newtcolorbox{exemplebox}[1][]{popbase,colframe=poporange,before upper={\popboxhead{Exemple}{poporange}{#1}}}
\newtcolorbox{definition}[1][]{popbase,colframe=popblue,before upper={\popboxhead{Définition}{popblue}{#1}}}
\newtcolorbox{propriete}[1][]{popbase,colframe=poppurple,before upper={\popboxhead{Propriété}{poppurple}{#1}}}
\newtcolorbox{methode}[1][]{popbase,colframe=popgold,before upper={\popboxhead{Méthode}{popgold}{#1}}}
\newtcolorbox{attention}[1][]{popbase,colframe=poppink,before upper={\popboxhead{Attention}{poppink}{#1}}}
\newtcolorbox{experience}[1][]{popbase,colframe=popblue,colback=popblueL,before upper={\popboxhead{Expérience}{popblue}{#1}}}
% « Le savais-tu ? » : carte violette pleine (contexte, culture, Cameroun)
\newtcolorbox{savaistu}[1][]{popbase,colback=poppurple,colframe=popink,
  fontupper=\popsemi\fontsize{10.4}{14.2}\selectfont\color{white},
  before upper={\poppill{Le savais-tu\,?}{white}{poppurple}\ifx\relax#1\relax\else\hspace{6pt}{\popxbold\color{white}#1}\fi\par\vspace{7pt}}}
% Exercice résolu : énoncé en haut, \tcblower puis la solution sur fond crème
\newcounter{popexo}\newcounter{popexr}
\newtcolorbox{exoresolu}[1][]{popbase,colframe=poporange,colbacklower=popcream,
  segmentation style={popink,line width=1.2pt,dashed},
  before upper={\stepcounter{popexr}\popboxhead{Exercice résolu \thepopexr}{poporange}{#1}},
  before lower={\poppill{Solution}{popink}{popcream}\par\vspace{6pt}}}
% Exercice (non résolu en place) — la correction est dans la partie Corrigés
\newtcolorbox{exercice}[1][]{popbase,colframe=popink,
  before upper={\stepcounter{popexo}\popboxhead{Exercice \thepopexo}{popink}{#1}}}
% Corrigé
\newtcolorbox{corrige}[1][]{popbase,colframe=popgreen,colback=popgreenL,
  before upper={\popboxhead{Corrigé}{popgreen}{#1}}}
% Objectifs / prérequis / bilan
\newtcolorbox{objectifs}{popbase,colframe=popink,colback=poporangeL,
  before upper={\popboxhead{Objectifs de la leçon}{popink}{}}}
\newtcolorbox{prerequis}{popbase,colframe=popink,colback=popblueL,
  before upper={\popboxhead{Prérequis}{popblue}{}}}
\newtcolorbox{bilan}{popbase,colframe=popink,colback=white,boxrule=2.8pt,
  before upper={\popboxhead{Fiche bilan}{poporange}{L'essentiel à connaître pour l'examen}}}
% Figure encadrée avec légende
\newtcolorbox{popfigure}[1][]{enhanced,breakable=false,colback=white,colframe=popline,boxrule=1.4pt,arc=8pt,
  left=10pt,right=10pt,top=10pt,bottom=8pt,before skip=10pt,after skip=12pt,halign=center,
  lower separated=false,halign lower=center,
  fontlower=\popsemi\fontsize{9}{11.5}\selectfont\color{popmuted},
  #1}
\newcommand{\legende}[1]{\par\vspace{6pt}{\popsemi\fontsize{9}{11.5}\selectfont\color{popmuted}#1}}

% ============================================================
% Signature OnBuch+
% ============================================================
\newcommand{\poplogoinline}[1]{{\poplogo\fontsize{#1}{#1}\selectfont\color{popink}OnBuch\color{poporange}+}}
\newcommand{\signatureonbuch}{%
  \begin{tcolorbox}[enhanced,colback=popink,colframe=popink,boxrule=2.2pt,arc=10pt,
      drop shadow=poporange,left=14pt,right=14pt,top=10pt,bottom=10pt,before skip=0pt,after skip=0pt]
    \tikz[baseline=-0.7ex]\node[circle,fill=popgreen,draw=popcream,line width=1.3pt,minimum size=22pt,inner sep=0]{\popcheck{white}};%
    \hspace{9pt}\begin{minipage}[c]{0.62\linewidth}
      {\popxbold\fontsize{12}{13}\selectfont\color{popcream}Signé\ \textbullet\ {\poplogo OnBuch\color{poporange}+}}\par\vspace{2pt}
      {\raggedright\popsemi\fontsize{8}{10.5}\selectfont\color{popline}\MakeUppercase{Équipe pédagogique OnBuch+}\\\MakeUppercase{Conforme au programme officiel MINESEC}\par}
    \end{minipage}\hfill
    {\poplogo\fontsize{22}{22}\selectfont\color{popcream}OnBuch\color{poporange}+}
  \end{tcolorbox}%
}

% ============================================================
% Page de garde
% #1 titre · #2 matière · #3 niveau · #4 module · #5 n° leçon · #6 durée ·
% #7 description / objectifs (texte ou itemize)
% ============================================================
\newcommand{\pagegarde}[7]{%
  \thispagestyle{empty}
  \newgeometry{a4paper,margin=1.7cm,top=1.6cm,bottom=1.4cm}%
  \begin{tikzpicture}[remember picture,overlay]
    \fill[poporange!18] ([xshift=-3.2cm,yshift=-3.6cm]current page.north east) circle (4.6cm);
    \fill[poppurple!14] ([xshift=2.4cm,yshift=7.6cm]current page.south west) circle (3.8cm);
  \end{tikzpicture}%
  {\poplogo\fontsize{30}{30}\selectfont\color{popink}OnBuch\color{poporange}+}\hfill
  \raisebox{0.6ex}{\poppill{Cours signé \textbullet\ Premium}{popink}{popcream}}\par
  \vspace{1pt}\popsquiggle{poppurple}\par
  \vspace{8pt}
  \begin{tcolorbox}[enhanced,colback=poporange,colframe=popink,boxrule=2.8pt,arc=13pt,
      drop shadow=popink,left=18pt,right=18pt,top=12pt,bottom=13pt,after skip=10pt]
    \poppill{#2 \textbullet\ #3}{popink}{popcream}\hspace{5pt}\poppill{#4}{white}{popink}\par\vspace{8pt}
    {\popdisplay\fontsize{14}{14}\selectfont\color{popink}LEÇON #5}\par\vspace{4pt}
    {\raggedright\hyphenpenalty=10000\exhyphenpenalty=10000\popdisplay\fontsize{25}{27}\selectfont\color{popcream}\MakeUppercase{#1}\par}
    \vspace{8pt}
    {\popxbold\fontsize{9.5}{9.5}\selectfont\color{popink}\MakeUppercase{Durée conseillée : #6}}
  \end{tcolorbox}
  \begin{tcolorbox}[enhanced,colback=white,colframe=popink,boxrule=2.2pt,arc=10pt,
      drop shadow=popink,left=16pt,right=16pt,top=10pt,bottom=10pt,after skip=8pt]
    \poppill{Dans cette leçon}{poppurple}{white}\par\vspace{5pt}
    \popsemi\fontsize{9.3}{12.6}\selectfont\color{popink}#7
  \end{tcolorbox}
  \noindent\begin{minipage}[t]{0.487\linewidth}
    \begin{tcolorbox}[enhanced,colback=popgreen,colframe=popink,boxrule=2.2pt,arc=10pt,
        drop shadow=popink,left=11pt,right=11pt,top=10pt,bottom=10pt,height=3.1cm,valign=center]
      \tcbox[colback=white,colframe=popink,boxrule=1.3pt,arc=5pt,boxsep=0pt,nobeforeafter,
        left=3pt,right=3pt,top=3pt,bottom=3pt,valign=center]{\qrcode[height=1.9cm]{https://play.google.com/store/apps/details?id=com.luvvix.onbuch&hl=fr}}%
      \hspace{8pt}\begin{minipage}[c]{0.5\linewidth}
        {\popdisplay\fontsize{11}{12.5}\selectfont\color{white}\MakeUppercase{Télécharge OnBuch}}\par\vspace{4pt}
        {\raggedright\popsemi\fontsize{8.4}{11}\selectfont\color{white}Gratuit \textbullet\ Google Play\\Tuteur IA, annales, concours, résultats\par}
      \end{minipage}
    \end{tcolorbox}
  \end{minipage}\hfill
  \begin{minipage}[t]{0.487\linewidth}
    \begin{tcolorbox}[enhanced,colback=poppurple,colframe=popink,boxrule=2.2pt,arc=10pt,
        drop shadow=popink,left=11pt,right=11pt,top=10pt,bottom=10pt,height=3.1cm,valign=center]
      \tikz[baseline=-0.7ex]\node[circle,fill=white,draw=popink,line width=1.3pt,minimum size=1.6cm,inner sep=0]{\popdisplay\fontsize{24}{24}\selectfont\color{poppurple}+};%
      \hspace{8pt}\begin{minipage}[c]{0.6\linewidth}
        {\popdisplay\fontsize{11}{12.5}\selectfont\color{white}\MakeUppercase{Passe à OnBuch+}}\par\vspace{4pt}
        {\raggedright\popsemi\fontsize{8.4}{11}\selectfont\color{white}Tous les cours signés, Tuteur IA illimité, fiches PDF — onglet OnBuch+ de l'app\par}
      \end{minipage}
    \end{tcolorbox}
  \end{minipage}
  \vspace{8pt}
  \popcontenu
  \vfill
  \signatureonbuch
  \restoregeometry
  \newpage
}

% Rangée « Ce cours contient » (page de garde)
\newcommand{\poptile}[3]{% #1 couleur #2 gros texte #3 légende
  \begin{tcolorbox}[enhanced,colback=#1,colframe=popink,boxrule=2pt,arc=9pt,shadow={2.5pt}{-2.5pt}{0pt}{popink},
      left=6pt,right=6pt,top=9pt,bottom=9pt,halign=center,valign=center,height=1.75cm,width=\linewidth,nobeforeafter]
    {\popdisplay\fontsize{12.5}{13}\selectfont\color{white}\MakeUppercase{#2}}\par\vspace{3pt}
    {\centering\popsemi\fontsize{7.8}{9.5}\selectfont\color{white}#3\par}
  \end{tcolorbox}}
\newcommand{\popcontenu}{%
  \par\noindent{\popxboldls\fontsize{7.5}{7.5}\selectfont\color{popmuted}CE COURS SIGNÉ CONTIENT}\par\vspace{7pt}
  \noindent\begin{minipage}[t]{0.235\linewidth}\poptile{popblue}{Cours}{complet, illustré, pas à pas}\end{minipage}\hfill
  \begin{minipage}[t]{0.235\linewidth}\poptile{poporange}{Méthodes}{et exercices résolus}\end{minipage}\hfill
  \begin{minipage}[t]{0.235\linewidth}\poptile{poppink}{Exercices}{type examen + corrigés}\end{minipage}\hfill
  \begin{minipage}[t]{0.235\linewidth}\poptile{popgreen}{Fiche bilan}{l'essentiel à retenir}\end{minipage}\par}

% Sommaire
\newtcolorbox{sommaire}{enhanced,colback=white,colframe=popink,boxrule=2.2pt,arc=10pt,
  drop shadow=popink,left=18pt,right=18pt,top=13pt,bottom=13pt,before skip=6pt,after skip=20pt,
  before upper={\poppill{Sommaire}{poppurple}{white}\par\vspace{9pt}\popsemi\fontsize{10.6}{14.5}\selectfont\color{popink}}}

% Signature de fin de cours — poussée tout en bas de la dernière page
\newcommand{\findecours}{%
  \par\vfill\nopagebreak
  \begin{center}\popsquiggle{poporange}\end{center}\vspace{4pt}
  \signatureonbuch
}
\AtBeginDocument{\def\llbracket{\mathopen{[\mkern-3.5mu[}}\def\rrbracket{\mathclose{]\mkern-3.5mu]}}} % intervalles d'entiers (glyphe ⟦ absent de la police)
% Glyphes absents de Fira Math : redéfinis à partir de glyphes disponibles
\AtBeginDocument{%
  \def\setminus{\mathbin{\backslash}}\let\smallsetminus\setminus
  \def\vdots{\mathord{\vbox{\baselineskip3.2pt\lineskiplimit0pt\kern4pt\hbox{.}\hbox{.}\hbox{.}}}}%
  \def\ddots{\mathinner{\mkern1mu\raise6.5pt\hbox{.}\mkern2mu\raise3.6pt\hbox{.}\mkern2mu\raise0.7pt\hbox{.}\mkern1mu}}%
  \def\triangle{\mathord{\scalebox{0.9}{$\symup{\Delta}$}}}%
  \def\nabla{\mathord{\raisebox{\depth}{\scalebox{1}[-1]{$\symup{\Delta}$}}}}%
  \def\checkmark{\popcheck{popdark}}%
  \def\star{\mathbin{\ast}}\def\bigstar{\mathord{\ast}}%
  \def\diamond{\mathbin{\scalebox{0.6}{$\Diamond$}}}%
  \def\bigcup{\mathop{\mathchoice{\raisebox{-0.3em}{\scalebox{1.7}{$\cup$}}}{\scalebox{1.25}{$\cup$}}{\cup}{\cup}}\displaylimits}%
  \def\bigcap{\mathop{\mathchoice{\raisebox{-0.3em}{\scalebox{1.7}{$\cap$}}}{\scalebox{1.25}{$\cap$}}{\cap}{\cap}}\displaylimits}%
  \def\lesseqqgtr{\mathrel{\lessgtr}}\def\bigoplus{\mathop{\oplus}\displaylimits}%
}
\providecommand{\ssi}{si et seulement si} % abréviation fréquente
\AtBeginDocument{\providecommand{\wideparen}{\overparen}} % arc de cercle

%% FIGFIX-BEGIN v4 — correctifs globaux des figures (docs/onbuch-generation/tools/figfix)
\usepackage{adjustbox}\usepackage{etoolbox}
% toute figure TikZ/pgfplots plus large que la ligne est réduite à la largeur disponible
\BeforeBeginEnvironment{tikzpicture}{\begin{adjustbox}{max width=\linewidth}}
\AfterEndEnvironment{tikzpicture}{\end{adjustbox}}
% légendes pgfplots lisibles par-dessus les courbes
\pgfplotsset{every axis/.append style={legend style={fill=white,fill opacity=0.92,text opacity=1,draw=black!15,inner sep=3pt}}}
% étiquettes d'arêtes (syntaxe edge["..."]) avec fond blanc
\tikzset{every edge quotes/.append style={fill=white,inner sep=1.2pt}}
%% FIGFIX-END
\begin{document}

\pagegarde{Contraction : Rédiger un résumé}{Français}{Tle Tech}{Français – classe de Terminale technique}{N04}{15 séances (fiche de progression harmonisée)}{Cette leçon outille l'élève de Terminale technique pour maîtriser la contraction de texte : repérer les idées essentielles d'un texte explicatif, les reformuler sans les déformer, puis rédiger un résumé cohérent et cohésif. Elle s'appuie sur l'étude du Vieux nègre et la médaille de Ferdinand Oyono et sur des documents authentiques de la presse camerounaise.
\vspace{4pt}
\begin{multicols}{2}\small
\begin{itemize}
\item À la fin de cette leçon, je sais identifier les caractéristiques d'un texte explicatif (visée informative, connecteurs logiques, objectivité)
\item À la fin de cette leçon, je sais distinguer les idées essentielles des idées secondaires (exemples, répétitions, digressions) dans un texte
\item À la fin de cette leçon, je sais reformuler une idée par synonymie, changement de catégorie grammaticale et transformation de phrase
\item À la fin de cette leçon, je sais respecter la règle du quart et le plan du texte de départ dans un résumé
\item À la fin de cette leçon, je sais garantir la cohésion et la cohérence de mon résumé grâce aux connecteurs et aux reprises nominales
\item À la fin de cette leçon, je sais analyser les tonalités satirique et polémique dans Le Vieux nègre et la médaille
\end{itemize}\end{multicols}}

\begin{sommaire}
\begin{multicols}{2}
\begin{enumerate}
\item Le texte explicatif : caractéristiques et tonalités
\item Cohésion, cohérence et progression du texte
\item Le sigle et l'emprunt : outils lexicaux de la contraction
\item Reformuler les idées essentielles
\item Rédiger, produire et améliorer le résumé
\item Le Vieux nègre et la médaille : contexte, exposé et bilan de l'œuvre
\item Activité d'intégration
\item Méthodes et exercices résolus
\item Exercices
\item Corrigés des exercices
\item Fiche bilan
\end{enumerate}
\end{multicols}
\end{sommaire}

\begin{objectifs}
\begin{itemize}
\item À la fin de cette leçon, je sais identifier les caractéristiques d'un texte explicatif (visée informative, connecteurs logiques, objectivité)
\item À la fin de cette leçon, je sais distinguer les idées essentielles des idées secondaires (exemples, répétitions, digressions) dans un texte
\item À la fin de cette leçon, je sais reformuler une idée par synonymie, changement de catégorie grammaticale et transformation de phrase
\item À la fin de cette leçon, je sais respecter la règle du quart et le plan du texte de départ dans un résumé
\item À la fin de cette leçon, je sais garantir la cohésion et la cohérence de mon résumé grâce aux connecteurs et aux reprises nominales
\item À la fin de cette leçon, je sais analyser les tonalités satirique et polémique dans Le Vieux nègre et la médaille
\item À la fin de cette leçon, je sais situer une œuvre dans son contexte historique et littéraire (colonisation, négritude)
\item À la fin de cette leçon, je sais rédiger et justifier une démarche de résumé applicable à des situations concrètes (note de synthèse, compte rendu)
\end{itemize}
\end{objectifs}

\begin{prerequis}
\begin{itemize}
\item La notion de paragraphe et de plan de texte (classe de Première)
\item Les connecteurs logiques et les articulateurs du texte (Seconde et Première)
\item Les classes grammaticales et les transformations de phrases (morphosyntaxe de Première)
\item La lecture méthodique : thème, thèse, procédés d'écriture
\item Lecture intégrale du Vieux nègre et la médaille (séances de lecture cursive)
\end{itemize}
\end{prerequis}

\newpage

\coursec{Le texte explicatif : caractéristiques et tonalités}

\noindent\textbf{Situation d'accroche.} À l'occasion du concours d'entrée à l'ENAM, un candidat de Yaoundé doit résumer en 150 mots un article de 800 mots du quotidien \emph{Cameroon Tribune} sur la décentralisation. Pressé par le temps, il recopie des phrases entières de l'auteur : il est sanctionné pour plagiat et manque de reformulation. Comment faire tenir l'essentiel d'un texte dans un quart de sa longueur sans trahir sa pensée ni copier ses mots ? C'est tout l'art du résumé, compétence exigée au Baccalauréat technique et dans la vie professionnelle (comptes rendus de réunion, notes de synthèse administratives, rapports techniques).

\medskip\noindent\textbf{Problématique.} Comment identifier les caractéristiques du texte explicatif pour en extraire l'essentiel et le reformuler, tout en distinguant les registres (satirique, polémique) qui peuvent s'y glisser, comme dans \emph{Le Vieux nègre et la médaille} ?

\courssub{Définition et visée du texte explicatif}

\begin{definition}[Texte explicatif]
Un \cle{texte explicatif} est un texte dont la visée est d'informer et d'expliquer objectivement un fait, une notion, un phénomène ou une situation, sans chercher à persuader ni à raconter une histoire. Il s'adresse à l'intelligence du lecteur pour lui faire comprendre \emph{comment} ou \emph{pourquoi} quelque chose fonctionne ou existe.
\end{definition}

\noindent Le texte explicatif répond à une intention \cle{informative} et \cle{objective}. L'auteur s'efface derrière le sujet : il ne donne pas son avis, il ne juge pas, il expose. On le rencontre dans les manuels scolaires, les articles scientifiques vulgarisés, les notices techniques, les rapports administratifs, les encyclopédies.

\begin{aretenir}[Caractéristiques principales]
\begin{itemize}
    \item \textbf{Visée informative :} transmettre un savoir, clarifier un concept, décrire un mécanisme.
    \item \textbf{Vocabulaire précis et dénotatif :} termes techniques, définitions, absence d'ambiguïté, peu d'adjectifs qualificatifs subjectifs.
    \item \textbf{Temps dominant :} \cle{présent de vérité générale} (ou gnomique) : « L'eau \textbf{bout} à 100°C », « La décentralisation \textbf{permet} une gestion de proximité ». Ce présent exprime une loi, une constante, une vérité intemporelle.
    \item \textbf{Connecteurs logiques (opérateurs argumentatifs) :} ils organisent la démonstration : cause (\emph{car, parce que, puisque, du fait de}), conséquence (\emph{donc, par conséquent, c'est pourquoi, de sorte que}), opposition/concession (\emph{mais, cependant, toutefois, bien que}), énumération (\emph{d'abord, ensuite, enfin}), illustration (\emph{par exemple, notamment}).
    \item \textbf{Impersonnalité :} usage de « on », de la voix passive, de sujets inanimés (« Ce phénomène s'explique par... »), absence du « je » d'opinion.
\end{itemize}
\end{aretenir}

\begin{exemplebox}[Extrait type : notice technique]
\textit{« Le fonctionnement du moteur à combustion interne repose sur un cycle en quatre temps. \textbf{Premièrement}, l'admission : le piston descend et aspire le mélange air-essence. \textbf{Deuxièmement}, la compression : le piston remonte et comprime le mélange. \textbf{Troisièmement}, l'explosion : la bougie produit une étincelle qui enflamme le mélange ; la détente des gaz repousse le piston. \textbf{Enfin}, l'échappement : le piston remonte et chasse les gaz brûlés. \textbf{C'est pourquoi} ce cycle se répète en continu tant que le moteur est alimenté. »}
\par\medskip\noindent\textbf{Analyse :} Visée informative (expliquer le fonctionnement). Vocabulaire technique (admission, compression, bougie, détente). Présent de vérité générale (repose, descend, aspire, remonte, produit, enflamme, repousse, chasse, se répète). Connecteurs énumératifs (premierement, deuxiemement, troisiemement, enfin) et causaux (c'est pourquoi). Impersonnalité totale.
\end{exemplebox}

\courssub{Structure type du texte explicatif}

\noindent Un texte explicatif bien construit suit une architecture en trois temps, garantissant la \cle{cohérence} (liens logiques entre les idées) et la \cle{cohésion} (liens linguistiques entre les énoncés).

\begin{popfigure}
\begin{tikzpicture}[
    box/.style={draw=popblue, thick, fill=popblueL, rounded corners, minimum height=1.2cm, text width=3.5cm, align=center, font=\small},
    arrow/.style={-Stealth, thick, popdark},
    title/.style={font=\bfseries\small, text=popdark}
]
\node[box, title, align=center] (intro){Introduction du sujet\\ \textit{Accroche, définition, délimitation,\\ énoncé du plan}};
\node[box, title, below=1.5cm of intro, align=center] (dev){Développement explicatif\\ \textit{Phases successives, causes/effets,\\ exemples, illustrations,\\ connecteurs logiques}};
\node[box, title, below=1.5cm of dev, align=center] (conc){Conclusion / Synthèse\\ \textit{Rappel de l'essentiel, ouverture,\\ conséquence pratique}};

\draw[arrow] (intro.south) -- (dev.north);
\draw[arrow] (dev.south) -- (conc.north);

% Annotations on sides
\node[left=0.5cm of intro, text width=2.5cm, align=right, font=\small\itshape, text=popmuted] {« Qu'est-ce que c'est ? »};
\node[left=0.5cm of dev, text width=2.5cm, align=right, font=\small\itshape, text=popmuted] {« Comment / Pourquoi ? »};
\node[left=0.5cm of conc, text width=2.5cm, align=right, font=\small\itshape, text=popmuted] {« En résumé / Et après ? »};
\end{tikzpicture}
\legende{Schéma de la structure type d'un texte explicatif : trois temps logiques pour une information claire.}
\end{popfigure}

\begin{methode}[Analyser la structure d'un texte explicatif]
\begin{enumerate}
    \item \textbf{Repérer l'introduction :} elle pose le sujet (définition, champ d'application), annonce le plan (souvent implicitement par des connecteurs : \emph{dans un premier temps, ensuite...}).
    \item \textbf{Découper le développement :} identifier les grandes parties (souvent des paragraphes distincts). Chaque partie traite un aspect : causes, conséquences, description d'étapes, classification, comparaison. Noter les connecteurs qui relient les parties.
    \item \textbf{Identifier la conclusion :} elle synthétise les points clés, peut ouvrir sur une application, une perspective future ou une question liée.
    \item \textbf{Vérifier la progression :} les idées s'enchaînent-elles logiquement (chronologique, causal, du général au particulier, du simple au complexe) ?
\end{enumerate}
\end{methode}

\begin{exoresolu}[Repérage de structure]
\textbf{Énoncé.} Voici un court texte. Identifiez les trois parties de la structure et soulignez les connecteurs qui les relient.

\textit{« La photosynthèse est le processus par lequel les plantes vertes transforment l'énergie lumineuse en énergie chimique. \textbf{Elle se déroule en deux phases principales.} \textbf{Dans un premier temps}, la phase photochimique : la lumière est captée par la chlorophylle et permet de scinder l'eau en oxygène (rejeté) et en hydrogène (stocké). \textbf{Dans un second temps}, la phase chimique : l'hydrogène stocké se combine avec le gaz carbonique de l'air pour former du glucose. \textbf{Au final}, la plante a produit sa matière organique et rejeté de l'oxygène indispensable à la respiration des êtres vivants. »}

\textbf{Solution détaillée.}
\begin{itemize}
    \item \textbf{Introduction :} « La photosynthèse est le processus par lequel les plantes vertes transforment l'énergie lumineuse en énergie chimique. » $\rightarrow$ Définition du phénomène.
    \item \textbf{Développement :} « Elle se déroule en deux phases principales. \textbf{Dans un premier temps}... \textbf{Dans un second temps}... » $\rightarrow$ Deux paragraphes (ou deux blocs) explicatifs, organisés chronologiquement (séquence temporelle). Connecteurs : \emph{dans un premier temps, dans un second temps}.
    \item \textbf{Conclusion :} « \textbf{Au final}, la plante a produit sa matière organique et rejeté de l'oxygène indispensable à la respiration des êtres vivants. » $\rightarrow$ Bilan (synthèse) + ouverture sur l'utilité écologique (respiration des êtres vivants). Connecteur : \emph{au final}.
\end{itemize}
\end{exoresolu}

\courssub{Distinction : texte explicatif, argumentatif, narratif}

\noindent La confusion entre ces trois types de textes est une source majeure d'erreurs en contraction : on résume un texte argumentatif comme s'il était explicatif (on perd la thèse), ou on traite un texte satirique comme purement informatif (on rate l'ironie). Le tableau ci-dessous fixe les repères décisifs.

\begin{popfigure}
\begin{tikzpicture}[
    cell/.style={draw=popline, minimum height=1.1cm, align=center, font=\small},
    header/.style={cell, fill=popblue, font=\bfseries\small, text=white},
    rowA/.style={cell, fill=popblueL},
    rowB/.style={cell, fill=popcream},
    rowC/.style={cell, fill=popblueL},
    rowD/.style={cell, fill=popcream},
    rowE/.style={cell, fill=popblueL},
]
% Column widths
\def\cA{2.8cm}
\def\cB{4.2cm}
\def\cC{4.2cm}
\def\cD{4.2cm}

% Header
\node[header, minimum width=\cA] (h1) at (0,0) {Critère};
\node[header, minimum width=\cB, anchor=west] (h2) at (h1.east) {Texte \textbf{explicatif}};
\node[header, minimum width=\cC, anchor=west] (h3) at (h2.east) {Texte \textbf{argumentatif}};
\node[header, minimum width=\cD, anchor=west] (h4) at (h3.east) {Texte \textbf{narratif}};

% Row 1: Visée
\node[rowA, minimum width=\cA, anchor=north] (r1c1) at (h1.south west) {Visée \cle{communicative}};
\node[rowA, minimum width=\cB, anchor=north west] (r1c2) at (r1c1.east) {Informer, expliquer, faire comprendre (objectivité)};
\node[rowA, minimum width=\cC, anchor=north west] (r1c3) at (r1c2.east) {Convaincre, persuader, faire adhérer à une thèse (subjectivité assumée)};
\node[rowA, minimum width=\cD, anchor=north west] (r1c4) at (r1c3.east) {Raconter, divertir, émouvoir, témoigner};

% Row 2: Temps dominant
\node[rowB, minimum width=\cA, anchor=north] (r2c1) at (r1c1.south west) {Temps dominant};
\node[rowB, minimum width=\cB, anchor=north west] (r2c2) at (r2c1.east) {Présent de vérité générale (gnomique)};
\node[rowB, minimum width=\cC, anchor=north west] (r2c3) at (r2c2.east) {Présent (énonciation), futur (engagement), passé (exemples historiques)};
\node[rowB, minimum width=\cD, anchor=north west] (r2c4) at (r2c3.east) {Passé simple / Imparfait (récit), Présent de narration};

% Row 3: Marques linguistiques
\node[rowC, minimum width=\cA, anchor=north] (r3c1) at (r2c1.south west) {Marques linguistiques clés};
\node[rowC, minimum width=\cB, anchor=north west] (r3c2) at (r3c1.east) {Vocabulaire dénotatif, précis, technique ; connecteurs logiques (cause, conséquence, opposition) ; impersonnalité (on, passif, sujets inanimés)};
\node[rowC, minimum width=\cC, anchor=north west] (r3c3) at (r3c2.east) {Modaux (devoir, pouvoir, falloir), évaluatifs (adjectifs, métaphores), connecteurs argumentatifs (donc, car, cependant), « je » ou « nous » d'opinion, questions rhétoriques};
\node[rowC, minimum width=\cD, anchor=north west] (r3c4) at (r3c3.east) {Verbes d'action, de parole, de pensée ; connecteurs temporels (puis, alors, soudain) ; descriptions, dialogues ; focalisation (interne/externe)};

% Row 4: Structure
\node[rowD, minimum width=\cA, anchor=north] (r4c1) at (r3c1.south west) {Structure type};
\node[rowD, minimum width=\cB, anchor=north west] (r4c2) at (r4c1.east) {Introduction (sujet, plan) $\rightarrow$ Développement explicatif (phases, causes/effets) $\rightarrow$ Conclusion (synthèse, ouverture)};
\node[rowD, minimum width=\cC, anchor=north west] (r4c3) at (r4c2.east) {Introduction (thèse, accroche) $\rightarrow$ Arguments + exemples (thèse/antithèse/synthèse ou accumulation) $\rightarrow$ Conclusion (réaffirmation, ouverture)};
\node[rowD, minimum width=\cD, anchor=north west] (r4c4) at (r4c3.east) {Situation initiale $\rightarrow$ Élément perturbateur $\rightarrow$ Péripéties $\rightarrow$ Élément de résolution $\rightarrow$ Situation finale};

% Row 5: Exemple de domaine
\node[rowE, minimum width=\cA, anchor=north] (r5c1) at (r4c1.south west) {Domaines fréquents};
\node[rowE, minimum width=\cB, anchor=north west] (r5c2) at (r5c1.east) {Manuels, notices, articles scientifiques, rapports techniques, encyclopédies};
\node[rowE, minimum width=\cC, anchor=north west] (r5c3) at (r5c2.east) {Éditoriaux, plaidoiries, discours politiques, essais, lettres d'opinion, publicités};
\node[rowE, minimum width=\cD, anchor=north west] (r5c4) at (r5c3.east) {Romans, contes, nouvelles, récits de vie, journaux intimes, légendes};
\end{tikzpicture}
\legende{Tableau comparatif des trois grands types de textes : repères pour ne pas les confondre lors de la lecture et de la contraction.}
\end{popfigure}

\begin{attention}[Piège fréquent : confondre explication et argumentation]
Un texte peut \emph{paraître} explicatif (vocabulaire précis, connecteurs logiques) tout en défendant une thèse. \textbf{Test décisif :} cherchez la \cle{thèse} (une opinion défendable, sujette à débat). Si l'auteur prend position (« Il est urgent de... », « La meilleure solution est... », « Nous pensons que... »), c'est un texte \textbf{argumentatif}, même s'il explique des faits pour les besoins de sa démonstration. Dans \emph{Le Vieux nègre et la médaille}, le discours du commandant \emph{explique} la « grandeur de la France » (registre explicatif de surface) mais \emph{argumente} pour justifier la colonisation (registre argumentatif/idéologique) et \emph{raille} Meka (registre satirique). Le résumé doit rendre compte de cette complexité.
\end{attention}

\courssub{Les tonalités satirique et polémique : quand l'explicatif bascule}

\noindent Un texte explicatif peut être le support d'une \cle{tonalité} particulière qui en modifie le sens profond. Deux tonalités sont au programme et s'illustrent parfaitement dans \emph{Le Vieux nègre et la médaille} de Ferdinand Oyono.

\begin{definition}[Tonalité satirique]
La \cle{tonalité satirique} est une attitude d'écriture qui raille et dénonce les défauts, les vices ou les travers d'une personne, d'un groupe ou d'une société par l'\cle{ironie}, la \cle{caricature} et le \cle{ridicule}. Elle vise à corriger par le rire (ou le sourire amer) en montrant l'absurde d'une situation. L'auteur dit le contraire de ce qu'il pense, ou exagère à dessein.
\end{definition}

\begin{definition}[Tonalité polémique]
La \cle{tonalité polémique} est une attitude d'écriture agressive, engagée, qui attaque violemment un adversaire, une idée, une institution. Elle use de l'invective, de la charge directe, de l'argumentation offensive. Elle ne rit pas : elle condamne. Elle est proche de l'argumentatif virulent.
\end{definition}

\begin{aretenir}[Différence clé]
\begin{itemize}
    \item \textbf{Satirique :} détournement, ironie, second degré, rire (jaune), implicite. « Quel \emph{merveilleux} discours ! » (pour dire qu'il est nul).
    \item \textbf{Polémique :} frontal, premier degré, colère, explicite. « Ce discours est une insulte à l'intelligence. »
    \item \textbf{Lien :} La satire \emph{peut} être une arme de la polémique (ironie destructrice), mais la polémique n'est pas nécessairement satirique.
\end{itemize}
\end{aretenir}

\begin{methode}[Repérer l'ironie (et la satire) dans un texte]
\begin{enumerate}
    \item \textbf{Chercher l'écart (le décalage) :} entre ce qui est \emph{dit} (littéral) et ce qui est \emph{pensé} (réel). C'est le cœur de l'\cle{antiphrase} (dire le contraire de ce qu'on pense).
    \item \textbf{Repérer les marqueurs d'exagération (hyperbole) :} éloges excessifs, adjectifs surévalués (\emph{magnifique, grandiose, illustre, divin}), superlatifs ironiques.
    \item \textbf{Identifier les contrastes :} entre le discours officiel/solennel et la réalité vécue/misérable ; entre le ton pompeux et le contenu vide ou cruel.
    \item \textbf{Analyser le contexte énonciatif :} qui parle ? à qui ? dans quelle circonstance ? (ex. : un colonisateur décorant un vieux colonisé).
    \item \textbf{Tester l'inversion :} si je retourne la phrase (je remplace l'éloge par une critique), le sens devient-il plus cohérent avec la situation réelle ? Si oui, c'est de l'ironie.
\end{enumerate}
\end{methode}

\begin{savaistu}[Ferdinand Oyono et \emph{Le Vieux nègre et la médaille} (1956)]
Ferdinand Oyono (1929--2010), haut fonctionnaire et diplomate camerounais, est une figure majeure de la littérature africaine d'expression française. Ce roman court, écrit à 27 ans, dénonce l'assimilation coloniale française par le portrait tragi-comique de Meka, un vieux chef « évolué » qui attend fièrement la médaille de la Légion d'honneur. La cérémonie de remise devient le théâtre d'une satire féroce : le discours du commandant (texte explicatif de surface sur la « mission civilisatrice ») est miné par l'ironie narrative d'Oyono qui en montre le grotesque et la violence symbolique. L'œuvre annonce \emph{Une vie de boy} et \emph{Le Chemin d'Europe}.
\end{savaistu}

\courssub{Application : la cérémonie de remise de la médaille dans \emph{Le Vieux nègre et la médaille}}

\noindent Le passage du discours du commandant (chapitre 6) est un modèle de \textbf{texte à tonalité satirique} qui emprunte les \emph{apparences} du texte explicatif (vocabulaire solennel, structure logique, présent de vérité générale) pour mieux le détourner.

\begin{exemplebox}[Analyse guidée : le discours du commandant (extrait reconstitué)]
\textit{« Mes amis, nous sommes réunis aujourd'hui pour honorer un homme qui a su, par son dévouement sans faille et sa fidélité exemplaire, mériter la reconnaissance de la Mère Patrie. \textbf{C'est donc avec une légitime fierté} que je remets cette médaille à Meka, \textbf{ce digne fils de la France}, \textbf{ce modèle de vertus civiques}. \textbf{Car}, n'oublions pas que la civilisation ne se décrète pas : elle se mérite par le travail et l'obéissance aux lois supérieures de la République. \textbf{Aussi}, que cet exemple inspire la jeunesse ! »}

\par\medskip\noindent\textbf{Analyse pas à pas :}
\begin{enumerate}
    \item \textbf{Apparence explicative :} Structure Introduction (sujet : honneur, Meka) $\rightarrow$ Développement (causes : dévouement, fidélité ; loi : civilisation = travail/obéissance) $\rightarrow$ Conclusion (exemple pour la jeunesse). Connecteurs logiques : \emph{c'est donc avec..., car, aussi}. Présent de vérité générale (\emph{se mérite, inspire}). Vocabulaire abstrait (\emph{reconnaissance, civilisation, vertus civiques, lois supérieures}).
    \item \textbf{Rupture satirique (ironie narrative) :} Le narrateur (Oyono) nous a montré avant cela la réalité : Meka a travaillé comme un forçat, a perdu ses terres, a vu ses fils partir à la guerre, vit dans la misère, est traité en enfant par l'administration. Le « dévouement » est une exploitation ; la « fidélité » est une soumission forcée ; la « Mère Patrie » est une marâtre.
    \item \textbf{Marqueurs d'ironie (antiphrases) :}
        \begin{itemize}
            \item « \textbf{ce digne fils de la France} » $\rightarrow$ Antiphrase : Meka est un « sujet » exploité, pas un fils chéri.
            \item « \textbf{ce modèle de vertus civiques} » $\rightarrow$ Antiphrase : sa « vertu » est sa résignation.
            \item « \textbf{la civilisation ne se décrète pas : elle se mérite par le travail et l'obéissance} » $\rightarrow$ Cynisme : la « civilisation » coloniale s'impose par la force et le travail forcé.
        \end{itemize}
    \item \textbf{Contraste (caricature) :} Le ton solennel, pompeux, paternaliste du commandant contraste avec la petitesse de la médaille (un bout de métal), la misère de Meka, l'indifférence de l'assistance. La caricature grossit le trait : le commandant incarne la bêtise coloniale satisfaite.
    \item \textbf{Effet produit :} Le lecteur rit (jaune) de l'aveuglement du commandant et ressent la tragédie de Meka. La satire dénonce le système colonial en montrant le ridicule de ses rites.
\end{enumerate}
\end{exemplebox}

\begin{exoresolu}[Entraînement : repérer l'ironie chez Oyono]
\textbf{Énoncé.} Relisez mentalement (ou sur le texte) le moment où Meka reçoit la médaille et où le commandant lui pince la joue en disant : « \textit{Bravo, mon vieux ! Tu as bien travaillé pour la France.} » Expliquez en trois lignes pourquoi cette réplique relève de la tonalité satirique (ironie + caricature).

\textbf{Solution détaillée.}
\begin{enumerate}
    \item \textbf{Ironie (antiphrase) :} « \textit{Tu as bien travaillé pour la France} » dit sur un ton affectueux (tutoiement, « mon vieux ») masque une réalité d'exploitation (travail forcé, impôts, recrutement forcé des fils). Le sens littéral (éloge) est contredit par le contexte (soumission).
    \item \textbf{Caricature :} Le geste de pincer la joue infantilise Meka (vieux chef respecté chez lui) et le réduit à un bon chien fidèle. La caricature grossit le trait paternaliste/raciste du colonisateur.
    \item \textbf{Dénonciation :} Cette familiarité humiliante, déguisée en récompense, résume la violence symbolique de l'assimilation : on dépossède l'Africain de sa dignité d'adulte pour l'intégrer comme « enfant » dans la « famille française ».
\end{enumerate}
\end{exoresolu}

\begin{exercice}[Entraînement à la distinction de registres]
\textbf{Consigne.} Pour chacun des trois courts extraits ci-dessous, indiquez : (1) le type de texte dominant (explicatif / argumentatif / narratif) ; (2) la tonalité éventuelle (neutre / satirique / polémique) ; (3) deux indices linguistiques justifiant votre réponse.

\textbf{Texte A.} « Le paludisme est une maladie parasitaire transmise par la piqûre du moustique anophèle femelle. Le parasite (plasmodium) se développe dans le foie puis dans les globules rouges qu'il détruit. La fièvre survient par crises. La prévention repose sur les moustiquaires imprégnées et l'assainissement. »

\textbf{Texte B.} « \textit{Quel magnifique spectacle que cette assemblée de notables, si empressés à lécher les bottes du représentant de l'Empire ! Leur servilité n'a d'égale que leur bêtise. Ils croient devenir français en reniant leurs pères. Pathétique comédie.} »

\textbf{Texte C.} « Il est inadmissible que le budget alloué à la santé rurale stagne depuis cinq ans pendant que les dépenses de prestige explosent. Le gouvernement doit, immédiatement, réorienter 15\% du budget souverain vers les centres de santé de district. C'est une question de survie pour nos populations. »
\end{exercice}

\begin{corrige}[Exercice n°1]
\textbf{Texte A.} (1) \textbf{Explicatif}. (2) \textbf{Neutre}. (3) Vocabulaire technique dénotatif (\emph{plasmodium, anophèle, globules rouges, moustiquaires imprégnées}) ; présent de vérité générale (\emph{est, se développe, détruit, survient, repose}) ; connecteurs logiques (\emph{puis, ensuite}) ; impersonnalité totale.

\textbf{Texte B.} (1) \textbf{Argumentatif} (sous forme d'invective) / \textbf{Narratif} (description d'une scène). (2) \textbf{Satirique} (et polémique). (3) Ironie / Antiphrase : « \emph{magnifique spectacle} » pour une scène humiliante ; « \emph{lécher les bottes} » (métaphore caricaturale de la soumission) ; hyperboles (\emph{n'a d'égale que leur bêtise}) ; jugement de valeur explicite (\emph{pathétique comédie}) ; « je » implicite du narrateur moqueur.

\textbf{Texte C.} (1) \textbf{Argumentatif}. (2) \textbf{Polémique}. (3) Modaux de nécessité/obligation (\emph{doit, immédiatement}) ; vocabulaire évaluatif négatif (\emph{inadmissible, stagnent, explosent}) ; thèse claire (\emph{réorienter 15\%...}) ; connecteurs argumentatifs (\emph{pendant que, c'est une question de}) ; « nos populations » (identification/engagement).
\end{corrige}

\begin{aretenir}[Bilan de la section 1]
\begin{itemize}
    \item Le \textbf{texte explicatif} informe objectivement (présent gnomique, vocabulaire précis, connecteurs logiques, structure Intro/Développement/Conclusion).
    \item Il se distingue de l'\textbf{argumentatif} (thèse, subjectivité, modaux) et du \textbf{narratif} (récit, temps du passé, verbes d'action).
    \item La \textbf{tonalité satirique} (ironie, antiphrase, caricature, second degré) et la \textbf{tonalité polémique} (attaque frontale, premier degré) peuvent investir un texte à structure explicative.
    \item Dans \emph{Le Vieux nègre et la médaille}, le discours officiel (explicatif de façade) est retourné par l'ironie narrative d'Oyono en satire coloniale.
    \item \textbf{Pour le résumé :} il faut identifier le type de texte \textbf{et} sa tonalité pour ne pas résumer l'ironie au premier degré (erreur fatale : résumer le discours du commandant comme un vrai éloge de Meka).
\end{itemize}
\end{aretenir}

\coursec{Cohésion, cohérence et progression du texte}

Dans la section précédente, nous avons étudié le texte explicatif : sa structure (thèse, arguments, exemples), ses marques linguistiques (connecteurs logiques, présent de vérité générale, lexique précis) et ses tonalités, notamment satirique et polémique. Nous savons donc désormais \emph{reconnaître} un texte explicatif et en dégager les idées essentielles. Mais pour \emph{contracter} ce texte, c'est-à-dire pour en rédiger un résumé, il ne suffit pas de sélectionner les idées : il faut les \textbf{relier} entre elles. Un résumé n'est pas une liste de phrases juxtaposées ; c'est un \cle{texte}, c'est-à-dire un tissu de phrases qui s'enchaînent naturellement. C'est précisément ce que garantissent trois notions fondamentales que nous allons étudier maintenant : la \cle{cohérence}, la \cle{cohésion} et la \cle{progression thématique}.

\courssub{La cohérence : l'unité de sens du texte}

Commençons par une intuition simple. Lisez ces trois phrases :

\emph{Le Cameroun a adopté un Plan National de Développement. Les éléphants se nourrissent de feuilles et d'écorces. Il fait beau ce matin à Douala.}

Chaque phrase, prise isolément, est correcte en français. Pourtant, l'ensemble ne forme pas un texte : on sent que quelque chose ne fonctionne pas. Il manque un \textbf{lien de sens} entre les phrases. À l'inverse, lisez :

\emph{Le Cameroun a adopté un Plan National de Développement. Ce plan vise à moderniser les infrastructures du pays. Il prévoit notamment la construction de routes et de barrages hydroélectriques.}

Ici, tout s'enchaîne : chaque phrase prolonge la précédente, et l'ensemble parle d'un seul et même sujet. Cette propriété s'appelle la cohérence.

\begin{definition}[La cohérence]
La \cle{cohérence} est l'organisation logique et sémantique des idées qui donne au texte son \textbf{unité de sens}. Un texte cohérent est un texte dont toutes les phrases se rapportent à un \textbf{thème unique} et s'organisent selon un ordre logique perceptible (chronologique, spatial, cause-conséquence, du général au particulier, etc.).
\end{definition}

La cohérence repose sur trois exigences :

\begin{itemize}
\item \textbf{L'unité thématique} : le texte traite d'un seul sujet. Toute phrase étrangère au sujet rompt la cohérence.
\item \textbf{L'ordre logique} : les idées s'enchaînent selon un plan (problème puis solution, cause puis conséquence, général puis particulier).
\item \textbf{L'absence de contradiction} : le texte ne doit pas affirmer une chose puis son contraire sans explication.
\end{itemize}

\begin{attention}[Cohérence et résumé]
Dans un résumé, la cohérence impose de respecter \textbf{l'ordre des idées du texte de départ}. Si l'auteur expose d'abord le problème, puis les causes, puis les solutions, le résumé doit suivre le même mouvement. Bouleverser cet ordre, c'est trahir la pensée de l'auteur — et c'est une faute sanctionnée à l'examen.
\end{attention}

\courssub{La cohésion : les liens linguistiques du texte}

Si la cohérence est une affaire de \emph{sens}, la cohésion est une affaire de \emph{langue}. Reprenons notre exemple : « Ce plan vise... Il prévoit... ». Les mots \emph{ce plan} et \emph{il} renvoient au \emph{Plan National de Développement} mentionné dans la première phrase : ce sont des fils linguistiques qui cousent les phrases entre elles.

\begin{definition}[La cohésion]
La \cle{cohésion} est l'ensemble des procédés linguistiques (\cle{connecteurs}, \cle{reprises} nominales et pronominales, \cle{substitutions}, \cle{ellipses}) qui relient les phrases entre elles et assurent la \textbf{continuité} du texte en surface.
\end{definition}

\begin{aretenir}[Cohérence et cohésion : le couple inséparable]
La cohérence est le lien de \textbf{sens} (profondeur) ; la cohésion est le lien de \textbf{forme} (surface). Un texte peut être cohérent sans être parfaitement cohésif (sens clair mais phrases sèches), mais un texte plein de connecteurs sans unité de sens reste incohérent. Le bon résumé exige les deux.
\end{aretenir}

\textbf{Les principaux outils de la cohésion}

\begin{enumerate}
\item \textbf{L'anaphore pronominale} : un pronom reprend un élément déjà exprimé. \emph{Le gouvernement a lancé le chantier. \textbf{Il} espère l'achever en trois ans.} (\emph{il} = le gouvernement)
\item \textbf{L'anaphore nominale} : un nom, souvent accompagné d'un déterminant démonstratif, reprend un élément antérieur. \emph{Le Plan National de Développement a été adopté. \textbf{Ce plan} ambitionne de...} ; \emph{...la construction de routes. \textbf{Ces infrastructures} désenclaveront les villages.}
\item \textbf{La substitution lexicale} : on reprend une idée par un synonyme, un hyperonyme ou une périphrase. \emph{Le vieux Meka reçut une médaille. \textbf{Cette récompense} / \textbf{cet honneur} le remplit d'orgueil.}
\item \textbf{Les connecteurs} : mots ou expressions qui marquent les relations logiques : \emph{d'abord, ensuite, en effet, car, cependant, donc, ainsi, enfin, c'est pourquoi}.
\item \textbf{Le champ lexical} : la répétition de mots de la même famille de sens tisse le texte (pour le développement : \emph{croissance, investissement, infrastructures, modernisation}).
\item \textbf{L'ellipse} : on supprime un élément déjà connu pour éviter la répétition. \emph{Meka reçut une médaille, ses amis [reçurent] des félicitations.}
\end{enumerate}

\begin{popfigure}
\begin{center}
\begin{tikzpicture}[font=\small]
% Paragraphe annoté : chaîne anaphorique
\node[anchor=west] (p1) at (0,0) {\textbf{Le Plan National de Développement}$_1$ a été adopté par le gouvernement.};
\node[anchor=west] (p2) at (0,-0.9) {\textbf{Ce plan}$_2$ vise à moderniser les infrastructures du pays.};
\node[anchor=west] (p3) at (0,-1.8) {\textbf{Il}$_3$ prévoit la construction de routes et de barrages.};
\node[anchor=west] (p4) at (0,-2.7) {\textbf{Ces ouvrages}$_4$ désenclaveront les régions rurales.};
% Flèches de reprise
\draw[->,popblue,thick] (p2.west) ++(0.15,-0.15) to[bend right=35] node[left,popblue,font=\scriptsize]{anaphore nominale} ($(p1.west)+(0.15,0.05)$);
\draw[->,poporange,thick] (p3.west) ++(0.15,-0.15) to[bend right=35] node[left,poporange,font=\scriptsize]{pronom} ($(p2.west)+(0.15,0.05)$);
\draw[->,popgreen,thick] (p4.west) ++(0.15,-0.15) to[bend right=35] node[left,popgreen,font=\scriptsize]{substitution} ($(p3.west)+(0.15,0.05)$);
% Cadre de la chaîne
\node[draw=poppurple,rounded corners,fit={(p1)(p4)},inner sep=6pt] (cadre) {};
\node[below=2pt of cadre,font=\scriptsize\itshape,popmuted] {Chaîne anaphorique : chaque maillon renvoie au précédent.};
\end{tikzpicture}
\end{center}
\legende{Diagramme en chaîne des reprises anaphoriques dans un paragraphe : 1. thème initial ; 2. reprise nominale (\emph{ce plan}) ; 3. reprise pronominale (\emph{il}) ; 4. substitution lexicale (\emph{ces ouvrages} reprenant \emph{routes et barrages}).}
\end{popfigure}

\begin{methode}[Vérifier la cohésion d'un résumé]
\begin{enumerate}
\item \textbf{Relire} le résumé à voix haute : toute phrase qui semble « tomber du ciel » signale une rupture.
\item \textbf{Souligner} tous les connecteurs (\emph{d'abord, ensuite, cependant, donc...}) : il doit y en avoir entre les grandes étapes du raisonnement.
\item \textbf{Entourer} les reprises nominales et pronominales (\emph{ce plan, il, cette mesure...}) : vérifier que chacune renvoie clairement à un élément précédent.
\item \textbf{Vérifier} qu'aucune phrase n'est isolée : chaque phrase doit être accrochée à la précédente par au moins un procédé de liaison.
\item \textbf{Corriger} les ruptures en ajoutant un connecteur ou une reprise nominale.
\end{enumerate}
\end{methode}

\begin{attention}[Le piège de la liste]
L'erreur la plus fréquente dans un résumé est d'\textbf{empiler les idées sans connecteurs} : \emph{Le texte parle du plan. Il parle des routes. Il parle des barrages. Il parle de l'emploi.} Ce n'est pas un texte, c'est une liste. Le correcteur du Baccalauréat technique sanctionne ce défaut : reliez toujours vos idées (\emph{d'abord... ensuite... enfin... ; non seulement... mais aussi...}).
\end{attention}

\courssub{La notion de progression thématique}

Pour comprendre la progression, il faut d'abord distinguer, dans chaque phrase, deux éléments :

\begin{definition}[Thème et rhème]
Dans une phrase, le \cle{thème} est ce dont on parle (l'information déjà connue, généralement placée au début) ; le \cle{rhème} est ce qu'on en dit (l'information nouvelle, généralement placée à la fin). Dans \emph{Le gouvernement \textbar{} a adopté le plan}, « le gouvernement » est le thème, « a adopté le plan » est le rhème.
\end{definition}

La \cle{progression thématique} décrit la manière dont les thèmes et les rhèmes s'enchaînent de phrase en phrase. On distingue trois types principaux.

\textbf{1. La progression à thème constant} : toutes les phrases parlent du même thème, qui est repris (par le même mot, un pronom ou un synonyme). \emph{Le plan est ambitieux. Il vise la modernisation du pays. Il concerne tous les secteurs.}

\textbf{2. La progression linéaire (ou en chaîne)} : le rhème d'une phrase devient le thème de la phrase suivante. \emph{Le gouvernement a adopté un plan. Ce plan prévoit des barrages. Ces barrages produiront de l'électricité.}

\textbf{3. La progression à thème dérivé} : les thèmes des phrases dérivent d'un hyperthème général, comme des branches issues d'un même tronc. \emph{L'économie camerounaise se transforme. L'agriculture se modernise. L'industrie se diversifie. Les services se développent.} (Tous les thèmes dérivent de « l'économie camerounaise ».)

\begin{propriete}[Lois de la progression thématique]
Dans une \textbf{progression linéaire}, le rhème d'une phrase devient le thème de la phrase suivante : le texte avance comme une chaîne, chaque maillon s'accrochant au précédent. Dans une \textbf{progression à thème constant}, toutes les phrases portent sur le même thème : le texte approfondit un seul sujet. Dans une \textbf{progression à thème dérivé}, chaque phrase éclaire un aspect d'un thème général : le texte procède par énumération organisée.
\end{propriete}

\begin{popfigure}
\begin{center}
\begin{tikzpicture}[font=\scriptsize, every node/.style={align=center}]
% Thème constant
\node[font=\small\bfseries,popblue] at (0,3.6) {Thème constant};
\node[draw,rounded corners,fill=popblueL] (t1) at (0,2.8) {T : Le plan};
\node[draw,rounded corners,fill=popblueL] (t2) at (0,1.9) {T : Il};
\node[draw,rounded corners,fill=popblueL] (t3) at (0,1.0) {T : Il};
\node[draw,rounded corners,fill=white] (r1) at (2.2,2.8) {R : est ambitieux};
\node[draw,rounded corners,fill=white] (r2) at (2.2,1.9) {R : modernise le pays};
\node[draw,rounded corners,fill=white] (r3) at (2.2,1.0) {R : crée des emplois};
\draw[->,popblue] (t1)--(t2); \draw[->,popblue] (t2)--(t3);
\draw[popmuted] (t1)--(r1); \draw[popmuted] (t2)--(r2); \draw[popmuted] (t3)--(r3);
% Linéaire
\node[font=\small\bfseries,poporange] at (6.4,3.6) {Progression linéaire};
\node[draw,rounded corners,fill=poporangeL] (a1) at (5.2,2.8) {T : Le gouvernement};
\node[draw,rounded corners,fill=white] (a2) at (8.2,2.8) {R : adopte un plan};
\node[draw,rounded corners,fill=poporangeL] (a3) at (5.2,1.9) {T : Ce plan};
\node[draw,rounded corners,fill=white] (a4) at (8.2,1.9) {R : prévoit des barrages};
\node[draw,rounded corners,fill=poporangeL] (a5) at (5.2,1.0) {T : Ces barrages};
\node[draw,rounded corners,fill=white] (a6) at (8.2,1.0) {R : produisent du courant};
\draw[->,poporange,thick] (a2.south) to[bend left=20] (a3.north);
\draw[->,poporange,thick] (a4.south) to[bend left=20] (a5.north);
\draw[popmuted] (a1)--(a2); \draw[popmuted] (a3)--(a4); \draw[popmuted] (a5)--(a6);
% Thème dérivé
\node[font=\small\bfseries,popgreen] at (1.5,-0.9) {Thème dérivé};
\node[draw,rounded corners,fill=popgreenL] (h) at (1.5,-1.7) {Hyperthème : l'économie};
\node[draw,rounded corners,fill=white] (d1) at (-1.4,-2.8) {L'agriculture...};
\node[draw,rounded corners,fill=white] (d2) at (1.5,-2.8) {L'industrie...};
\node[draw,rounded corners,fill=white] (d3) at (4.4,-2.8) {Les services...};
\draw[->,popgreen] (h)--(d1); \draw[->,popgreen] (h)--(d2); \draw[->,popgreen] (h)--(d3);
\end{tikzpicture}
\end{center}
\legende{Schéma des trois types de progression thématique. À gauche, le thème constant (T identique repris) ; à droite, la progression linéaire (le rhème R devient le thème suivant) ; en bas, la progression à thème dérivé (plusieurs thèmes issus d'un hyperthème).}
\end{popfigure}

\textbf{Effet sur la construction du sens.} Le choix d'une progression n'est pas neutre : la progression linéaire donne un texte \emph{dynamique}, qui avance étape par étape (idéal pour un raisonnement ou un récit) ; le thème constant donne un texte \emph{descriptif ou explicatif}, qui examine un sujet sous toutes ses faces ; le thème dérivé donne un texte \emph{organisé en tableau}, commode pour présenter les aspects d'une situation. Dans un résumé, il faut \textbf{reproduire la progression dominante du texte de départ} pour en respecter la démarche.

\begin{exoresolu}[Identifier les types de progression]
\textbf{Énoncé.} Voici un extrait adapté d'un article de presse camerounaise sur le Plan National de Développement :

\emph{(1) Le Plan National de Développement a été présenté hier à Yaoundé. (2) Ce plan ambitionne de faire du Cameroun un pays émergent. (3) Il mise d'abord sur les infrastructures. (4) Les infrastructures routières seront réhabilitées sur tout le territoire. (5) Leur réhabilitation facilitera le transport des produits agricoles vers les villes.}

Indiquez la progression qui relie les phrases (1) à (3), puis celle qui relie les phrases (3) à (5). Justifiez.

\tcblower
\textbf{Solution détaillée.}
\begin{itemize}
\item \textbf{Phrases (1) à (3) : progression à thème constant.} Le thème « Plan National de Développement » de la phrase (1) est repris par l'anaphore nominale « Ce plan » en (2), puis par le pronom « Il » en (3). Les trois phrases parlent du même sujet : le plan.
\item \textbf{Phrases (3) à (5) : progression linéaire.} Le rhème de (3), « les infrastructures », devient le thème de (4) (« Les infrastructures routières ») ; le rhème de (4), « réhabilitées / réhabilitation », devient le thème de (5) (« Leur réhabilitation »). Chaque information nouvelle devient le point de départ de la phrase suivante.
\end{itemize}
\textbf{Conclusion.} L'article commence par approfondir un thème unique (le plan), puis enchaîne les idées en chaîne : c'est le mouvement classique du texte explicatif journalistique, que le résumé devra reproduire.
\end{exoresolu}

\begin{exercice}[À vous de jouer]
Lisez ce court texte : \emph{La formation technique est une priorité nationale. Elle prépare les jeunes aux métiers de l'industrie. Ces métiers sont aujourd'hui très recherchés par les entreprises. Les entreprises camerounaises manquent en effet de techniciens qualifiés.}

1. Relevez deux reprises anaphoriques (une pronominale, une nominale).
2. Identifiez le type de progression dominant du passage et justifiez.
3. Rédigez en deux phrases liées le résumé de ce passage, en utilisant au moins un connecteur.
\end{exercice}

\begin{corrige}[Exercice : progression et reprises]
1. Reprise pronominale : « Elle » (phrase 2) reprend « La formation technique ». Reprise nominale : « Ces métiers » (phrase 3) reprend « les métiers de l'industrie » ; « Les entreprises camerounaises » (phrase 4) reprend « les entreprises ».
2. Progression linéaire dominante : le rhème de chaque phrase devient le thème de la suivante (formation technique $\rightarrow$ métiers $\rightarrow$ entreprises).
3. Exemple de résumé : « La formation technique, priorité nationale, prépare les jeunes aux métiers de l'industrie. Or ces techniciens qualifiés sont aujourd'hui très recherchés par les entreprises camerounaises. » (Connecteur : \emph{or} ; reprise : \emph{ces techniciens qualifiés}.)
\end{corrige}

\begin{aretenir}[L'essentiel de la section]
\begin{itemize}
\item La \cle{cohérence} = unité de sens (un seul thème, un ordre logique).
\item La \cle{cohésion} = liens linguistiques (anaphores, substitutions, connecteurs, champs lexicaux, ellipses).
\item La \cle{progression thématique} = enchaînement des thèmes et rhèmes : à thème constant, linéaire ou à thème dérivé.
\item Dans un résumé : respecter l'ordre des idées, relier chaque phrase à la précédente, reproduire la progression du texte source.
\end{itemize}
\end{aretenir}

Ces outils d'analyse du texte trouveront une application directe dans les sections suivantes : le sigle et l'emprunt, que nous étudierons ensuite, sont précisément des instruments lexicaux qui servent la contraction en condensant l'expression sans rompre ni la cohésion ni la cohérence du résumé.

\coursec{Le sigle et l'emprunt : outils lexicaux de la contraction}

Dans la section précédente, nous avons étudié la \cle{cohésion} et la \cle{cohérence} du texte : un texte bien construit est un tissu où chaque phrase est reliée aux autres par des reprises, des connecteurs et une progression logique. Mais rédiger un résumé, c'est aussi jouer sur le \cle{lexique} : il faut dire la même chose en moins de mots, sans perdre la précision. Deux outils lexicaux nous y aident puissamment, surtout dans les textes administratifs, économiques et journalistiques que l'on rencontre au Probatoire et au Baccalauréat technique : le \cle{sigle} et l'\cle{emprunt}. Cette section leur est consacrée.

\courssub{Le sigle : définition et observation}

\begin{definition}[Le sigle]
Un \cle{sigle} est une abréviation constituée des initiales d'un groupe de mots ou d'une expression. Il s'écrit en général en majuscules, sans points entre les lettres.
\begin{itemize}
\item MINESEC = \textbf{M}inistère des \textbf{E}nseignements \textbf{Sec}ondaires ;
\item CNUED = \textbf{C}onférence des \textbf{N}ations \textbf{U}nies sur le \textbf{C}ommerce et le \textbf{D}éveloppement ;
\item GICAM = \textbf{G}roupement \textbf{I}nterpatronal du \textbf{Cam}eroun.
\end{itemize}
\end{definition}

\textbf{Observation (démarche d'investigation).} Prenons un titre de la presse camerounaise : \emph{« La SONARA et le FEICOM signent une convention sous le regard du MINESEC »}. Trois questions se posent :
\begin{enumerate}
\item Ces mots en majuscules renvoient-ils à des réalités précises ? Oui : chacun désigne une institution identifiable.
\item Comment se prononcent-ils ? On dit \emph{la So-na-ra} comme un mot ordinaire, mais on épelera \emph{effe-i-com} ou \emph{em-i-ne-sec} lettre par lettre selon l'usage.
\item Peut-on les remplacer par leur forme développée ? Oui, mais la phrase devient alors très lourde.
\end{enumerate}
De cette observation naît la distinction fondamentale entre deux types de sigles.

\begin{definition}[Sigle et acronyme]
\begin{itemize}
\item Le \cle{sigle} proprement dit se prononce \textbf{lettre par lettre} : SNI (Société Nationale d'Investissement), CRTV, ONG.
\item L'\cle{acronyme} est un sigle que l'on prononce \textbf{comme un mot ordinaire} : ONU (\emph{o-nu}), SONARA (\emph{so-na-ra}), SMIC.
\end{itemize}
L'acronyme finit parfois par devenir un mot du lexique courant : beaucoup de locuteurs ignorent que \emph{laser} ou \emph{OVNI} sont nés comme des sigles.
\end{definition}

\begin{exemplebox}[Sigles camerounais courants]
\begin{center}
\begin{tabularx}{\linewidth}{|l|X|l|}
\hline
\rowcolor{popblueL} \textbf{Sigle} & \textbf{Forme développée} & \textbf{Prononciation} \\
\hline
MINESEC & Ministère des Enseignements Secondaires & lettre par lettre \\
\hline
FEICOM & Fonds Spécial d'Équipement et d'Intervention Intercommunale & lettre par lettre \\
\hline
SONARA & Société Nationale de Raffinage & acronyme \\
\hline
CRTV & Cameroon Radio Television & lettre par lettre \\
\hline
\end{tabularx}
\end{center}
\end{exemplebox}

\begin{popfigure}
\begin{tikzpicture}[font=\small]
\draw[fill=popblueL, draw=popblue, rounded corners] (0,0) rectangle (4.6,0.9);
\node[align=center] at (2.3,0.45) {\textbf{Sigle} : MINESEC};
\draw[fill=popgreenL, draw=popgreen, rounded corners] (5.2,0) rectangle (9.8,0.9);
\node[align=center] at (7.5,0.45) {\textbf{Forme développée}};
\draw[fill=poporangeL, draw=poporange, rounded corners] (10.4,0) rectangle (14.6,0.9);
\node[align=center] at (12.5,0.45) {\textbf{Nature}};
\foreach \y/\a/\b/\c in {
  -1.1/MINESEC/Ministère des Enseignements Secondaires/sigle épelé,
  -2.2/FEICOM/Fonds Spécial d'Équipement \ldots{} Intercommunale/sigle épelé,
  -3.3/SONARA/Société Nationale de Raffinage/acronyme,
  -4.4/CRTV/Cameroon Radio Television/sigle épelé
}{
  \draw[fill=popcream, draw=popline] (0,\y-0.45) rectangle (4.6,\y+0.45);
  \node at (2.3,\y) {\textbf{\a}};
  \draw[fill=popcream, draw=popline] (5.2,\y-0.45) rectangle (9.8,\y+0.45);
  \node[align=center] at (7.5,\y) {\b};
  \draw[fill=popcream, draw=popline] (10.4,\y-0.45) rectangle (14.6,\y+0.45);
  \node[align=center] at (12.5,\y) {\emph{\c}};
}
\end{tikzpicture}
\legende{Tableau à double entrée : quatre sigles camerounais courants, leur forme développée et leur nature (sigle épelé lettre par lettre ou acronyme prononcé comme un mot).}
\end{popfigure}

\courssub{Le comportement grammatical des sigles}

Le sigle obéit à des règles précises qu'il faut maîtriser pour ne pas perdre de points dans un résumé.

\begin{propriete}[Règles d'emploi des sigles]
\begin{enumerate}
\item \textbf{Invariabilité} : en général, le sigle ne prend \textbf{pas de marque du pluriel}. On écrit \emph{des ONG}, \emph{des PDG}, jamais \emph{des ONGs} ni \emph{des ONG'S}. Seuls quelques acronymes lexicalisés acceptent le pluriel (\emph{des lasers}).
\item \textbf{Article et genre} : le sigle prend l'article du \textbf{nom principal} de l'expression développée. On dit \emph{la SONARA} (la Société), \emph{le FEICOM} (le Fonds), \emph{l'ONU} (l'Organisation).
\item \textbf{Féminisation} : le sigle suit le genre du nom noyau ; on ne le féminise pas arbitrairement. On écrit \emph{la CRTV} parce que le nom noyau, \emph{Television}, est de genre féminin dans l'usage français.
\item \textbf{Typographie} : majuscules, sans points, sans italique pour les sigles d'institutions nationales.
\end{enumerate}
\end{propriete}

\begin{attention}[Erreur fréquente : le pluriel des sigles]
Beaucoup d'élèves ajoutent un \og s \fg{} aux sigles au pluriel sans vérifier l'usage : \emph{les ONGS}, \emph{les MINESECS}. C'est une faute. Le sigle est \textbf{invariable} : on écrit \emph{des ONG}, sans lettre finale ajoutée arbitrairement. De même, évitez l'apostrophe anglaise (\emph{des ONG's}), qui est un anglicisme fautif en français normatif.
\end{attention}

\courssub{L'emprunt : définition et sources}

\begin{definition}[L'emprunt]
Un \cle{emprunt} est un mot ou une expression importé d'une autre langue et intégré au lexique français, avec ou sans adaptation graphique et phonétique. Le français parlé et écrit au Cameroun emprunte à plusieurs sources :
\begin{itemize}
\item \textbf{à l'anglais} : \emph{management}, \emph{briefing}, \emph{marketing}, \emph{deal} ;
\item \textbf{aux langues camerounaises} : \emph{ndolè} (plat de feuilles), \emph{miondo} (bâton de manioc), \emph{mboa}, \emph{sawa} ;
\item \textbf{au pidgin et au camfranglais} : \emph{choco} (voler), \emph{mbarga}, \emph{ashia} (expression de compassion).
\end{itemize}
\end{definition}

\begin{popfigure}
\begin{tikzpicture}[
  font=\small,
  racine/.style={rectangle, rounded corners, draw=popdark, fill=popblueL, thick, align=center},
  branche/.style={rectangle, rounded corners, draw=popblue, fill=popcream, align=center},
  feuille/.style={rectangle, rounded corners, draw=popgreen, fill=popgreenL, align=center, font=\footnotesize},
  lien/.style={-{Stealth[length=2.5mm]}, thick, popmuted}
]
\node[racine, align=center] (R) at (0,0){\textbf{Les emprunts}\\ \textbf{en français}};
\node[branche, align=center] (A) at (-5.2,-2.2){Emprunts\\ à l'anglais};
\node[branche, align=center] (B) at (0,-2.2){Emprunts aux\\ langues nationales};
\node[branche, align=center] (C) at (5.2,-2.2){Emprunts au pidgin\\ et au camfranglais};
\node[feuille, align=center] (A1) at (-5.2,-4.2){management\\ briefing\\ marketing};
\node[feuille, align=center] (B1) at (0,-4.2){ndolè\\ miondo\\ sawa};
\node[feuille, align=center] (C1) at (5.2,-4.2){choco\\ ashia\\ mbarga};
\draw[lien] (R) -- (A);
\draw[lien] (R) -- (B);
\draw[lien] (R) -- (C);
\draw[lien] (A) -- (A1);
\draw[lien] (B) -- (B1);
\draw[lien] (C) -- (C1);
\end{tikzpicture}
\legende{Arbre de classification des emprunts dans le français en usage au Cameroun : trois grandes sources (anglais, langues nationales, pidgin/camfranglais) avec des exemples pour chaque branche.}
\end{popfigure}

\begin{savaistu}[Le camfranglais, créativité... à doser]
Le \cle{camfranglais}, parlé par de nombreux jeunes Camerounais, illustre une remarquable créativité lexicale : il mélange français, anglais, pidgin et langues nationales, crée des sigles et des emprunts nouveaux (\emph{go} pour \og fille \fg{}, \emph{way} pour \og affaire \fg{}). C'est un phénomène sociolinguistique riche. Mais attention : dans un \textbf{résumé en français normatif}, comme dans toute copie d'examen, le camfranglais est \textbf{à proscrire}. On n'écrira jamais \emph{le gars a choco l'argent} dans une contraction de texte, mais \emph{le jeune homme a volé l'argent}.
\end{savaistu}

\courssub{Rôle des sigles et des emprunts dans la contraction}

Pourquoi étudier ces deux outils dans une leçon sur le résumé ? Parce que la contraction de texte obéit à une double exigence : \textbf{l'économie} (dire en moins de mots) et la \textbf{précision} (ne rien déformer).

\begin{propriete}[Fonctions des sigles et emprunts dans le résumé]
\begin{enumerate}
\item \textbf{Économie de mots} : écrire \emph{la SONARA} au lieu de \emph{la Société Nationale de Raffinage} fait gagner quatre mots à chaque occurrence. Sur un résumé limité à un quart du texte, le gain est considérable.
\item \textbf{Précision référentielle} : le sigle désigne une institution unique, sans ambiguïté ; l'emprunt (\emph{ndolè}, \emph{management}) nomme une réalité qu'aucun mot français ne rend exactement. Les supprimer, c'est appauvrir le texte.
\item \textbf{Fidélité au texte source} : un résumé ne doit pas remplacer les termes de l'auteur par des périphrases personnelles ; conserver ses sigles et ses emprunts, c'est respecter sa pensée.
\end{enumerate}
\end{propriete}

\begin{methode}[Utiliser les sigles dans un résumé]
\begin{enumerate}
\item \textbf{Repérer} dans le texte source tous les sigles et toutes les expressions institutionnelles longues.
\item \textbf{Développer une fois} si nécessaire : à la première occurrence, si le sigle est peu connu ou absent du texte, donner la forme développée suivie du sigle entre parenthèses : \emph{le Fonds Spécial d'Équipement et d'Intervention Intercommunale (FEICOM)}. Si le sigle figure déjà dans le texte source, on peut le reprendre directement.
\item \textbf{Utiliser ensuite le sigle seul} pour toutes les reprises : cela fait gagner des mots à chaque occurrence.
\item \textbf{Vérifier la grammaire} : article correct (\emph{la} SONARA, \emph{le} FEICOM), pas de \og s \fg{} au pluriel, majuscules respectées.
\item \textbf{Compter les mots} : s'assurer que l'économie réalisée sert la limite imposée par la consigne.
\end{enumerate}
\end{methode}

\begin{exoresolu}[Contracter un passage riche en sigles et emprunts]
\textbf{Énoncé.} Résumez en environ 25 mots le passage suivant (58 mots), en exploitant les sigles et les emprunts :

\emph{« Le Fonds Spécial d'Équipement et d'Intervention Intercommunale a organisé hier à Yaoundé un briefing à l'intention des maires. Le management du Fonds a annoncé que la Société Nationale de Raffinage financerait l'équipement des marchés où l'on vend le ndolè et le miondo. »}

\tcblower
\textbf{Solution détaillée.}
\begin{enumerate}
\item \emph{Repérage} : deux expressions institutionnelles longues peuvent devenir des sigles (FEICOM, SONARA) ; deux emprunts à l'anglais (\emph{briefing}, \emph{management}) et deux emprunts aux langues nationales (\emph{ndolè}, \emph{miondo}) doivent être conservés car ils sont précis et courts.
\item \emph{Première occurrence} : on développe une fois si utile, puis on sigle. Ici, pour gagner un maximum de mots, on peut passer directement aux sigles, qui sont connus du lecteur camerounais.
\item \emph{Résumé proposé} : \og Le FEICOM a réuni les maires en briefing à Yaoundé : son management annonce que la SONARA financera l'équipement des marchés de ndolè et de miondo. \fg{} (27 mots).
\end{enumerate}
Le passage de 58 à 27 mots doit beaucoup aux sigles (8 mots économisés) et à la conservation des emprunts, qui évitent de longues périphrases (\emph{plat de feuilles de ndolè} aurait alourdi le résumé).
\end{exoresolu}

\begin{exercice}[Recensement dans la presse camerounaise]
Relevez dans un article du journal \emph{Cameroon Tribune} ou d'un texte administratif (circulaire, arrêté, communiqué) :
\begin{enumerate}
\item cinq sigles, en précisant pour chacun la forme développée et la nature (sigle épelé ou acronyme) ;
\item trois emprunts, en indiquant leur langue d'origine ;
\item rédigez ensuite un résumé de l'article en 40 mots au maximum, en appliquant la méthode d'emploi des sigles.
\end{enumerate}
\end{exercice}

\begin{corrige}[Exercice : recensement dans la presse]
Exemple de corrigé à partir d'un article sur l'éducation :
\begin{enumerate}
\item Sigles relevés : MINESEC (Ministère des Enseignements Secondaires, sigle épelé) ; MINESUP (Ministère de l'Enseignement Supérieur, sigle épelé) ; GCE (\emph{General Certificate of Education}, sigle épelé) ; CAPIEMP (Concours d'Accès aux Postes d'Inspecteurs de l'Enseignement Maternel et Primaire, sigle épelé) ; SYNES (Syndicat National des Enseignants du Supérieur, sigle épelé).
\item Emprunts : \emph{briefing} (anglais), \emph{curriculum} (latin, via l'anglais), \emph{ndolè} (langue du groupe duala).
\item Résumé possible (38 mots) : \og Le MINESEC a présidé un briefing sur la réforme du curriculum. Le CAPIEMP recrutera davantage d'enseignants, tandis que le SYNES demande au MINESUP de meilleures conditions de travail. \fg{}
\end{enumerate}
Critères de réussite : sigles correctement développés, aucun pluriel fautif, articles adaptés (\emph{le} MINESEC, \emph{la} SONARA), limite de mots respectée.
\end{corrige}

\begin{aretenir}[L'essentiel de la section]
\begin{itemize}
\item Le \cle{sigle} est une abréviation formée des initiales d'une expression (MINESEC, FEICOM) ; prononcé comme un mot, il devient un \cle{acronyme} (ONU, SONARA).
\item Le sigle est en général \textbf{invariable} (des ONG), prend l'article du nom principal (\emph{la} SONARA, \emph{le} FEICOM) et s'écrit en majuscules sans points.
\item L'\cle{emprunt} est un mot importé d'une autre langue : anglais (\emph{management}, \emph{briefing}), langues nationales (\emph{ndolè}, \emph{miondo}), pidgin.
\item Dans un résumé, sigles et emprunts assurent \textbf{économie} et \textbf{précision} : on développe une fois si nécessaire, puis on utilise le sigle ; on conserve les emprunts plutôt que de périphraser.
\item Le camfranglais, source vivante d'emprunts, est exclu du résumé en français normatif.
\end{itemize}
\end{aretenir}

Ces outils lexicaux maîtrisés, nous pouvons maintenant passer au cœur de la contraction : la \emph{reformulation des idées essentielles}, qui fera l'objet de la section suivante.

\coursec{Reformuler les idées essentielles}

Après avoir étudié les outils lexicaux que sont le sigle et l'emprunt — qui permettent de condenser l'information sans en perdre la substance —, nous abordons maintenant le cœur opérationnel de la contraction : la reformulation. C'est l'étape où le lecteur devient auteur : il doit restituer la pensée d'autrui avec ses propres mots, en ne gardant que l'indispensable. C'est un exercice de haute précision linguistique, sanctionné au Probatoire par le respect strict de la règle du quart et de la fidélité au texte source.

\courssub{Étape 1 : Lecture analytique, thème et thèse}

Avant toute écriture, la contraction exige une lecture \cle{analytique} (ou \cle{lecture active}). Il ne s'agit pas de « lire pour savoir de quoi ça parle », mais de disséquer le texte pour en extraire la structure argumentative.

\begin{definition}[Thème et thèse]
Le \textbf{thème} est le sujet général du texte (ex. : « l'exode rural », « la corruption », « l'agriculture camerounaise »). Il se formule en un nom ou un groupe nominal. La \textbf{thèse} est la position de l'auteur sur ce thème, l'idée-force qu'il défend et qu'il cherche à faire admettre. Elle se formule en une phrase affirmative complète (ex. : « L'exode rural vide les campagnes de leurs forces vives et menace la sécurité alimentaire »).
\end{definition}

\begin{methode}[Repérer le thème et la thèse]
\begin{enumerate}
\item \textbf{Lire une première fois} pour l'impression globale : de quoi parle-t-on ? (Thème).
\item \textbf{Lire une seconde fois} en surlignant les connecteurs logiques (\emph{donc, par conséquent, cependant, en effet, or, mais}) et les énoncés modaux (\emph{il faut que, il est nécessaire, je pense que}).
\item \textbf{Identifier la phrase-thèse} : souvent en introduction (annonce du plan) ou en conclusion (bilan). Elle résume l'intention argumentative.
\item \textbf{Vérifier} : chaque paragraphe doit apporter un argument ou une illustration au service de cette thèse.
\end{enumerate}
\end{methode}

\begin{exemplebox}[Application sur un extrait fictif]
\textit{« Le Cameroun possède d'immenses potentialités agricoles. Pourtant, le pays importe massivement du riz et du blé. Cette paradoxale situation s'explique par le désintérêt de la jeunesse pour l'agriculture, perçue comme une activité de subsistance peu rentable. Il est urgent de moderniser le secteur par la mécanisation et l'accès au crédit pour assurer notre souveraineté alimentaire. »}

\begin{itemize}
\item \textbf{Thème :} L'agriculture camerounaise / la souveraineté alimentaire.
\item \textbf{Thèse :} La modernisation de l'agriculture (mécanisation, crédit) est indispensable pour corriger le paradoxe d'un pays riche qui importe sa nourriture et garantir la souveraineté alimentaire.
\end{itemize}
\end{exemplebox}

\courssub{Distinction idées essentielles / idées secondaires}

C'est le tri décisif. Tout ce qui ne porte pas la thèse est candidat à l'élimination.

\begin{definition}[Idée essentielle vs Idée secondaire]
Une \textbf{idée essentielle} (ou \cle{idée directrice}) est un argument, un fait structurant ou une étape du raisonnement sans laquelle la thèse s'effondre. Sa suppression rend le texte incompréhensible ou trahit la pensée de l'auteur.
Une \textbf{idée secondaire} (ou \cle{idée accessoire}) sert à appuyer, illustrer, nuancer ou décorer l'idée essentielle : exemples chiffrés, citations d'autorité, anecdotes, comparaisons, répétitions, énumérations détaillées, figures de style.
\end{definition}

\begin{methode}[Test de l'élimination pour identifier l'essentiel]
\begin{enumerate}
\item Segmentez le texte en unités de sens (paragraphes ou groupes de paragraphes).
\item Pour chaque unité, posez la question : \textbf{« Si je supprime cette phrase, la thèse tient-elle encore ? »}
\item \textbf{OUI} $\rightarrow$ Idée secondaire (exemple, illustration, reformulation interne). \textbf{Action :} Supprimer ou généraliser.
\item \textbf{NON} $\rightarrow$ Idée essentielle (argument clé, cause, conséquence, thèse). \textbf{Action :} Conserver et reformuler.
\end{enumerate}
\end{methode}

\begin{attention}[Piège : Confondre « important » et « essentiel »]
Un exemple frappant, une statistique choc, une belle métaphore peuvent sembler « importants » rhétoriquement, mais ils restent \textbf{secondaires} structurellement. Au Probatoire, recopier un exemple détaillé (ex. : « M. Dupont, cultivateur à Bafia, a vu sa récolte doubler grâce au tracteur ») au lieu d'écrire « La mécanisation améliore les rendements » coûte des points (non-respect de la contraction, hors-sujet partiel).
\end{attention}

\begin{popfigure}
\begin{tikzpicture}[
    node distance=1.2cm and 1.5cm,
    box/.style={draw=popdark, thick, rounded corners, fill=popcream, text width=3.5cm, align=center, minimum height=1.8cm, font=\small},
    arrow/.style={-Stealth, thick, popblue},
    filter/.style={draw=poporange, thick, trapezium, trapezium left angle=70, trapezium right angle=110, fill=poporangeL, text width=3cm, align=center, minimum height=1.2cm, font=\small\bfseries}
]
\node[box, align=center] (source){\textbf{Texte Original}\\\emph{(ex. 400 mots)}\\\textit{Idées essentielles + Secondaires + Style}};
\node[filter, below=of source, align=center] (elim){\textbf{ÉLIMINATION}\\\textit{Exemples, citations,}\\\textit{répétitions, figures,\-énumérations}};
\node[filter, below=of elim, align=center] (gen){\textbf{GÉNÉRALISATION}\\\textit{Listes $\rightarrow$ Génériques}\\\textit{Détails $\rightarrow$ Catégories}};
\node[filter, below=of gen, align=center] (ref){\textbf{REFORMULATION}\\\textit{Synonymie, nominalisation,}\\\textit{condensation, voix passive}};
\node[box, below=of ref, fill=popgreenL, draw=popgreen, align=center] (resume){\textbf{RÉSUMÉ}\\\emph{(ex. 100 mots)}\\\textit{Idées essentielles seules}\\\textit{Langage neutre, propre}\\\textit{Ratio 1/4 respecté}};
\draw[arrow] (source) -- (elim) node[midway, right, font=\tiny, text=popmuted] {Tri};
\draw[arrow] (elim) -- (gen) node[midway, right, font=\tiny, text=popmuted] {Abstraction};
\draw[arrow] (gen) -- (ref) node[midway, right, font=\tiny, text=popmuted] {Réécriture};
\draw[arrow] (ref) -- (resume) node[midway, right, font=\tiny, text=popmuted] {Finalisation};
% Entonnoir visuel
\draw[popline, thick] (source.south west) -- (elim.north west) -- (gen.north west) -- (ref.north west) -- (resume.south west);
\draw[popline, thick] (source.south east) -- (elim.north east) -- (gen.north east) -- (ref.north east) -- (resume.south east);
\end{tikzpicture}
\legende{Schéma du processus de contraction en entonnoir : du texte source (400 mots) au résumé (100 mots) par élimination, généralisation et reformulation successives.}
\end{popfigure}

\courssub{Techniques d'élimination : alléger sans trahir}

L'élimination est la première passe de nettoyage. Elle vise le \cle{réducteur} (tout ce qui alourdit le rapport signal/bruit).

\begin{propriete}[Opérations d'élimination autorisées]
\begin{enumerate}
\item \textbf{Suppression des exemples et illustrations} : « \textit{Paul, Pierre et Jacques ont réussi...} » $\rightarrow$ \textit{Plusieurs acteurs ont réussi...} (ou suppression pure si l'idée est déjà dite).
\item \textbf{Suppression des citations et références} : « \textit{Comme l'écrivait Montaigne : "..."} » $\rightarrow$ \textit{Les philosophes soulignent que...} (ou suppression).
\item \textbf{Suppression des figures de style} : métaphores, comparaisons, hyperboles, ironie. On garde le sens littéral. « \textit{Il a mangé la vie à pleines dents} » $\rightarrow$ \textit{Il a profité intensément de la vie.}
\item \textbf{Suppression des redondances et pléonasmes} : « \textit{Monter en haut, sortir dehors, collaboration étroite} » $\rightarrow$ \textit{Monter, sortir, collaboration.}
\item \textbf{Condensation des énumérations détaillées} : Voir généralisation ci-dessous.
\end{enumerate}
\end{propriete}

\courssub{Techniques de généralisation : monter en abstraction}

La généralisation remplace du concret, du particulier, du multiple par du générique, de l'abstrait, de l'unitaire. C'est le passage de l'\cle{extension} (l'ensemble des objets désignés) à la \cle{compréhension} (les propriétés communes).

\begin{methode}[La généralisation par l'hyperonyme]
\begin{enumerate}
\item Repérer une liste (énumération) ou un exemple précis.
\item Identifier la catégorie commune (l'hyperonyme).
\item Remplacer la liste par l'hyperonyme (+ quantificateur si besoin : \textit{divers, plusieurs, de nombreux}).
\end{enumerate}
\end{methode}

\begin{exemplebox}[Exemples de généralisation]
\begin{center}
\begin{tabularx}{\linewidth}{|l|X|X|}\hline
\textbf{Source (Particulier / Liste)} & \textbf{Opération} & \textbf{Cible (Générique / Hyperonyme)} \\ \hline
Bananes, manioc, taro, ignames, patates douces & Regroupement catégorie alimentaire & \textbf{Produits vivriers} (ou \emph{tubercules et bananes plantains}) \\ \hline
L'enseignant, l'infirmier, le policier, le juge & Regroupement catégorie socioprofessionnelle & \textbf{Fonctionnaires} (ou \emph{agents de l'État}) \\ \hline
À Yaoundé, à Douala, à Bafoussam, à Garoua & Regroupement géographique & \textbf{Dans les grandes métropoles} (ou \emph{sur l'ensemble du territoire}) \\ \hline
Il a labouré, semé, désherbé, récolté & Regroupement séquentiel action & \textbf{Il a conduit le cycle cultural} \\ \hline
\end{tabularx}
\end{center}
\end{exemplebox}

\begin{attention}[Piège : Sur-généralisation / Perte de sens]
Remplacer « \textit{Le riz, le maïs et le mil} » par « \textit{Les céréales} » est correct. Les remplacer par « \textit{Les aliments} » est une \textbf{sur-généralisation} : on perd l'information cruciale (céréales vs tubercules). Le résumé doit rester \textbf{informativement équivalent} sur le plan thématique.
\end{attention}

\courssub{Techniques de reformulation : réécrire avec ses mots}

C'est l'étape noble : produire un texte \emph{neuf} syntaxiquement et lexicalement, mais \emph{identique} sémantiquement. Interdiction formelle du \cle{copier-coller} (plagiat sanctionné à l'examen).

\begin{definition}[Reformulation (Paraphrase)]
Opération consistant à exprimer le même sens propositionnel qu'un énoncé source, par des formes linguistiques différentes (lexique, syntaxe, morphologie), sans ajout, ni omission, ni distorsion de l'information pertinente.
\end{definition}

\begin{propriete}[Les quatre leviers syntaxico-lexicaux de la reformulation]
\begin{enumerate}
\item \textbf{Synonymie / Paraphrase lexicale} : Remplacer un mot par un équivalent contextuel.
    \begin{itemize}
    \item \textit{« Le coût de la vie \textbf{augmente} »} $\rightarrow$ \textit{« La vie \textbf{chérît} / \textbf{devient plus chère} »}
    \item \textit{« Il \textbf{a démontré} »} $\rightarrow$ \textit{« Il \textbf{a prouvé} / \textbf{a établi} »}
    \end{itemize}
\item \textbf{Transposition de catégorie (Nominalisation / Verbalisation)} : Transformer un verbe en nom (nominalisation) ou un nom en verbe. C'est le levier n°1 de la condensation.
    \begin{itemize}
    \item \textit{« Il \textbf{a résisté} courageusement »} (Groupe Verbal) $\rightarrow$ \textit{« \textbf{Sa résistance} a été courageuse »} (Groupe Nominal). Gain : suppression du sujet « il », de l'adverbe « courageusement » (absorbé par l'adjectif).
    \item \textit{« La \textbf{décision} du ministre »} $\rightarrow$ \textit{« Le ministre \textbf{a décidé} »} (Verbalisation, utile si le ministre est déjà le sujet de la phrase précédente).
    \end{itemize}
\item \textbf{Transformation Actif / Passif (Diathèse)} : Permet de changer le focus thématique et d'éviter la répétition du sujet agent.
    \begin{itemize}
    \item \textit{« L'État \textbf{construit} des routes »} (Actif) $\rightarrow$ \textit{« Des routes \textbf{sont construites} par l'État »} (Passif) ou mieux \textit{« Des routes \textbf{sont construites} »} (Passif agentless, si l'agent est connu).
    \end{itemize}
\item \textbf{Condensation / Subordination / Participiales} : Réduire des propositions subordonnées complétives ou relatives en infinitifs, participes, ou prépositions.
    \begin{itemize}
    \item \textit{« \textbf{Quand il a fini} ses études, \textbf{il est parti} »} $\rightarrow$ \textit{« \textbf{Après avoir fini} ses études, \textbf{il est parti} »} $\rightarrow$ \textit{« \textbf{Ses études finies}, \textbf{il est parti} »} (Absolu) $\rightarrow$ \textit{« \textbf{Après ses études}, \textbf{il est parti} »} (Prépositionnel).
    \item \textit{« L'homme \textbf{qui dirige} le projet »} $\rightarrow$ \textit{« L'homme \textbf{dirigeant} le projet »} (Participe présent) ou \textit{« Le \textbf{dirigeant} du projet »} (Nominalisation).
    \end{itemize}
\end{enumerate}
\end{propriete}

\begin{popfigure}
\begin{tikzpicture}[
    node distance=0.8cm and 1cm,
    src/.style={draw=popblue, thick, fill=popblueL, rounded corners, text width=5.5cm, align=left, font=\small, minimum height=1.5cm},
    tgt/.style={draw=popgreen, thick, fill=popgreenL, rounded corners, text width=5.5cm, align=left, font=\small, minimum height=1.5cm},
    arr/.style={-Stealth, thick, poppurple, bend left=15},
    lbl/.style={midway, above, font=\tiny\bfseries, text=poppurple, fill=white, inner sep=1pt}
]
\node[src] (s1) {\textbf{Source} : « \textit{Le gouvernement \textbf{a décidé} de \textbf{baisser} les prix \textbf{parce que} la population \textbf{manifestait}.} »};
\node[tgt, right=of s1] (t1) {\textbf{Reformulée} : « \textit{La \textbf{décision} du gouvernement de \textbf{baisser} les prix \textbf{fait suite aux} manifestations populaires.} »};
\draw[arr] (s1) -- (t1) node[lbl] {Nominalisation + Préposition};
\node[src, below=of s1] (s2) {\textbf{Source} : « \textit{Les agriculteurs \textbf{cultivent} le cacao, le café, le coton et l'hévéa.} »};
\node[tgt, right=of s2] (t2) {\textbf{Reformulée} : « \textit{Les agriculteurs \textbf{produisent} des \textbf{cultures d'exportation}.} »};
\draw[arr] (s2) -- (t2) node[lbl] {Synonymie + Généralisation};
\node[src, below=of s2] (s3) {\textbf{Source} : « \textit{Comme \textbf{le sol est fertile}, les rendements \textbf{sont élevés}.} »};
\node[tgt, right=of s3] (t3) {\textbf{Reformulée} : « \textit{\textbf{La fertilité du sol} \textbf{entraîne} des rendements élevés.} »};
\draw[arr] (s3) -- (t3) node[lbl] {Nominalisation cause + Verbe causal};
\node[src, below=of s3] (s4) {\textbf{Source} : « \textit{Il \textbf{a étudié} le dossier \textbf{avant de} \textbf{le signer}.} »};
\node[tgt, right=of s4] (t4) {\textbf{Reformulée} : « \textit{\textbf{Après l'étude} du dossier, \textbf{il l'a signé}.} »};
\draw[arr] (s4) -- (t4) node[lbl] {Nominalisation + Antéposition};
\end{tikzpicture}
\legende{Tableau de transformations types : chaque flèche illustre un levier majeur (nominalisation, généralisation, condensation) pour passer de la phrase source à la phrase reformulée, plus dense.}
\end{popfigure}

\courssub{La règle du quart : contrainte quantitative stricte}

Au Probatoire / Baccalauréat Technique, la contraction n'est pas une « synthèse libre ». C'est un exercice normé.

\begin{aretenir}[La Règle du Quart (Consigne d'examen)]
Le résumé ne doit pas dépasser \textbf{le quart (1/4) de la longueur du texte original}, avec une tolérance de \textbf{$\pm 10\%$} sur cette cible.
\begin{itemize}
\item \textbf{Comptage :} On compte les \textbf{mots} (séparés par des espaces). Les nombres, sigles, symboles (\%, \$), traits d'union comptent pour un mot.
\item \textbf{Calcul cible :} $N_{\text{cible}} = \frac{N_{\text{original}}}{4}$.
\item \textbf{Intervalle acceptable :} $\left[ 0,9 \times N_{\text{cible}} \ ; \ 1,1 \times N_{\text{cible}} \right]$.
\end{itemize}
\end{aretenir}

\begin{exemplebox}[Calcul chiffré : Texte de 800 mots]
\begin{align*}
N_{\text{original}} &= 800 \text{ mots} \\
N_{\text{cible}} &= 800 / 4 = 200 \text{ mots} \\
\text{Tolérance } \pm 10\% &= 20 \text{ mots} \\
\textbf{Intervalle acceptable :} &\quad \textbf{180 mots $\le$ Résumé $\le$ 220 mots}
\end{align*}
\textbf{Attention :} Un résumé de 179 mots = \textbf{Hors norme} (trop court, risque d'omission d'idées essentielles). Un résumé de 221 mots = \textbf{Hors norme} (trop long, échec de la contraction). Les deux sont pénalisés.
\end{exemplebox}

\begin{attention}[Erreur fréquente : L'opinion personnelle]
Le résumé est un \textbf{texte tiers} (discours rapporté). Toute marque de subjectivité de l'élève est interdite :
\begin{itemize}
\item \textit{« \textbf{À mon avis}, l'auteur a raison... »} $\rightarrow$ \textbf{Interdit.}
\item \textit{« \textbf{Heureusement}, le gouvernement a agit... »} $\rightarrow$ \textbf{Interdit.}
\item \textit{« \textbf{Malheureusement}, la situation empire. »} $\rightarrow$ \textbf{Interdit.}
\end{itemize}
Utilisez des introducteurs neutres : \textit{« L'auteur montre que... », « Le texte explique que... », « Selon l'auteur, ... »}.
\end{attention}

\courssub{Entraînement guidé : Article Cameroon Tribune (400 mots $\rightarrow$ 100 mots)}

Appliquons la chaîne complète sur un texte simulé type « Cameroon Tribune » (style journalistique, factuel, arguments chiffrés).

\begin{exoresolu}[Contraction d'un article de presse : « Agriculture : le défi de la mécanisation »]
\textbf{Texte original (environ 400 mots) :}

\begin{itshape}« Le Ministre de l'Agriculture et du Développement Rural, Gabriel Mbairobe, a présidé hier à Yaoundé la cérémonie de lancement officiel du Programme National de Mécanisation de l'Agriculture (PRONAMA). Ce programme ambitieux, financé à hauteur de 25 milliards de FCFA par l'État et ses partenaires techniques et financiers, vise à doter le pays de 500 tracteurs, 1 000 charrues et 2 000 semoirs d'ici à 2027. 

Lors de son allocution, le Ministre a souligné que l'agriculture camerounaise reste encore largement dominée par l'outil aratoire manuel, à savoir la houe et la machette. Cette situation, a-t-il dit, explique les faibles rendements observés dans les exploitations familiales qui représentent pourtant plus de 80\% de la production nationale. « Il est temps de tourner la page de l'agriculture de subsistance pour entrer dans l'ère de l'agriculture moderne, compétitive et rémunératrice », a-t-il déclaré avec force. 

Le PRONAMA prévoit également la création de 50 Centres de Services Agricoles Mécanisés (CESEAM) répartis dans les dix régions. Ces centres auront pour mission de louer le matériel aux groupements de producteurs à des coûts subventionnés, d'assurer la maintenance et de former les opérateurs. Le Ministre a insisté sur la rigueur dans la gestion de ces équipements pour éviter les dérives constatées par le passé, où du matériel public se retrouvait dans des fermes privées. 

Les organisations paysannes présentes, notamment la CNPC et la CAPRO, ont salué l'initiative tout en demandant une mise en œuvre rapide et transparente. Elles ont aussi réclamé l'amélioration des pistes rurales pour évacuer les productions. En conclusion, le Ministre a réaffirmé la volonté du Gouvernement de faire de l'agriculture le pilier de l'émergence à l'horizon 2035. »\end{itshape}

\textbf{Étape 1 : Analyse \& Repérage (Thème / Thèse / Idées essentielles)}
\begin{itemize}
\item \textbf{Thème :} Lancement du PRONAMA (mécanisation agriculture).
\item \textbf{Thèse :} Le PRONAMA (matériel, centres, formation) doit moderniser l'agriculture familiale (houe/machette $\rightarrow$ faible rendement) pour la rendre compétitive, sous condition de gestion rigoureuse et d'infrastructures routières.
\item \textbf{Idées essentielles (5 arguments) :}
    \begin{enumerate}
    \item Lancement PRONAMA (25 Mds FCFA) : 500 tracteurs, 1000 charrues, 2000 semoirs (2027).
    \item Constat : Agriculture dominée par outil manuel (houe/machette) $\rightarrow$ faibles rendements (exploitations familiales = 80\% prod).
    \item Objectif : Passer agriculture subsistance $\rightarrow$ agriculture moderne/competitive.
    \item Dispositif : 50 CESEAM (location subventionnée, maintenance, formation) dans 10 régions.
    \item Conditions réussite : Gestion rigoureuse (anti-détournement) + Pistes rurales (évacuation).
    \end{enumerate}
\item \textbf{Idées secondaires (à éliminer/généraliser) :} Noms propres (Gabriel Mbairobe, Yaoundé, CNPC, CAPRO, 2035), citations directes (« Il est temps... »), chiffres détaillés du matériel (remplacés par « matériel agricole »), anecdote détournement passé, allocution/salutations.
\end{itemize}

\textbf{Étape 2 : Rédaction du résumé (Application élimination, généralisation, reformulation)}

\textit{« Le Ministre de l'Agriculture a lancé le Programme National de Mécanisation (PRONAMA), doté de 25 milliards FCFA pour acquérir du matériel agricole (tracteurs, charrues, semoirs) d'ici 2027. Face à la prédominance de l'outil manuel dans les exploitations familiales (80\% de la production), cause de faibles rendements, ce programme vise à moderniser l'agriculture. Il prévoit la création de 50 centres régionaux de services mécanisés pour la location subventionnée, la maintenance et la formation. La réussite conditionne une gestion rigoureuse du matériel et l'amélioration des pistes rurales pour l'évacuation des récoltes. »}

\textbf{Étape 3 : Comptage et Vérification}
\begin{itemize}
\item Mots comptés : \textbf{98 mots}. (Cible : 100 $\pm$ 10 $\rightarrow$ [90 ; 110]).
\item \textbf{Conforme.} Ratio respecté. Idées essentielles toutes présentes. Zéro copier-coller. Zéro opinion. Neutre.
\end{itemize}
\tcblower
\textbf{Solution détaillée :} Voir l'analyse ci-dessus. Le processus a consisté à : (1) Supprimer le nom du ministre, le lieu, la date, les noms des organisations paysannes, l'horizon 2035, les citations. (2) Généraliser « 500 tracteurs, 1000 charrues, 2000 semoirs » en « matériel agricole » (ou « équipement motorisé »). (3) Nominaliser « l'agriculture reste dominée par... » $\rightarrow$ « la prédominance de l'outil manuel » ; « explique les faibles rendements » $\rightarrow$ « cause de faibles rendements ». (4) Condenser « qui auront pour mission de louer... d'assurer... de former » $\rightarrow$ « pour la location..., la maintenance et la formation ». (5) Fusionner les conditions (gestion + pistes) en une phrase finale.
\end{exoresolu}

\courssub{Synthèse : Check-list de la reformulation réussie}

\begin{aretenir}[Check-list avant remise de copie]
\begin{enumerate}
\item \textbf{Fidélité :} Toutes les idées essentielles (arguments structurants) sont présentes ? Aucune idée secondaire (exemple, citation) n'a survécu ?
\item \textbf{Reformulation :} Aucune phrase copiée-collée ? Vocabulaire personnel (synonymie) ? Syntaxe propre (nominalisation, condensation) ?
\item \textbf{Neutralité :} Aucun « je », « à mon avis », « heureusement », adjectifs subjectifs ?
\item \textbf{Longueur :} Comptage effectué ? Dans l'intervalle $[0,9 \times N/4 \ ; \ 1,1 \times N/4]$ ?
\item \textbf{Cohérence :} Le résumé se lit-il comme un texte autonome, fluide, logiquement articulé (connecteurs conservés/adaptés) ?
\end{enumerate}
\end{aretenir}

\begin{savaistu}[Le résumé dans la vie professionnelle technique]
En BTS, DUT, ou sur le terrain (technicien supérieur, assistant ingénieur), la « note de synthèse » ou le « rapport d'activité » sont des contractions quotidiennes. Un chef de chantier ne lit pas les 50 pages du rapport journalier du géomètre ; il lit le résumé d'une page : « \textit{Terrassement achevé zone A. 2 pannes détectées sur canalisation DN200. Reprise prévue J+2. Besoin : 3 tuyaux + 1 brise-roche.} » C'est exactement la même compétence : \textbf{extraire le signal, éliminer le bruit, reformuler pour le décideur.}
\end{savaistu}

\begin{exercice}[Entraînement individuel]
Contractez le texte suivant (environ 360 mots) en un résumé de \textbf{90 mots $\pm$ 10} (soit entre 81 et 99 mots). Appliquez la méthode : 1. Thème/Thèse. 2. Idées essentielles (liste). 3. Élimination/Généralisation. 4. Reformulation (nominalisation). 5. Comptage.

\begin{itshape}« La Société Nationale des Hydrocarbures (SNH) a publié son rapport annuel 2023. Le chiffre d'affaires s'établit à 3 200 milliards de FCFA, en hausse de 12\% par rapport à 2022. Le Directeur Général, M. Adolphe Moudiki, attribue cette performance à la hausse du prix du baril sur le marché international, passé de 85 à 95 dollars en moyenne, et à l'optimisation de la production sur les champs matures grâce aux techniques de récupération assistée (injection d'eau, de gaz). 

Cependant, le rapport souligne des nuages à l'horizon. La production nationale continue de baisser, passant de 35 millions de barils en 2022 à 32 millions en 2023, du fait de l'épuisement naturel des gisements historiques comme Kole ou Ebomé. Aucun nouveau champ majeur n'est entré en production depuis 2018. Le taux de remplacement des réserves reste inférieur à 50\%. 

Pour inverser la tendance, la SNH mise sur l'offshore profond (projet Gaz du Sud, champ Bowle) et la diversification vers les énergies renouvelables (solaire, biomasse). Le DG appelle l'État à sécuriser le cadre fiscal pour attirer les investisseurs majors (TotalEnergies, Perenco, BW Energy) qui hésitent à engager les milliards nécessaires à l'exploration profonde. La transparence de la gestion des revenus pétroliers, via l'Initiative pour la Transparence dans les Industries Extractives (ITIE), reste un atout majeur pour la crédibilité du pays. »\end{itshape}
\end{exercice}

\begin{corrige}[Exercice n°1 : SNH Rapport 2023]
\textbf{Thème :} Bilan SNH 2023 / Avenir pétrolier Cameroun.
\textbf{Thèse :} Malgré un CA record 2023 (hausse prix + optimisation), la production baisse (épuisement gisements, faible remplacement réserves) ; l'avenir dépend de l'offshore profond et d'un cadre fiscal attractif pour les majors.

\textbf{Résumé proposé (92 mots) :}
\textit{« Le rapport annuel 2023 de la SNH affiche un chiffre d'affaires record de 3 200 milliards FCFA (+12\%), porté par la hausse du baril et l'optimisation des champs matures. Pourtant, la production recule (32 vs 35 millions de barils) du fait de l'épuisement des gisements historiques et d'un taux de remplacement des réserves inférieur à 50\% depuis 2018. La relance conditionne le développement de l'offshore profond (Gaz du Sud, Bowle) et la diversification énergétique, qui exigent un cadre fiscal sécurisé pour attirer les majors. La transparence ITIE demeure un atout de crédibilité. »}

\textbf{Vérification :} 92 mots $\in$ [81 ; 99]. Idées essentielles : CA record + causes ; Baisse prod + causes ; Solutions (offshore, diversification, fiscal) ; ITIE. Zéro citation, zéro nom propre secondaire (Moudiki, TotalEnergies...), généralisation « majors » pour investisseurs. Reformulation : nominalisation (« attribue cette performance à » $\rightarrow$ « porté par »), condensation.
\end{corrige}

\coursec{Rédiger, produire et améliorer le résumé}

Dans la section précédente, nous avons appris à \emph{reformuler} les idées essentielles d'un texte : remplacer les mots de l'auteur par nos propres mots, condenser les phrases, supprimer les exemples et les détails. Mais une liste de phrases reformulées, même correctes, ne fait pas encore un \cle{résumé}. Il reste à les organiser, à les relier, à les rédiger dans un texte continu et fluide, puis à relire et corriger ce texte jusqu'à ce qu'il soit irréprochable. C'est cette dernière étape du travail — la plus décisive à l'examen — que nous étudions maintenant : passer de la reformulation à la \emph{rédaction}, puis de la rédaction à l'\emph{amélioration}.

\courssub{Vue d'ensemble : les cinq étapes de production du résumé}

Rédiger un résumé n'est pas un acte unique mais un processus en cinq étapes enchaînées. L'élève qui saute une étape (par exemple qui rédige sans avoir repéré les idées, ou qui rend sa copie sans relecture) s'expose à des pertes de points considérables. L'organigramme ci-dessous fixe cette démarche, que vous devez pouvoir reproduire de mémoire le jour du Probatoire ou du Baccalauréat technique.

\begin{popfigure}
\begin{tikzpicture}[
  node distance=0.55cm,
  etape/.style={rectangle, rounded corners=3pt, draw=popblue, thick, fill=popblueL, text width=3.1cm, align=center, minimum height=0.9cm, font=\small},
  detail/.style={rectangle, draw=popmuted, fill=popcream, text width=5.6cm, align=left, font=\footnotesize},
  fleche/.style={-{Stealth[length=2.5mm]}, thick, popdark}
]
\node[etape, align=center] (e1){\textbf{1. LIRE}\\le texte deux fois};
\node[etape, below=of e1, align=center] (e2){\textbf{2. REPÉRER}\\les idées essentielles};
\node[etape, below=of e2, align=center] (e3){\textbf{3. REFORMULER}\\avec ses propres mots};
\node[etape, below=of e3, align=center] (e4){\textbf{4. RÉDIGER}\\un texte lié et compté};
\node[etape, below=of e4, align=center] (e5){\textbf{5. RELIRE}\\et améliorer};
\draw[fleche] (e1) -- (e2);
\draw[fleche] (e2) -- (e3);
\draw[fleche] (e3) -- (e4);
\draw[fleche] (e4) -- (e5);
\node[detail, right=0.6cm of e1] {Comprendre le sens global, dégager la thèse de l'auteur.};
\node[detail, right=0.6cm of e2] {Souligner les idées, supprimer exemples, répétitions, détails.};
\node[detail, right=0.6cm of e3] {Synonymes, changements de formulation, condensation.};
\node[detail, right=0.6cm of e4] {Phrases complètes, connecteurs logiques, respect du quart.};
\node[detail, right=0.6cm of e5] {Vérifier fidélité, cohésion, longueur, langue.};
\end{tikzpicture}
\legende{Organigramme des cinq étapes de production du résumé. Chaque étape prépare la suivante ; aucune ne peut être supprimée sans risque pour la note finale.}
\end{popfigure}

\begin{aretenir}[L'enchaînement des étapes]
Les trois premières étapes (lire, repérer, reformuler) ont été étudiées dans les sections précédentes. Cette section est consacrée aux deux dernières : \textbf{rédiger} le résumé et le \textbf{relire} pour l'améliorer. Au Baccalauréat technique, la rédaction et la relecture représentent souvent la moitié des points de l'épreuve de contraction.
\end{aretenir}

\courssub{Rédiger le résumé : du plan du texte au texte continu}

\textbf{L'intuition d'abord.} Imaginez que vous racontez à un ami le contenu d'un long article que vous venez de lire. Vous ne récitez pas des phrases isolées : vous enchaînez naturellement les idées (« d'abord... ensuite... c'est pourquoi... enfin... »). Le résumé fonctionne exactement ainsi : c'est un \emph{nouveau texte}, court, autonome, qui se lit tout seul, sans le texte de départ.

\begin{definition}[Le résumé rédigé]
Le \cle{résumé} est un texte bref (en général le \textbf{quart} du texte original, avec une tolérance de plus ou moins 10\,\%), rédigé en phrases complètes et liées, qui restitue \emph{fidèlement} et \emph{avec ses propres mots} les idées essentielles du texte de départ, sans commentaire personnel ni ajout.
\end{definition}

\textbf{Le plan du résumé est calqué sur le plan du texte.} C'est le principe directeur de la rédaction. Si le texte de départ développe trois idées (par exemple : les causes du chômage des jeunes, ses conséquences, les solutions proposées), le résumé comportera trois moments, dans le même ordre. On ne réorganise jamais la pensée de l'auteur : le résumé n'est pas une dissertation, il ne juge pas, il ne discute pas, il \emph{suit}.

\begin{methode}[Rédiger le résumé en cinq gestes]
\begin{enumerate}
\item \textbf{Suivre le plan du texte} : une idée essentielle du texte = une phrase (parfois deux) du résumé, dans le même ordre.
\item \textbf{Rédiger une phrase d'introduction} annonçant l'auteur, le titre et la thèse du texte.
\item \textbf{Rédiger en phrases complètes} (sujet, verbe, complément), jamais en simples mots ou fragments.
\item \textbf{Relier les phrases par des connecteurs logiques} adaptés au sens : d'abord / ensuite / enfin (énumération) ; car / en effet (cause) ; donc / c'est pourquoi (conséquence) ; mais / cependant (opposition).
\item \textbf{Compter les mots} au fur et à mesure et ajuster : supprimer ou condenser si l'on dépasse, développer une idée si l'on est trop court.
\end{enumerate}
\end{methode}

\textbf{La phrase d'introduction.} Le résumé s'ouvre par une phrase qui présente le texte résumé. Elle doit contenir trois informations : l'\emph{auteur}, le \emph{titre} (et éventuellement le genre ou la date) et la \emph{thèse} (l'idée directrice). Modèles de formulation :

\begin{itemize}
\item « Dans \emph{Le Vieux nègre et la médaille}, Ferdinand Oyono dénonce, à travers le personnage de Meka, l'hypocrisie de la colonisation. »
\item « Dans cet extrait d'\emph{Une si longue lettre}, Mariama Bâ évoque le poids des traditions dans la vie des femmes. »
\end{itemize}

\begin{attention}[Erreur fréquente : la formule répétitive]
Ne commencez pas chaque phrase du résumé par « Dans ce texte, l'auteur dit que... » ou « L'auteur affirme que... ». Cette répétition alourdit le texte et gaspille des mots. Variez les formulations : « Selon l'auteur... », « Il montre ensuite que... », « Il conclut en... », ou mieux encore, enchaînez directement les idées avec des connecteurs : « En effet... », « C'est pourquoi... », « Enfin... ». La formule de présentation de l'auteur n'est utile qu'\emph{une seule fois}, dans la phrase d'introduction.
\end{attention}

\textbf{Assurer la cohésion du résumé.} Un résumé cohérent repose sur trois outils travaillés dans les sections de langue :

\begin{itemize}
\item les \cle{connecteurs logiques}, qui explicitent les relations entre les idées (addition, cause, conséquence, opposition, chronologie) ;
\item les \cle{reprises nominales} (pronoms, synonymes, mots généraux), qui évitent les répétitions : « Meka... le vieillard... cet homme... il... » ;
\item la \cle{progression thématique} maîtrisée : chaque phrase reprend un élément déjà connu (le thème) pour ajouter une information nouvelle (le propos), ce qui fait avancer le texte sans rupture.
\end{itemize}

\begin{exemplebox}[Passage de notes reformulées à un résumé rédigé]
Notes reformulées (idées essentielles d'un texte sur l'école) : \emph{école = outil de promotion sociale ; mais inégalités d'accès en milieu rural ; solutions : construire des lycées techniques, former des enseignants.}

Résumé rédigé : « Dans cet article, l'auteur présente l'école comme le principal instrument de promotion sociale des jeunes. Cependant, il constate que l'accès à l'éducation reste inégal, notamment dans les zones rurales. C'est pourquoi il propose de construire davantage de lycées techniques et de former plus d'enseignants. »

On observe : une phrase d'introduction, un connecteur d'opposition (\emph{cependant}), un connecteur de conséquence (\emph{c'est pourquoi}), et une reprise nominale (\emph{il} pour \emph{l'auteur}).
\end{exemplebox}

\textbf{Le respect du nombre de mots.} La consigne d'examen impose presque toujours le \textbf{quart} du texte. Pour un texte de 400 mots, le résumé doit compter environ 100 mots, entre 90 et 110. Comptez mot à mot (les mots composés et les formes élidées comme \emph{l'}, \emph{d'} comptent pour un mot ; un sigle comme MINESEC compte pour un mot).

\begin{attention}[Erreur fréquente : le non-respect du nombre de mots]
Dépasser fortement le nombre de mots imposé (par exemple 150 mots au lieu de 100) ou rester très en dessous (60 mots) entraîne une \textbf{pénalité systématique} au Baccalauréat : le correcteur cesse de lire au-delà de la limite ou sanctionne le non-respect de la consigne. Un résumé trop long prouve que vous n'avez pas su supprimer ; un résumé trop court prouve que vous avez abandonné des idées essentielles. Comptez toujours vos mots et indiquez le total à la fin de votre copie : « (102 mots) ».
\end{attention}

\courssub{Relire et améliorer le résumé : les quatre critères du correcteur}

\textbf{L'intuition d'abord.} Un menuisier qui fabrique une chaise ne la livre pas sans vérifier qu'elle est stable, bien poncée et aux bonnes dimensions. De même, un résumé ne se rend jamais « premier jet ». La relecture est une étape de production à part entière, qui peut rapporter plusieurs points.

\begin{methode}[La relecture en quatre vérifications]
Relire son résumé, c'est se mettre à la place du correcteur et vérifier successivement quatre critères :
\begin{enumerate}
\item \textbf{La fidélité} : toutes les idées essentielles du texte sont-elles présentes ? Ai-je ajouté une idée qui n'est pas dans le texte, ou un jugement personnel ? Ai-je déformé la pensée de l'auteur ?
\item \textbf{La cohésion} : les phrases sont-elles reliées par des connecteurs et des reprises nominales ? Le texte se lit-il d'un trait, sans rupture ?
\item \textbf{La longueur} : le nombre de mots respecte-t-il le quart (plus ou moins 10\,\%) ? Le total est-il indiqué ?
\item \textbf{La langue} : orthographe, accords (sujet-verbe, adjectifs), conjugaison, ponctuation sont-ils corrects ? Les mots du texte ont-ils été réellement remplacés par mes propres mots ?
\end{enumerate}
\end{methode}

\begin{propriete}[Barème type de l'épreuve de résumé au Baccalauréat technique]
La correction du résumé s'appuie sur une grille critériée. Un barème fréquent, sur 20 points, est le suivant : fidélité et repérage des idées essentielles : 8 points ; reformulation effective (usage de ses propres mots) : 6 points ; cohésion et organisation du texte : 3 points ; correction de la langue : 3 points. Ce barème montre que la relecture (cohésion + langue) pèse à elle seule 6 points : négliger la relecture, c'est renoncer à près d'un tiers de la note.
\end{propriete}

La grille d'auto-évaluation ci-dessous reprend ces critères. Utilisez-la en classe et entraînez-vous à noter vos propres résumés avant de les rendre.

\begin{popfigure}
\begin{tikzpicture}
\draw[fill=popblue, draw=popblue] (0,0) rectangle (12.6,0.7);
\node[text=white, font=\bfseries\small, anchor=west] at (0.15,0.35) {Critère d'évaluation du résumé};
\node[text=white, font=\bfseries\small, anchor=east] at (12.45,0.35) {Barème /20};
\foreach \y/\txt/\pts/\coul in {
  -0.7/{Fidélité au sens : toutes les idées essentielles, sans ajout ni déformation}/{8 pts}/popblueL,
  -1.4/{Reformulation effective : phrases réécrites avec ses propres mots}/{6 pts}/popcream,
  -2.1/{Cohésion : connecteurs logiques, reprises nominales, progression claire}/{3 pts}/popblueL,
  -2.8/{Respect du quart : nombre de mots conforme et indiqué (dans la fidélité)}/{---}/popcream,
  -3.5/{Correction langagière : orthographe, grammaire, ponctuation}/{3 pts}/popblueL}{
  \draw[fill=\coul, draw=popline] (0,\y) rectangle (12.6,\y+0.7);
  \node[font=\footnotesize, anchor=west, text width=10.2cm] at (0.15,\y+0.35) {\txt};
  \node[font=\footnotesize\bfseries, anchor=east] at (12.45,\y+0.35) {\pts};
}
\end{tikzpicture}
\legende{Grille d'évaluation du résumé sur 20 points. Le respect du nombre de mots est vérifié à l'intérieur du critère de fidélité : tout écart excessif y est sanctionné.}
\end{popfigure}

\begin{exoresolu}[Corriger collectivement le résumé d'un élève]
Énoncé. Texte de départ (idées essentielles) : un article affirme que le sport scolaire développe la santé des élèves, qu'il améliore leur concentration en classe, mais qu'il reste négligé faute d'infrastructures ; l'auteur propose d'équiper les lycées. Consigne : résumé en 50 mots environ. Voici le résumé d'un élève (58 mots) :

\emph{« Dans ce texte, l'auteur dit que le sport est bon pour la santé. L'auteur dit aussi que le sport aide les élèves à mieux se concentrer. L'auteur dit enfin que le sport est négligé car il n'y a pas de terrains. L'auteur dit qu'il faut équiper les lycées, ce qui est une excellente idée. »}

Identifiez les erreurs et proposez une version améliorée.
\tcblower
\textbf{Analyse des erreurs (démarche de remédiation).}
\begin{enumerate}
\item \textbf{Répétition de la formule} « L'auteur dit » quatre fois : lourdeur et perte de mots.
\item \textbf{Ajout personnel} : « ce qui est une excellente idée » est un jugement de l'élève, interdit dans un résumé (atteinte à la fidélité).
\item \textbf{Cohésion faible} : aucun connecteur logique, les phrases sont juxtaposées.
\item \textbf{Longueur} : 58 mots, légèrement au-dessus de la cible, alors que des mots sont gaspillés en répétitions.
\end{enumerate}
\textbf{Version améliorée (47 mots)} : « Dans cet article, l'auteur affirme que le sport scolaire développe la santé des élèves et améliore leur concentration. Cependant, il constate que cette discipline reste négligée faute d'infrastructures. C'est pourquoi il propose d'équiper les lycées. »

Cette version est fidèle (les quatre idées, aucun ajout), cohérente (\emph{cependant}, \emph{c'est pourquoi}), reformulée (\emph{discipline}, \emph{constate}) et conforme à la longueur exigée.
\end{exoresolu}

\courssub{Le compte rendu de la démarche : justifier ses choix par écrit}

Après la production et la correction, le programme prévoit un \cle{compte rendu} : un court texte écrit dans lequel vous expliquez et justifiez votre démarche. Cet exercice développe la capacité officielle « rédiger et justifier une démarche, une réponse ou une conclusion ».

\begin{definition}[Le compte rendu de démarche]
Le compte rendu de démarche est un texte bref et argumenté dans lequel l'élève explique \emph{comment} il a construit son résumé : quelles idées il a conservées et pourquoi, quels passages il a supprimés (exemples, répétitions, détails) et pourquoi, comment il a reformulé et relié les idées.
\end{definition}

\begin{methode}[Rédiger un compte rendu de démarche en trois paragraphes]
\begin{enumerate}
\item \textbf{Les choix de conservation} : « J'ai conservé les trois idées essentielles du texte, à savoir..., car elles constituent la thèse et les arguments de l'auteur. »
\item \textbf{Les choix de suppression} : « J'ai supprimé l'exemple du paragraphe 2 et les répétitions du paragraphe 4, car ils illustrent ou reprennent des idées déjà exprimées sans ajouter d'information nouvelle. »
\item \textbf{Les choix de reformulation et de liaison} : « J'ai remplacé l'expression... par..., et j'ai relié les idées par les connecteurs... afin de marquer l'opposition puis la conséquence. »
\end{enumerate}
\end{methode}

\begin{exemplebox}[Extrait de compte rendu d'élève]
« Pour résumer l'article sur le sport scolaire, j'ai d'abord conservé les quatre idées essentielles : les bienfaits du sport sur la santé et sur la concentration, son abandon faute d'infrastructures et la solution proposée. J'ai ensuite supprimé les statistiques du deuxième paragraphe, qui ne faisaient qu'illustrer la première idée. Enfin, j'ai reformulé "il n'y a pas de terrains" par "faute d'infrastructures" et relié les idées par \emph{cependant} et \emph{c'est pourquoi} pour respecter la logique du texte. Mon résumé compte 47 mots, conformément à la consigne. »
\end{exemplebox}

\begin{exercice}[Production complète]
Lisez attentivement le texte suivant (environ 120 mots) : « Le téléphone portable a transformé la vie des jeunes Camerounais. Grâce à lui, ils communiquent instantanément avec leurs proches, même dans les villages les plus éloignés. Il facilite aussi l'accès à l'information et aux services bancaires mobiles, qui permettent d'envoyer de l'argent sans se déplacer. Cependant, cet outil comporte des dangers : de nombreux élèves y consacrent des heures au détriment de leurs études, et les réseaux sociaux les exposent parfois à de fausses informations. Il appartient donc aux parents et aux enseignants d'éduquer les jeunes à un usage raisonnable du téléphone. »

\begin{enumerate}
\item Rédigez un résumé de ce texte en 30 mots environ (plus ou moins 10\,\%).
\item Relisez votre résumé à l'aide de la grille d'évaluation et notez-vous sur 20.
\item Rédigez un compte rendu de démarche justifiant vos choix de suppression et de reformulation.
\end{enumerate}
\end{exercice}

\begin{corrige}[Exercice : production complète]
\textbf{1. Résumé possible (32 mots)} : « Dans ce texte, l'auteur montre que le téléphone portable facilite la communication et l'accès aux services. Cependant, il nuit aux études et expose aux fausses informations. Il préconise donc d'éduquer les jeunes à un usage raisonnable. » (Les quatre idées sont présentes, reliées par \emph{cependant} et \emph{donc}, sans ajout personnel.)

\textbf{2. Auto-évaluation attendue} : fidélité (8/8 si les quatre idées sont là), reformulation (6/6 si aucune phrase n'est copiée), cohésion (3/3 avec les connecteurs), langue (3/3 si aucune faute). Toute phrase copiée du texte fait perdre des points de reformulation.

\textbf{3. Compte rendu attendu} : l'élève indique qu'il a conservé les quatre idées essentielles (avantages, dangers, solution), supprimé les exemples et précisions (villages éloignés, services bancaires mobiles, réseaux sociaux), reformulé « au détriment de leurs études » par « il nuit aux études », et vérifié que son texte compte 32 mots, conforme au quart demandé.
\end{corrige}

\begin{aretenir}[L'essentiel de la section]
Rédiger un résumé, c'est suivre le plan du texte, ouvrir par une phrase d'introduction (auteur, titre, thèse), rédiger une phrase par idée essentielle, relier par des connecteurs et compter les mots. Améliorer un résumé, c'est le relire selon quatre critères : fidélité, reformulation, cohésion, langue. Le compte rendu de démarche justifie par écrit les choix de conservation, de suppression et de reformulation. Ces compétences, évaluées sur 20 points au Baccalauréat technique, s'acquièrent par l'entraînement régulier et l'auto-évaluation.
\end{aretenir}

\coursec{Le Vieux nègre et la médaille : contexte, exposé et bilan de l'œuvre}

Après avoir maîtrisé les techniques de la contraction de texte — identification des idées essentielles, reformulation, rédaction et amélioration du résumé — nous mobilisons désormais ces compétences pour achever l'étude de l'œuvre au programme : \emph{Le Vieux nègre et la médaille} de Ferdinand Oyono. Cette section propose une approche globale : contextualisation historique et littéraire, analyse thématique et stylistique (avec un focus sur la satire), préparation de l'exposé oral (exposé n° 3) et réalisation d'un bilan écrit sous forme de synthèse résumée.

\courssub{Ferdinand Oyono : vie, œuvre et engagement}

\begin{savaistu}[Un destin d'écrivain et d'homme d'État]
Ferdinand Oyono (1929--2010) écrit \emph{Le Vieux nègre et la médaille} en 1956, alors qu'il est étudiant en droit à Paris. Il n'a que 27 ans. Ce roman, son deuxième après \emph{Une vie de boy} (1956, prix Sainte-Beuve 1957), révèle une maturité politique et littéraire saisissante. Après l'indépendance du Cameroun (1er janvier 1960), il embrasse une carrière diplomatique de haut niveau : ambassadeur en France, au Royaume-Uni, aux Nations Unies, puis ministre des Affaires étrangères (1992--1997). Son parcours illustre la transition de la plume engagée contre le colonialisme à l'action au service de l'État camerounais souverain.
\end{savaistu}

\begin{definition}[Ferdinand Oyono (1929--2010)]
Écrivain et homme d'État camerounais, figure majeure de la littérature africaine d'expression française des années 1950. Auteur d'une trilogie romanesque fondatrice : \emph{Une vie de boy} (1956, prix Sainte-Beuve), \emph{Le Vieux nègre et la médaille} (1956) et \emph{Chemin d'Europe} (1960). Son œuvre dénonce avec une ironie féroce l'aliénation coloniale et l'illusion de l'assimilation.
\end{definition}

\courssub{Contexte historique : le Cameroun sous administration française (1916--1960)}

\begin{popfigure}
\begin{tikzpicture}[scale=0.9, every node/.style={font=\small}]
% Axe principal
\draw[popdark, thick, ->] (0,0) -- (15,0) node[right] {\textbf{Chronologie historique}};
% Graduations et dates
\foreach \x/\date/\event in {
  0/1916/Arrivée des Français (mandat SDN),
  2.5/1944/Conférence de Brazzaville,
  3.5/1945/Fin de la Seconde Guerre mondiale,
  5/1948/Création de l'UPC (Um Nyobé),
  7/1955/Insurrection nationaliste{,} répression,
  9/1956/\textbf{Publication du roman} (Loi-cadre Defferre),
  11/1958/Autonomie interne (Ahmadou Ahidjo),
  13/1960/\textbf{Indépendance} (1er janv.)
} {
    \draw[popdark, thick] (\x, -0.15) -- (\x, 0.15);
    \node[below, align=center, text width=2.2cm] at (\x, -0.5) {\date};
    \node[above, align=center, text width=2.5cm, popmuted] at (\x, 0.6) {\event};
}
% Zone roman
\draw[poporange, thick, dashed] (8.5, -1.2) rectangle (9.5, 1.5);
\node[poporange, rotate=90, anchor=south] at (8.3, 0.1) {\small Contexte d'écriture};
% Flèche Nord (indicative)
\node[popblue] at (14.5, 1.2) {\Large $\uparrow$ N};
\end{tikzpicture}
\legende{Frise chronologique : du mandat français (1916) à l'indépendance (1960). Le roman paraît en 1956, année charnière (Loi-cadre Defferre, début de l'autonomie), au cœur de la répression anti-UPC.}
\end{popfigure}

Pour comprendre la charge satirique du roman, il faut saisir la réalité du \textbf{code de l'indigénat} (travail forcé, impôt de capitation, justice sommaire, interdiction de circuler sans laissez-passer) et la politique d'\textbf{assimilation} : la promesse d'accéder à la citoyenneté française à condition de renoncer à sa culture. Le roman se situe au moment où cette promesse se révèle mensongère. L'\textbf{UPC} (Union des Populations du Cameroun), fondée en 1948 par Ruben Um Nyobé, revendique l'indépendance immédiate et la réunification. La répression brutale (1955--1960) fait des dizaines de milliers de morts. Oyono écrit dans l'urgence de ce basculement historique.

\courssub{Contexte littéraire : négritude et roman anticolonial}

\begin{definition}[La Négritude]
Mouvement littéraire et culturel né dans les années 1930 (Aimé Césaire, Léopold Sédar Senghor, Léon-Gontran Damas) qui revendique la valeur, la beauté et la dignité des civilisations noires face à l'idéologie coloniale de supériorité blanche. \cle{Refus de l'assimilation}, \cle{réhabilitation de l'Afrique précoloniale}, \cle{langage poétique surréalisant}.
\end{definition}

Si Oyono partage la révolte anticoloniale de la Négritude, il s'en distingue par le \textbf{genre} (le roman réaliste et satirique, non la poésie) et par le \textbf{ton} : l'ironie, le rire, le grotesque remplacent le lyrisme. \emph{Le Vieux nègre et la médaille} s'inscrit dans la vague du \textbf{roman anticolonial des années 1950} (\emph{L'Aventure ambiguë} de Cheikh Hamidou Kane, 1961 ; \emph{Les Soleils des Indépendances} d'Ahmadou Kourouma, 1968) qui analyse l'échec de la rencontre coloniale non pas comme un choc métaphysique, mais comme une escroquerie administrative et humaine.

\courssub{Résumé de l'intrigue : de l'illusion à la révolte}

\begin{exemplebox}[Résumé synthétique de l'intrigue]
Meka, vieux chef de village respecté, reçoit la visite du commandant de cercle qui lui annonce la remise de la \textbf{médaille de la Légion d'honneur} pour services rendus à la France (recrutement de porteurs, travail forcé, collaboration). Meka prépare une grande fête, s'endette, offre des cadeaux aux notables et aux missionnaires (Père Vandermayer). Le jour de la cérémonie, l'administrateur ne vient pas ; il envoie un sous-officier ivre qui épingle la médaille n'importe comment. Pire, on apprend que les terres du village sont confisquées pour une plantation européenne. Meka, humilié, rentre chez lui. Sa femme Kelara lui révèle qu'il a été trahi par ses propres "amis" (le catéchiste, le chef de poste). Dans un ultime sursaut de dignité, Meka arrache la médaille, la jette dans les latrines et décide de rejoindre la forêt (le maquis/UPC).
\end{exemplebox}

\courssub{Thématiques majeures}

\begin{center}
\begin{tabularx}{\linewidth}{|l|X|X|}\hline
\rowcolor{popblueL}
\textbf{Thème} & \textbf{Manifestation dans le roman} & \textbf{Portée critique} \\ \hline
\cle{Illusion de l'assimilation} & Meka croit que la médaille = reconnaissance/égalité. & Démonstration par l'absurde : l'assimilation est un leurre pour mieux exploiter. \\ \hline
\cle{Hypocrisie coloniale} & Discours du commandant ("enfants de la France") vs réalité (spoliation, mépris). & Déconstruction du "mythe civilisateur". \\ \hline
\cle{Hypocrisie religieuse} & Père Vandermayer bénit la médaille, ignore la spoliation, protège ses intérêts. & Critique de la complicité Église/État colonial. \\ \hline
\cle{Dignité retrouvée} & Geste final : jeter la médaille, partir au maquis. & Le refus de l'avilissement fonde l'humanité nouvelle. \\ \hline
\end{tabularx}
\end{center}

\courssub{La satire dans l'œuvre : mécanismes et cibles}

La satire est l'arme principale d'Oyono. Elle ne se contente pas de divertir ; elle \textbf{démasque} en grossissant le trait jusqu'à l'absurde.

\begin{propriete}[Procédés satiriques chez Oyono]
\begin{enumerate}
    \item \textbf{L'ironie narrative} : Le narrateur adopte le point de vue naïf de Meka pour mieux en montrer le décalage avec la réalité (focalisation interne ironique).
    \item \textbf{Le grotesque et la caricature} : Les administrateurs sont ridicules (ventrus, ivres, incompétents), les missionnaires cupides.
    \item \textbf{Le décalage langage/réalité} : Le style administratif pompeux ("services éminents") contraste avec la boue, la faim, la violence.
    \item \textbf{L'antiphrase} : Dire le contraire de ce qu'on pense. Ex : "La France généreuse" pour qualifier la puissance spoliantrice.
\end{enumerate}
\end{propriete}

\begin{exemplebox}[La médaille : symbole satirique central]
La médaille offerte à Meka symbolise la \textbf{récompense dérisoire de la soumission}, immédiatement dévalorisée par l'humiliation finale (remise par un sous-officier ivre, épinglée de travers). Elle condense tout le mensonge colonial : une breloque en toc pour une vie de servitude et la perte de sa terre.
\end{exemplebox}

\begin{popfigure}
\begin{tikzpicture}[scale=0.85, every node/.style={font=\small, align=center}]
% Nœuds personnages
\node[draw, rounded corners, fill=popblueL, thick, align=center] (Meka) at (0,0){\textbf{Meka}\\\emph{Vieux chef, naïf puis révolté}};
\node[draw, rounded corners, fill=poporangeL, thick, align=center] (Kelara) at (-3.5,-2.5){\textbf{Kelara}\\\emph{Femme lucide, voix de la vérité}};
\node[draw, rounded corners, fill=poppurpleL, thick, align=center] (Commandant) at (4,2){\textbf{Commandant}\\\emph{Administrateur hypocrite, absent}};
\node[draw, rounded corners, fill=popgreenL, thick, align=center] (Vandermayer) at (4,-2){\textbf{Père Vandermayer}\\\emph{Missionnaire complice, intéressé}};
\node[draw, rounded corners, fill=poppinkL, thick, align=center] (Engamba) at (-3.5,2){\textbf{Engamba}\\\emph{Interprète/mediateur, double jeu}};
% Relations
\draw[popdark, thick, ->] (Meka) -- node[above, sloped, pos=0.5] {Soumission initiale} (Commandant);
\draw[popdark, thick, ->] (Commandant) -- node[above, sloped, pos=0.5] {Promesse médaille / Spoliation} (Meka);
\draw[popdark, thick, ->] (Meka) -- node[left, pos=0.5] {Confiance / Écoute} (Kelara);
\draw[popdark, thick, ->] (Kelara) -- node[left, pos=0.5] {Révélation / Lucidité} (Meka);
\draw[popdark, thick, dashed, ->] (Meka) -- node[above, sloped, pos=0.5] {Collaboration religieuse} (Vandermayer);
\draw[popdark, thick, dashed, ->] (Vandermayer) -- node[below, sloped, pos=0.5] {Bénédiction intéressée} (Meka);
\draw[popdark, thick, dotted, ->] (Meka) -- node[above, sloped, pos=0.5] {Médiation linguistique} (Engamba);
\draw[popdark, thick, dotted, ->] (Engamba) -- node[below, sloped, pos=0.5] {Trahison / Manipulation} (Meka);
% Légende styles
\begin{scope}[xshift=9cm, yshift=-3.5cm]
\draw[popdark, thick, ->] (0,0.3) -- (1.5,0.3) node[right] {Relation principale (pouvoir)};
\draw[popdark, thick, dashed, ->] (0,0) -- (1.5,0) node[right] {Relation religieuse (complicité)};
\draw[popdark, thick, dotted, ->] (0,-0.3) -- (1.5,-0.3) node[right] {Médiation / Trahison};
\end{scope}
\end{tikzpicture}
\legende{Schéma des relations entre personnages principaux. Meka est au centre du système colonial (flèches pleines), manipulé par l'administration, l'Église et les intermédiaires. Kelara est le contre-pouvoir lucide.}
\end{popfigure}

\courssub{L'exposé n° 3 : présentation orale structurée}

L'exposé n° 3 évalue votre capacité à \textbf{présenter oralement un aspect précis de l'œuvre} (personnage, thème, procédé satirique) en vous appuyant sur un support écrit rigoureux. Ce n'est pas une lecture de résumé, mais une \textbf{analyse argumentée}.

\begin{methode}[Réussir l'exposé littéraire (Exposé n° 3)]
\begin{enumerate}
    \item \textbf{Choix du sujet ciblé} : Ne pas traiter "le roman" mais \emph{un} aspect : "La figure de Kelara, conscience lucide", "La satire du commandant de cercle", "Le symbolisme de la médaille", "L'échec de l'assimilation chez Meka".
    \item \textbf{Introduction écrite (sur support)} :
        \begin{itemize}
            \item Accroche (citation courte, fait historique, question).
            \item Sujet posé + définition des termes clés.
            \item Annonce du plan (2 ou 3 parties équilibrées).
        \end{itemize}
    \item \textbf{Développement illustré} :
        \begin{itemize}
            \item Chaque partie = 1 idée force + 1 citation précise (numéro de page/édition) + 1 analyse du procédé (ironie, lexique, syntaxe).
            \item Utiliser les connecteurs logiques (\emph{en outre, par conséquent, en revanche}).
        \end{itemize}
    \item \textbf{Conclusion personnelle argumentée} :
        \begin{itemize}
            \item Réponse synthétique à la problématique.
            \item Ouverture (autre œuvre d'Oyono, actualité, portée universelle).
        \end{itemize}
    \item \textbf{Support visuel (optionnel mais recommandé)} : Schéma, frise, citation projetée, carte mentale.
    \item \textbf{Prise de parole} : Regard sur l'auditoire, voix posée, temps maîtrisé (5--10 min), notes sur cartes (pas lecture intégrale).
\end{enumerate}
\end{methode}

\begin{attention}[Piège fréquent : le résumé déguisé]
\emph{Erreur :} Passer 80 \% du temps à raconter l'histoire de Meka ("Il fait la fête, puis il est triste, puis il part").
\emph{Correction :} L'intrigue ne sert que d'\textbf{appui} à l'analyse. Exemple : "La fête (chap. 4) est le théâtre de la satire : l'énumération des mets (\emph{‘vin de palme, poulets, manioc’}) contraste avec la misère des villageois, soulignant l'absurdité de l'hommage rendu à un oppresseur."
\end{attention}

\begin{exoresolu}[Sujet d'exposé : « La satire du Père Vandermayer »]
\textbf{Introduction} : Dans \emph{Le Vieux nègre et la médaille}, l'Église n'est pas un refuge mais un rouage du système colonial. Le Père Vandermayer incarne cette complicité. \textbf{Problématique} : Comment la satire dévoile-t-elle l'hypocrisie missionnaire ? \textbf{Plan} : 1. L'instrumentalisation de la religion. 2. La cupidité masquée en charité. 3. L'échec spirituel face à la révolte de Meka.

\textbf{Développement} :
1. \emph{Instrumentalisation} : Vandermayer bénit la médaille ("\emph{Dieu a béni la France}"). Citation p. 45 : "\emph{La médaille, c'est la croix du Christ sur la poitrine d'un noir}". \textbf{Analyse} : Antiphrase sacrilège ; confusion volontaire entre sacré et décorations républicaines pour légitimer le pouvoir colonial.
2. \emph{Cupidité} : Il réclame sa part du festin, s'inquiète pour sa plantation. Lexique du commerce ("\emph{affaire, bénéfice, compte}") dans la bouche d'un prêtre. \textbf{Analyse} : Grotesque ; le religieux devient commerçant.
3. \emph{Échec} : Face à Meka qui jette la médaille, Vandermayer ne trouve que des menaces ("\emph{Tu iras en enfer}"). Meka répond : "\emph{L'enfer, c'est ici avec vous}". \textbf{Analyse} : Renversement final ; le "sauvage" détient la vérité morale.

\textbf{Conclusion} : Vandermayer n'est pas un homme de foi trahi, mais un agent colonial conscient. La satire d'Oyono dépasse la caricature : elle accuse un système où la religion sert d'opium aux peuples (référence à Marx, reprise par la négritude politique). \textbf{Ouverture} : Comparer avec le Père Drumont dans \emph{Une vie de boy} (même fonction satirique).
\tcblower
\textbf{Évaluation (grille indicative)} : Sujet ciblé (2/2) -- Plan clair (2/2) -- Citations précises + analyse (6/6) -- Langage oral soigné (2/2) -- Temps respecté (2/2) -- Conclusion/ouverture (2/2) -- Support écrit remis (2/2) = 18/20.
\end{exoresolu}

\courssub{Bilan de l'étude : synthèse écrite et contraction finale}

Le bilan de l'unité prend la forme d'une \textbf{synthèse écrite récapitulative} qui mobilise explicitement la technique du résumé (contraction) apprise dans les sections 4 et 5. Vous devez produire un texte d'environ \textbf{200--250 mots} (soit environ 1/4 du texte source si on considère l'ensemble des notes de cours comme source) qui articule : contexte, intrigue, thèmes, satire, portée.

\begin{methode}[Rédiger le bilan synthétique (contraction appliquée à l'œuvre)]
\begin{enumerate}
    \item \textbf{Collecte des idées maîtresses} (relecture des fiches de cours) : Listez 8--10 idées forces (1 par séance/thème).
    \item \textbf{Hiérarchisation et regroupement} : Fusionnez les idées redondantes. Groupez par : Contexte/Genèse -- Intrigue/Structure -- Thèmes/Thèses -- Style/Satire -- Portée/Actualité.
    \item \textbf{Rédaction du "macro-texte"} : Rédigez un paragraphe par groupe en utilisant vos propres mots (reformulation), sans citation longue, avec des connecteurs de progression (\emph{Dans un premier temps, Par ailleurs, En définitive}).
    \item \textbf{Contrôle de la contraction} : Comptez les mots. Visez 200--250. Supprimez les exemples de détail, gardez les concepts. Vérifiez la cohérence (anaphores, connecteurs) et la correction (orthographe, syntaxe).
    \item \textbf{Titre} : Donnez un titre informatif (ex. \emph{Bilan d'étude : Le Vieux nègre et la médaille, satire de l'assimilation coloniale}).
\end{enumerate}
\end{methode}

\begin{exoresolu}[Bilan synthétique type (≈ 220 mots)]
\textbf{Bilan d'étude : Le Vieux nègre et la médaille, satire de l'assimilation coloniale}

Publié en 1956 au cœur de la guerre d'indépendance camerounaise, \emph{Le Vieux nègre et la médaille} de Ferdinand Oyono dénonce par la satire l'illusion de l'assimilation française. Le roman situe l'action dans un village sous administration coloniale où le code de l'indigénat régit le quotidien. Meka, vieux chef loyal, reçoit la promesse de la Légion d'honneur pour sa collaboration. La préparation de la fête mobilise la communauté et l'institution missionnaire, incarnée par le Père Vandermayer, complice intéressé. La cérémonie tourne à la farce : l'administrateur brille par son absence, un sous-officier ivre remet la médaille de travers, et l'on apprend la spoliation des terres villageoises. Humilié, Meka bascule : éclairé par sa femme Kelara, il rejette la breloque dans les latrines et rejoint le maquis.

L'œuvre structure sa critique par une ironie narrative constante : le point de vue naïf du héros révèle le cynisme des dominants. La satire cible l'hypocrisie administrative (discours paternaliste vs violence foncière), religieuse (bénédiction de la médaille, cupidité du prêtre) et les intermédiaires africains (Engamba). Le symbole de la médaille condense le mensonge colonial : récompense dérisoire pour une aliénation totale. La révolte finale de Meka, du rire à la dignité, inscrit le roman dans la lignée anticoloniale des années 1950, prolongée par l'engagement diplomatique de l'auteur. Oyono démontre que la liberté ne se mendie pas, elle se conquiert par le refus lucide de l'avilissement.
\tcblower
\textbf{Analyse de la contraction} : Le texte source (cours complet) représente ~3000 mots. Le bilan en fait ~220 (taux ~7 \%). Les idées principales (contexte, intrigue, thèmes, style, portée) sont toutes présentes. Les détails (noms des plats, dialogues exacts, description de la fête) sont éliminés. La progression logique (contexte $\to$ intrigue $\to$ analyse $\to$ portée) assure la cohérence.
\end{exoresolu}

\courssub{Exercices d'entraînement et remédiation}

\begin{exercice}[Entraînement à l'exposé n° 3]
Préparez le support écrit (introduction + plan détaillé avec 2 citations analysées par partie + conclusion) pour l'un des sujets suivants :
\begin{enumerate}
    \item « Kelara, figure de la résistance féminine et de la lucidité politique. »
    \item « Le corps du vieux nègre : de l'instrument de travail au lieu de la révolte. »
    \item « La fête du chapitre 4 : apogée et effondrement de l'illusion assimilationniste. »
\end{enumerate}
Durée de l'oral visée : 7 minutes.
\end{exercice}

\begin{exercice}[Contraction bilan : version examen]
En vous basant \textbf{uniquement} sur le bilan synthétique type ci-dessus (220 mots), produisez une \textbf{contraction extrême de 80 mots maximum} (résumé de résumé), conservant : auteur/date/contexte, intrigue noyau, thèse centrale, procédé dominant, portée.
\end{exercice}

\begin{corrige}[Exercice 2 -- Contraction 80 mots]
\textbf{Proposition (78 mots)} :
En 1956, au Cameroun en guerre, Ferdinand Oyono publie \emph{Le Vieux nègre et la médaille}. Meka, chef collaborateur, attend la Légion d'honneur. La cérémonie vire à la farce : administrateur absent, médaille remise ivre, terres spoliées. Humilié, Meka rejette la breloque et rejoint le maquis. Par l'ironie narrative et le grotesque, le roman satirise l'assimilation (leurre), la complicité Église/État et l'aliénation. La révolte finale du "vieux nègre" restaure sa dignité et annonce l'indépendance. Œuvre fondatrice de l'anticolonialisme réaliste.
\end{corrige}

\begin{aretenir}[Bilan des compétences validées dans cette unité]
\begin{itemize}
    \item \textbf{Connaissances} : Contexte historico-littéraire (Cameroun 1950, Négritude, Oyono), intrigue, thèmes, satire.
    \item \textbf{Savoir-faire réception} : Analyser les procédés satiriques (ironie, grotesque, antiphrase), caractériser le texte explicatif/argumentatif dans l'œuvre.
    \item \textbf{Savoir-faire production} : Préparer/présenter un exposé littéraire structuré (oral + support écrit). Rédiger une synthèse contractée (bilan) mobilisant les techniques de reformulation, hiérarchisation, cohérence.
    \item \textbf{Langue} : Identifier sigles/emprunts (UPC, ONU, Légion d'honneur), analyser cohésion/cohérence/progression dans les extraits étudiés.
\end{itemize}
\end{aretenir}

\newpage

\coursec{Activité d'intégration}

\courssub{Objectifs de l'activité}

Cette activité d'intégration clôt la leçon sur la contraction de texte. Elle place l'élève dans une situation professionnelle authentique, proche de celles qu'il rencontrera dans l'administration, l'entreprise ou le journalisme : produire une \cle{note de synthèse} à partir d'un dossier documentaire. À l'issue de l'activité, l'élève doit être capable de :

\begin{itemize}
\item repérer les \cle{idées essentielles} d'un texte explicatif et d'un discours officiel ;
\item éliminer les exemples, les répétitions et les détails secondaires ;
\item reformuler les idées essentielles avec ses propres mots, en respectant la \cle{fidélité} au texte de départ ;
\item rédiger un résumé \cle{cohésif} et \cle{cohérent}, en utilisant correctement les connecteurs et les reprises nominales ;
\item respecter la \cle{contrainte du quart} (ici : 150 mots maximum pour environ 600 mots de départ) ;
\item employer correctement les \cle{sigles} (MINESEC, MINEDUB) et les \cle{emprunts} courants du français administratif ;
\item justifier par écrit ses choix de suppression et de reformulation ;
\item s'auto-évaluer à l'aide d'une grille et participer à une correction collective de remédiation.
\end{itemize}

\courssub{Situation}

Yaoundé, région du Centre, fin septembre. Tu es stagiaire au cabinet du Délégué régional des Enseignements secondaires du Centre (DRESEC). Demain matin, le Délégué doit présenter, en cinq minutes, un point sur la rentrée scolaire devant le Gouverneur de région. Il te confie un dossier de presse composé de trois documents et te demande de rédiger une \cle{note de synthèse} de \textbf{150 mots maximum}, claire, fidèle et directement lisible à l'oral.

\textbf{Document 1 — Article de \emph{Cameroon Tribune} (environ 600 mots).} « Rentrée scolaire : le Centre affiche complet. » L'article annonce que la rentrée s'est déroulée « dans le calme » dans la région du Centre. Le MINESEC y signale un effectif global de 512 400 élèves dans les enseignements secondaires de la région, contre 489 700 l'année précédente, soit une hausse d'environ 4,6 \%. L'article cite longuement un proviseur de lycée de Yaoundé qui décrit les files d'attente devant son établissement, raconte l'histoire d'une élève de Première arrivée à 5 h du matin, et détaille la cérémonie de remise des manuels. Il mentionne aussi un problème : 37 salles de classe restent inachevées dans trois départements (Mfoundi, Méfou-et-Afamba, Nyong-et-Kellé), et 214 enseignants manquent encore à l'appel, principalement en mathématiques et en sciences physiques. L'article se termine par des anecdotes sur la circulation à Yaoundé le jour de la rentrée.

\textbf{Document 2 — Tableau statistique du MINESEC.} Effectifs d'élèves du secondaire par région (en milliers) :

\begin{center}
\begin{tabularx}{\linewidth}{|l|X|X|X|}\hline
\rowcolor{popblueL} \textbf{Région} & \textbf{Année N-1} & \textbf{Année N} & \textbf{Évolution} \\ \hline
Centre & 489,7 & 512,4 & + 4,6 \% \\ \hline
Littoral & 402,1 & 415,8 & + 3,4 \% \\ \hline
Ouest & 298,5 & 305,2 & + 2,2 \% \\ \hline
Extrême-Nord & 351,9 & 368,0 & + 4,6 \% \\ \hline
\end{tabularx}
\end{center}

\textbf{Document 3 — Extrait du discours officiel du MINEDUB.} Le ministre déclare : « L'École camerounaise reste le socle de notre développement. Le Gouvernement a recruté 2 000 enseignants du primaire et du secondaire cette année. J'exhorte les chefs d'établissement à veiller à l'assiduité des élèves et à la ponctualité des enseignants. Aucun enfant ne doit être exclu de l'école pour des raisons financières. »

\begin{attention}[Vérifier les chiffres avant de les recopier]
Un résumé doit être \cle{fidèle}, y compris dans les données chiffrées. Avant de retenir un pourcentage, vérifie-le : pour la région du Centre, la hausse se calcule ainsi :
\[ \frac{512,4 - 489,7}{489,7} \times 100 = \frac{22,7}{489,7} \times 100 \approx 4,6 \% \]
Le taux annoncé par l'article est donc exact. De même, pour l'Extrême-Nord : $(368,0 - 351,9) / 351,9 \times 100 \approx 4,6 \%$. Le Centre et l'Extrême-Nord affichent bien la même hausse, la plus forte du tableau. Recopier un chiffre faux ou déformé est une faute grave de fidélité, lourdement sanctionnée à l'examen.
\end{attention}

\courssub{Consignes}

Travail en groupes de quatre. Durée : deux séances.

\begin{enumerate}
\item \textbf{Lecture et repérage.} Lis attentivement les trois documents. Souligne dans chacun les idées essentielles : les faits, les chiffres, les décisions. Entoure les exemples, les anecdotes et les redondances qui pourront être supprimés.
\item \textbf{Classement.} Construis un tableau à deux colonnes : à gauche, les idées à conserver ; à droite, les passages à éliminer, avec la raison de chaque élimination (exemple, anecdote, répétition, détail hors sujet).
\item \textbf{Reformulation.} Reformule chaque idée essentielle avec tes propres mots, sans copier les phrases des documents. Vérifie que tu n'ajoutes aucune idée nouvelle et que tu ne déformes aucun chiffre.
\item \textbf{Rédaction.} Rédige la note de synthèse en \textbf{150 mots maximum} (soit environ le quart du document principal). Elle doit comporter une phrase d'introduction annonçant le sujet, un développement organisé (effectifs, réussites, difficultés, mesures) et une brève conclusion. Utilise au moins trois \cle{connecteurs} différents et écris correctement les sigles MINESEC et MINEDUB (majuscules, sans point).
\item \textbf{Justification.} Rédige un court paragraphe (5 à 8 lignes) justifiant deux suppressions et deux reformulations que tu as opérées.
\item \textbf{Auto-évaluation et remédiation.} Évalue ta production à l'aide de la grille ci-dessous, puis participe à la correction collective. La meilleure note sera lue en plénière comme modèle.
\end{enumerate}

\begin{center}
\begin{tabularx}{\linewidth}{|l|X|c|}\hline
\rowcolor{popblueL} \textbf{Critère} & \textbf{Attendu} & \textbf{Points} \\ \hline
Fidélité & Aucune idée ajoutée ni déformée ; chiffres exacts & /4 \\ \hline
Sélection & Idées essentielles conservées, exemples éliminés & /4 \\ \hline
Reformulation & Mots personnels, pas de phrases copiées & /4 \\ \hline
Cohésion & Connecteurs variés, reprises nominales correctes & /4 \\ \hline
Contrainte & 150 mots maximum ; sigles correctement écrits & /4 \\ \hline
\end{tabularx}
\end{center}

\courssub{Ressources à mobiliser}

\begin{aretenir}[Ce qu'il faut se rappeler]
\begin{itemize}
\item Un \cle{résumé} est fidèle (aucune idée ajoutée), bref (environ le \cle{quart} du texte de départ), reformulé (mots personnels) et cohésif (connecteurs, reprises nominales).
\item On supprime : les exemples, les anecdotes, les citations longues, les répétitions, les détails chiffrés secondaires. On conserve : la thèse, les arguments, les faits et chiffres essentiels.
\item Les \cle{connecteurs} utiles : d'abord, ensuite, en effet, cependant, toutefois, en outre, par conséquent, enfin.
\item Un \cle{sigle} s'écrit en majuscules, sans point : le MINESEC (Ministère des Enseignements secondaires), le MINEDUB (Ministère de l'Éducation de base). Un \cle{emprunt} est un mot venu d'une langue étrangère (par exemple \emph{planning}, \emph{budget}), fréquent dans le français administratif : il faut l'employer à bon escient et sans faute.
\item Le texte explicatif informe et explique de manière objective ; le résumé conserve ce ton neutre, sans commentaire personnel.
\end{itemize}
\end{aretenir}

\courssub{Démarche et corrigé commenté}

\begin{exoresolu}[Rédaction de la note de synthèse]
Énoncé : à partir des trois documents du dossier, rédiger une note de synthèse de 150 mots maximum pour le Délégué régional, puis justifier les choix opérés.
\tcblower
\textbf{Étape 1 — Repérage des idées essentielles.} Document 1 : rentrée calme au Centre ; 512 400 élèves contre 489 700 (+ 4,6 \%) ; 37 salles inachevées ; 214 enseignants manquants (mathématiques, sciences physiques). À éliminer : l'anecdote de l'élève de 5 h du matin, la cérémonie de remise des manuels, la circulation à Yaoundé (exemples et détails hors sujet). Document 2 : le Centre connaît la plus forte hausse, à égalité avec l'Extrême-Nord (+ 4,6 \%). Document 3 : recrutement de 2 000 enseignants ; exhortation à l'assiduité ; aucune exclusion financière.

\textbf{Étape 2 — Reformulation.} « 512 400 élèves » peut devenir « plus d'un demi-million d'élèves », mais le chiffre exact est ici jugé essentiel et conservé ; « manquent à l'appel » devient « font défaut » ; « affiche complet » devient « enregistre une forte fréquentation ».

\textbf{Étape 3 — Rédaction (modèle, 148 mots).}

\emph{Note de synthèse — Rentrée scolaire dans la région du Centre.}

\emph{La rentrée scolaire s'est déroulée dans le calme dans la région du Centre. D'abord, les effectifs progressent : selon le MINESEC, les enseignements secondaires accueillent 512 400 élèves, contre 489 700 l'an dernier, soit une hausse de 4,6 \%, la plus forte du pays à égalité avec l'Extrême-Nord. Cependant, des difficultés demeurent : 37 salles de classe restent inachevées dans trois départements et 214 enseignants font défaut, surtout en mathématiques et en sciences physiques. En outre, le MINEDUB a annoncé le recrutement de 2 000 enseignants et rappelé qu'aucun enfant ne doit être exclu de l'école pour des raisons financières. Enfin, les chefs d'établissement sont invités à veiller à l'assiduité des élèves et à la ponctualité des enseignants. La rentrée est donc globalement réussie, malgré des besoins persistants en infrastructures et en personnel.}

\textbf{Étape 4 — Justification des choix.} J'ai supprimé l'anecdote de l'élève arrivée à 5 h du matin car c'est un exemple qui illustre la forte affluence sans ajouter d'information. J'ai supprimé le passage sur la circulation à Yaoundé car il est hors sujet. J'ai reformulé « manquent à l'appel » par « font défaut » pour éviter la copie et adopter un registre administratif. J'ai condensé le tableau statistique en une seule comparaison (la plus forte hausse) afin de respecter la contrainte des 150 mots.

\textbf{Étape 5 — Vérification.} 148 mots (contrainte respectée) ; aucune idée ajoutée ; chiffres exacts et vérifiés ; sigles en majuscules sans point ; quatre connecteurs variés (d'abord, cependant, en outre, enfin) ; ton neutre et explicatif.
\end{exoresolu}

\begin{methode}[Les cinq étapes du résumé, à réutiliser à l'examen]
\begin{enumerate}
\item \textbf{Lire} le texte deux fois et dégager le thème et la thèse.
\item \textbf{Sélectionner} les idées essentielles ; barrer exemples, anecdotes, citations et répétitions.
\item \textbf{Reformuler} chaque idée avec ses propres mots, sans rien ajouter ni déformer.
\item \textbf{Rédiger} en respectant la contrainte de longueur (le quart), avec introduction, développement organisé et conclusion brève.
\item \textbf{Relire} : compter les mots, vérifier les chiffres, les sigles, les connecteurs et l'orthographe.
\end{enumerate}
\end{methode}

\courssub{Bilan de l'activité}

Cette activité a permis de mobiliser, dans une situation professionnelle réaliste, l'ensemble des savoir-faire de la leçon. Elle a d'abord montré que repérer les idées essentielles suppose de distinguer le fait de l'exemple : les anecdotes de l'article de presse, pourtant séduisantes, n'ont pas leur place dans une note administrative. Elle a ensuite confirmé que la reformulation n'est pas une paraphrase mot à mot, mais une réécriture fidèle, adaptée au destinataire et au registre de la situation de communication. Elle a enfin mis en évidence le rôle décisif de la \cle{cohésion} : sans connecteurs ni reprises nominales, le résumé devient une simple liste de phrases juxtaposées. La contrainte du quart, loin d'être arbitraire, oblige à hiérarchiser l'information, compétence indispensable dans la vie administrative et professionnelle. La grille d'auto-évaluation et la correction collective de remédiation ont permis à chaque élève d'identifier ses points faibles (fidélité, sélection, reformulation, cohésion ou respect de la contrainte) et de les corriger avant l'évaluation certificative du Probatoire et du Baccalauréat technique, où la contraction de texte est une épreuve régulièrement proposée.

\newpage

\coursec{Méthodes et exercices résolus}

\coursec{Contraction : Rédiger un résumé}

\courssub{Le texte explicatif : caractéristiques et tonalités}

\begin{methode}[Identifier le texte explicatif et ses tonalités]
    \textbf{Objectif :} Reconnaître un texte explicatif et repérer les tonalités satirique ou polémique.
    \begin{enumerate}
        \item \textbf{Repérer la finalité :} Le texte vise à \cle{informer}, \cle{expliquer} un phénomène, un mécanisme ou une idée. L'auteur s'efface derrière l'objectivité (sauf tonalité marquée).
        \item \textbf{Analyser la structure :} Introduction (sujet, thèse, plan) $\rightarrow$ Développement (arguments, exemples, données) $\rightarrow$ Conclusion (bilan, ouverture).
        \item \textbf{Identifier les marques linguistiques :}
            \begin{itemize}
                \item Temps : Présent de vérité générale, Futur (prévisions), Passé composé (faits avérés).
                \item Connecteurs logiques : \emph{donc, par conséquent, ainsi, en effet, or, mais}.
                \item Vocabulaire précis, termes techniques, définitions, reformulations (\emph{c'est-à-dire, autrement dit}).
            \end{itemize}
        \item \textbf{Détecter la tonalité satirique :} Ironie, humour, caricature, décalage entre le ton léger et la gravité du sujet $\rightarrow$ critique indirecte.
        \item \textbf{Détecter la tonalité polémique :} Attaque directe, arguments violents, réfutation de l'adversaire, vocabulaire chargé affectivement, usage de l'hyperbole, de l'accumulation.
    \end{enumerate}
    \textbf{Reflexe Bac :} Citer explicitement les indices textuels (mots, phrases, ponctuation) qui prouvent la tonalité.
\end{methode}

\begin{exoresolu}[Analyse d'un extrait explicatif à tonalité satirique]
    \textbf{Énoncé :} Lisez l'extrait suivant et répondez aux questions.
    \medskip
    \emph{« Nos honorables députés ont découvert, avec une stupéfaction feinte, que l'argent public avait des ailes. Ils ont donc voté, dans l'urgence et l'unanimité la plus touchante, une loi « anti-détournement » qui porte le chiffre 42. Le texte prévoit, tenez-vous bien, que tout fonctionnaire pris la main dans le sac devra... rembourser ce qu'il a volé. Une sanction terrible, assurément, qui fera trembler les plus endurcis. On murmure déjà que la prochaine session étudiera le projet de loi « anti-gravité » pour empêcher les valises de s'envoler toutes seules. »}
    \begin{enumerate}
        \item Quelles sont les caractéristiques du texte explicatif présentes ici ?
        \item Identifiez la tonalité employée et justifiez par trois procédés.
        \item Quel est l'effet visé sur le lecteur ?
    \end{enumerate}
    \tcblower
    \textbf{Solution détaillée :}
    \begin{enumerate}
        \item \textbf{Caractéristiques du texte explicatif :}
            \begin{itemize}
                \item \textbf{Sujet exposé :} La lutte contre la corruption / le vote d'une loi.
                \item \textbf{Structure informative :} Contexte (découverte), Action (vote loi n°42), Contenu de la loi (remboursement), Perspective (loi anti-gravité).
                \item \textbf{Marques linguistiques :} Présent de narration (\emph{ont découvert, ont voté, prévoit}), connecteurs (\emph{donc, tenez-vous bien, assurément}), vocabulaire du champ lexical juridique/administratif (\emph{députés, loi, fonctionnaire, sanction, session, projet de loi}).
            \end{itemize}
        \item \textbf{Tonalité : \cle{Satirique} (ironique).}
            \textbf{Procédés :}
            \begin{enumerate}
                \item \textbf{L'ironie / l'antiphrase :} \emph{« stupéfaction feinte »}, \emph{« unanimité la plus touchante »}, \emph{« sanction terrible »} (pour un simple remboursement), \emph{« fera trembler les plus endurcis »}. Le sens littéral est contredit par le contexte.
                \item \textbf{L'hyperbole / la caricature :} \emph{« loi anti-gravité »}, \emph{« valises [qui] s'envolent toutes seules »}. Réduction absurde du problème complexe de la corruption à un problème physique.
                \item \textbf{L'euphémisme / le sous-entendu :} \emph{« l'argent public avait des ailes »} (détournement), \emph{« pris la main dans le sac »} (flagrant délit). Le ton badin masque une accusation grave.
            \end{enumerate}
        \item \textbf{Effet visé :} Susciter le \cle{rire critique} et l'indignation. Le lecteur, complice de l'ironie, prend conscience de l'inefficacité réelle des mesures (simple remboursement = impunité) et de la complicité des législateurs (\emph{« honorables députés »} ironique). La satire dénonce l'hypocrisie du pouvoir.
    \end{enumerate}
    \textbf{Conclusion :} Ce texte explique le mécanisme de l'impunité législative sous couvert d'une fausse explication objective, utilisant la satire pour frapper les esprits.
\end{exoresolu}

\courssub{Cohésion, cohérence et progression du texte}

\begin{methode}[Analyser la cohésion, la cohérence et la progression]
    \textbf{Objectif :} Expliquer comment le texte « tient ensemble » (cohésion), « a du sens » (cohérence) et « avance » (progression).
    \begin{enumerate}
        \item \textbf{Cohésion (liens de surface, grammaticaux/lexicaux) :} Repérer les \cle{chaînes de référence} (pronoms, déictiques, reprises nominales, synonymes), les \cle{connecteurs} (logiques, temporels), les \cle{champs lexicaux} récurrents, l'\cle{anaphore} / \cle{cataphore}, l'\cle{ellipse}.
        \item \textbf{Cohérence (liens de fond, sémantiques) :} Vérifier la \cle{pertinence} du propos (thème unique), l'\cle{organisation logique} (cause/effet, thèse/antithèse, général/particulier), la \cle{non-contradiction}. Identifier le \cle{schéma textuel} (narratif, explicatif, argumentatif).
        \item \textbf{Progression (dynamique du texte) :} Classer le type de progression :
            \begin{itemize}
                \item \textbf{Linéaire (chaînée) :} Thème de la phrase $n$ = Rhème de la phrase $n-1$ (A $\rightarrow$ B, B $\rightarrow$ C).
                \item \textbf{Parallèle (à thème constant) :} Même thème de départ pour plusieurs phrases (A $\rightarrow$ B, A $\rightarrow$ C, A $\rightarrow$ D).
                \item \textbf{Progressive (à rhème constant / développement) :} Le rhème s'enrichit.
                \item \textbf{Par rebondissement / rupture :} Changement brutal de perspective (marqué par \emph{mais, or, cependant}).
            \end{itemize}
        \item \textbf{Synthèse :} Rédiger : « La cohésion est assurée par... La cohérence repose sur... La progression est de type... ce qui permet de... »
    \end{enumerate}
    \textbf{Reflexe Bac :} Ne pas confondre cohésion (forme) et cohérence (sens). Citer les mots-outils exacts.
\end{methode}

\begin{exoresolu}[Étude de la texture textuelle]
    \textbf{Énoncé :} Analysez la cohésion, la cohérence et la progression dans le paragraphe suivant :
    \medskip
    \emph{« [1] La déforestation en Amazonie s'accélère. [2] Elle menace directement la biodiversité unique de cette forêt. [3] En effet, la disparition des arbres prive les espèces animales de leur habitat naturel. [4] De plus, l'érosion des sols appauvrit la terre. [5] Par conséquent, les populations locales perdent leurs moyens de subsistance. [6] Cependant, des programmes de reforestation voient le jour. [7] Ils offrent un espoir fragile. »}
    \begin{enumerate}
        \item Relevez trois procédés de cohésion.
        \item Justifiez la cohérence de ce paragraphe.
        \item Identifiez le type de progression dominant et son évolution.
    \end{enumerate}
    \tcblower
    \textbf{Solution détaillée :}
    \begin{enumerate}
        \item \textbf{Procédés de cohésion :}
            \begin{enumerate}
                \item \textbf{Référence anaphorique (pronominal) :} \emph{« Elle »} (ph. 2) reprend \emph{« La déforestation »} (ph. 1) ; \emph{« Ils »} (ph. 7) reprend \emph{« des programmes de reforestation »} (ph. 6).
                \item \textbf{Reprise nominale / Synonymie :} \emph{« la disparition des arbres »} (ph. 3) = \emph{« La déforestation »} (ph. 1) ; \emph{« l'érosion des sols »} (ph. 4) conséquence lexicale de la disparition des arbres.
                \item \textbf{Connecteurs logiques :} \emph{« En effet »} (explication), \emph{« De plus »} (addition d'argument), \emph{« Par conséquent »} (conséquence), \emph{« Cependant »} (opposition/rupture).
            \end{enumerate}
        \item \textbf{Cohérence :} Le texte est \cle{cohérent} car il développe un unique \cle{thème} (conséquences de la déforestation) selon une logique \cle{causale} chaînée : Cause (déforestation) $\rightarrow$ Conséquences écologiques (biodiversité, sols) $\rightarrow$ Conséquences humaines (subsistance) $\rightarrow$ Solution partielle (reforestation). Aucune contradiction interne. Schéma textuel : \cle{Explicatif} (constat $\rightarrow$ causes/effets $\rightarrow$ perspective).
        \item \textbf{Progression :}
            \begin{itemize}
                \item Ph. 1 $\rightarrow$ 2 : \textbf{Linéaire} (Thème 1 = Rhème 1 : \emph{Déforestation} $\rightarrow$ \emph{Menace biodiversité}).
                \item Ph. 2 $\rightarrow$ 3 $\rightarrow$ 4 $\rightarrow$ 5 : \textbf{Parallèle (à thème constant)}. Le thème constant est \emph{« La déforestation / Elle / La disparition / L'érosion / Par conséquent »} (cause unique). Les rhèmes s'accumulent (biodiversité, habitat, sols, populations).
                \item Ph. 5 $\rightarrow$ 6 : \textbf{Rupture / Rebondissement} (marqué par \emph{« Cependant »}). Passage du constat négatif à l'action positive.
                \item Ph. 6 $\rightarrow$ 7 : \textbf{Linéaire} (\emph{Programmes} $\rightarrow$ \emph{Espoir}).
            \end{itemize}
            \textbf{Dominant :} Progression \textbf{parallèle} (accumulation des effets) suivie d'une \textbf{rupture} argumentative.
    \end{enumerate}
    \textbf{Conclusion :} La texture textuelle sert une démonstration cause/effet solide, brisée in fine par une ouverture contrastée.
\end{exoresolu}

\courssub{Le sigle et l'emprunt : outils lexicaux de la contraction}

\begin{methode}[Maîtriser le sigle et l'emprunt pour condenser]
    \textbf{Objectif :} Utiliser la siglaison et l'emprunt lexical comme stratégies de \cle{condensation} (gain de place) et de \cle{précision} technique.
    \begin{enumerate}
        \item \textbf{Le Sigle (Acronyme / Initialisme) :}
            \begin{itemize}
                \item \textbf{Formation :} Initiales des mots clés d'un syntagme (ex: \cle{ONU} = Organisation des Nations Unies). Prononcé lettre par lettre (initialisme) ou comme un mot (acronyme : \cle{UNESCO}, \cle{NASA}).
                \item \textbf{Règle de rédaction :} \textbf{1ère occurrence} = Forme développée + sigle entre parenthèses. \textbf{Occurrences suivantes} = Sigle seul.
                \item \textbf{Accord :} Le sigle s'accorde en genre/nombre selon le noyau du syntagme d'origine (ex: \emph{les ONG} [organisations], \emph{le GIP} [groupement]).
                \item \textbf{Utilité contraction :} Remplace un groupe nominal long par un mot court.
            \end{itemize}
        \item \textbf{L'Emprunt lexical :}
            \begin{itemize}
                \item \textbf{Définition :} Mot importé d'une langue étrangère (souvent l'anglais en technique/sciences) sans traduction (ex: \cle{software}, \cle{startup}, \cle{benchmark}, \cle{deadline}, \cle{management}).
                \item \textbf{Intégration :} Peut rester en italique (emprunt brut) ou s'adapter (prise de genre, pluriel français : \emph{des softwares} / \emph{des logiciels}).
                \item \textbf{Utilité contraction :} Désigne un concept complexe en un mot unique là où le français nécessiterait une périphrase (ex: \emph{cloud computing} vs \emph{informatique en nuage / externalisation des données sur serveurs distants}).
            \end{itemize}
        \item \textbf{Attention :} Ne pas abuser. En résumé, définir le sigle à la 1ère occurrence si le lecteur peut ne pas le connaître. Privilégier le terme français officiel s'il existe (Terminologie).
    \end{enumerate}
    \textbf{Formule à retenir :} \cle{Sigle} = Condensation morphologique (initiales). \cle{Emprunt} = Condensation sémantique (concept dense).
\end{methode}

\begin{exoresolu}[Application : Sigles et emprunts dans un résumé technique]
    \textbf{Énoncé :} Vous devez résumer un rapport technique de 500 mots en 80 mots. Le texte original contient : « L'Organisation Mondiale de la Santé (OMS) a publié un rapport sur la COVID-19. Le Directeur Général (DG) a insisté sur la nécessité d'un "lockdown" précoce. Les Pays à Revenu Intermédiaire (PRI) manquent de "ventilateurs". »
    \begin{enumerate}
        \item Réécrivez ce passage en appliquant les règles de la contraction (sigles, emprunts) pour un public averti.
        \item Justifiez le choix de garder ou non l'emprunt « lockdown » et « ventilateurs ».
    \end{enumerate}
    \tcblower
    \textbf{Solution détaillée :}
    \begin{enumerate}
        \item \textbf{Proposition de réécriture contractée (42 mots) :}
            \emph{« L'\textbf{OMS} alerte sur la \textbf{COVID-19}. Son \textbf{DG} préconise un \textbf{lockdown} précoce. Les \textbf{PRI} font face à une pénurie de \textbf{ventilateurs}. »}
            \begin{itemize}
                \item \textbf{OMS, DG, PRI, COVID-19} : Sigles/Acronymes standards. 1ère occurrence implicite car contexte « public averti / résumé technique » $\rightarrow$ gain de place maximal.
                \item \textbf{Lockdown} : Emprunt anglais conservé. Terme technique précis, passé dans l'usage scientifique/médiatique francophone (confinement strict). Plus court que « confinement strict et total ».
                \item \textbf{Ventilateurs} : Emprunt lexical intégré (anglais \emph{ventilator}), terme technique exact (respirateur artificiel). Précis et monosyllabique (gain de place vs « appareils d'assistance respiratoire »).
            \end{itemize}
        \item \textbf{Justification :}
            \begin{itemize}
                \item \textbf{Lockdown} : \textbf{Garder}. Public averti (technique/médical). Terme de référence international. Évite l'ambiguïté de « confinement » (qui peut être partiel). Gain de 2 mots.
                \item \textbf{Ventilateurs} : \textbf{Garder}. Terme technique consacré (Dictionnaire de l'Académie, terminologie hospitalière). « Respirateurs » est synonyme mais « ventilateurs » est l'emprunt technique standard pour la ventilation mécanique.
            \end{itemize}
            \textbf{Note :} Si le résumé visait le grand public, on écrirait : \emph{« L'Organisation Mondiale de la Santé (OMS)... son Directeur Général... un confinement strict... Les pays à revenus intermédiaires... des respirateurs artificiels. »}
    \end{enumerate}
    \textbf{Conclusion :} Le choix dépend du \cle{public cible} (expert vs profane) et de la \cle{contrainte de longueur}. L'expertise technique valide l'emprunt et le sigle nu.
\end{exoresolu}

\courssub{Reformuler les idées essentielles}

\begin{methode}[Réformuler : Sélectionner, Supprimer, Généraliser]
    \textbf{Objectif :} Passer du texte source au \cle{contenu} du résumé sans copier-coller (reformulation = changement de forme, conservation du sens).
    \begin{enumerate}
        \item \textbf{Étape 1 : Segmentation et Hiérarchisation (Macro-structure).}
            \begin{itemize}
                \item Découper en \cle{macro-propositions} (idées forces / paragraphes).
                \item Éliminer : \cle{Exemples}, \cle{Illustrations}, \cle{Détails} chiffrés (sauf si essentiels), \cle{Répétitions}, \cle{Digressions}, \cle{Commentaires méta-textuels} (\emph{« nous verrons que... »}).
            \end{itemize}
        \item \textbf{Étape 2 : Reformulation micro-textuelle (Opérations linguistiques).}
            \begin{itemize}
                \item \textbf{Généralisation / Hyperonymie :} Remplacer une liste (hyponymes) par la catégorie (hyperonyme). Ex: \emph{« lions, tigres, léopards »} $\rightarrow$ \emph{« félins »} / \emph{« grands prédateurs »}.
                \item \textbf{Synthèse / Fusion :} Grouper plusieurs phrases en une. Ex: \emph{« Il pleut. Le sol est mouillé. »} $\rightarrow$ \emph{« La pluie mouille le sol. »}
                \item \textbf{Transposition grammaticale :}
                    \begin{itemize}
                        \item Verbe $\rightarrow$ Nom (nominalisation) : \emph{« Il a décidé »} $\rightarrow$ \emph{« Sa décision »}.
                        \item Proposition subordonnée $\rightarrow$ Groupe prépositionnel / Participial : \emph{« Parce qu'il pleuvait, il est rentré »} $\rightarrow$ \emph{« Du fait de la pluie, il est rentré »} / \emph{« Il est rentré, la pluie le forçant »}.
                    \end{itemize}
                \item \textbf{Suppression du sujet agent / Passif / Impersonnel :} \emph{« Les ingénieurs ont construit le pont »} $\rightarrow$ \emph{« La construction du pont »} / \emph{« Le pont a été construit »}.
                \item \textbf{Préposition / Adverbe au lieu de conjonction :} \emph{« Bien qu'il soit tard »} $\rightarrow$ \emph{« Malgré l'heure tardive »}.
            \end{itemize}
        \item \textbf{Étape 3 : Vérification (Fidélité / Neutralité).} Aucun avis personnel, aucune info ajoutée, sens inversé ? Ordre logique respecté ?
    \end{enumerate}
    \textbf{Reflexe Bac :} Travailler sur brouillon : 1. Souligner idées forces. 2. Écrire mots-clés en marge. 3. Reconstruire phrases nouvelles.
\end{methode}

\begin{exoresolu}[Exercice de reformulation progressive]
    \textbf{Énoncé :} Reformulez le passage suivant en une seule phrase de 25 mots max, en appliquant : généralisation, nominalisation, suppression de l'agent.
    \medskip
    \emph{« Les techniciens de la SONARA ont procédé à la maintenance préventive de la turbine n°3. Ils ont démonté le rotor. Ils ont vérifié l'état des pales. Ils ont remplacé les joints d'étanchéité usés. Ils ont remonté l'ensemble. L'opération a duré 48 heures. Elle a permis d'éviter une panne critique. »}
    \tcblower
    \textbf{Solution détaillée :}
    \textbf{Analyse des idées forces :}
    \begin{itemize}
        \item Sujet/Agent : Techniciens SONARA (supprimable).
        \item Action principale : Maintenance préventive turbine n°3.
        \item Étapes (détails) : Démontage, vérification, remplacement joints, remontage $\rightarrow$ \textbf{Généralisation} : \emph{« révision complète »} ou \emph{« maintenance »}.
        \item Durée : 48h (détail chiffré, garder si contrainte temps, sinon supprimer).
        \item Finalité / Conséquence : Éviter panne critique (essentiel).
    \end{itemize}
    \textbf{Opérations de reformulation :}
    \begin{enumerate}
        \item \textbf{Nominalisation :} \emph{« ont procédé à la maintenance »} $\rightarrow$ \emph{« La maintenance préventive »}.
        \item \textbf{Généralisation des étapes :} \emph{« démonté, vérifié, remplacé, remonté »} $\rightarrow$ \emph{« comprenant révision et remplacement des joints »} ou simplement inclus dans \emph{« maintenance »}.
        \item \textbf{Suppression agent / Passif / Infinitif :} \emph{« La maintenance... a permis d'éviter... »}.
    \end{enumerate}
    \textbf{Propositions (22 mots) :}
    \emph{« La maintenance préventive de la turbine n°3, incluant le remplacement des joints d'étanchéité, a duré 48h et évité une panne critique. »}
    \textbf{Variante plus contractée (16 mots) :}
    \emph{« La révision préventive de la turbine n°3 (remplacement joints) a prévenu une panne critique en 48h. »}
    \textbf{Vérification :} Sens conservé ? Oui. Agent supprimé ? Oui. Généralisation (étapes $\rightarrow$ révision) ? Oui. Nominalisation ? Oui. Longueur respectée ? Oui.
\end{exoresolu}

\courssub{Rédiger, produire et améliorer le résumé}

\begin{methode}[Rédiger le résumé : Méthode en 5 temps (Bac)]
    \textbf{Objectif :} Produire un texte autonome, fidèle, concis, cohérent, respectant la contrainte de longueur (\emph{généralement 1/4 à 1/3 du texte original, ou nombre de mots imposé}).
    \begin{enumerate}
        \item \textbf{Temps 1 : Lecture active \& Repérage (15-20\%).} Lire 2 fois. Souligner : \cle{Sujet}, \cle{Thèse/Problématique}, \cle{Arguments principaux} (1 par §), \cle{Conclusion}. Chiffrer : mots source / mots cible.
        \item \textbf{Temps 2 : Prise de notes / Schéma (20\%).} Notez les \cle{mots-clés} (pas des phrases) dans l'ordre. Structure : \emph{Thèse $\rightarrow$ Arg 1 $\rightarrow$ Arg 2 $\rightarrow$ ... $\rightarrow$ Conclusion}. \textbf{Barrez} exemples/détails.
        \item \textbf{Temps 3 : Rédaction du brouillon (40\%).}
            \begin{itemize}
                \item \textbf{Phrase d'amorce (obligatoire) :} \emph{« Dans son texte « Titre » (Auteur, Date), X montre/démontre/analyse que [Thèse principale]. »}
                \item \textbf{Enchaînement :} Replier les notes en phrases complètes. Utiliser \cle{connecteurs} (\emph{Dans un premier temps, Ensuite, Par ailleurs, Enfin, Ainsi}). Appliquer \textbf{reformulation} (nominalisation, généralisation, passif/impersonnel).
                \item \textbf{Neutralité :} Pas de « je », pas d'adjectifs subjectifs (\emph{« intéressant », « choquant »}), temps au \textbf{présent} (ou passé si récit historique).
            \end{itemize}
        \item \textbf{Temps 4 : Comptage \& Ajustement (15\%).} Compter les mots.
            \begin{itemize}
                \item \textbf{Trop long :} Couper adjectifs/adverbes, fusionner phrases, généraliser listes, supprimer connecteurs lourds.
                \item \textbf{Trop court :} Développer un argument clé, ajouter une conséquence logique, préciser un chiffre clé.
            \end{itemize}
        \item \textbf{Temps 5 : Relecture finale (10\%).} \cle{Cohérence} (ça se comprend seul ?), \cle{Cohésion} (liens logiques), \cle{Orthographe/Grammaire}, \cle{Respect consigne} (mots, titre, auteur).
    \end{enumerate}
    \textbf{Formule d'amorce type :} \emph{« L'auteur, dans « Titre » (Année), soutient que [Thèse] en argumentant que [Arg1], [Arg2], pour conclure que [Conclusion]. »}
\end{methode}

\begin{exoresolu}[Production d'un résumé type Bac (Contrainte : 60 mots $\pm$ 10\%)]
    \textbf{Énoncé :} Résumez le texte suivant en \textbf{60 mots maximum} (texte source $\approx$ 220 mots).
    \medskip
    \emph{« L'intelligence artificielle (IA) transforme radicalement le secteur de la santé. D'abord, les algorithmes de diagnostic assisté par ordinateur (DAO) permettent une détection précoce des cancers avec une précision supérieure à l'œil humain. Ensuite, la chirurgie robotisée, pilotée à distance par des experts, offre une précision micrométrique, réduisant les séquelles post-opératoires. De plus, l'analyse massive de données (Big Data) facilite la médecine personnalisée, adaptant les traitements au profil génétique du patient. Cependant, ces avancées soulèvent des enjeux éthiques majeurs : la protection des données sensibles, la responsabilité en cas d'erreur algorithmique et le risque d'une médecine à deux vitesses, réservée aux pays riches. En conclusion, l'IA est un levier puissant pour la santé globale, à condition d'un encadrement juridique strict et d'une coopération internationale équitable. »}
    \tcblower
    \textbf{Solution détaillée :}
    \textbf{Temps 1-2 : Analyse \& Notes.}
    \begin{itemize}
        \item \textbf{Sujet/Thèse :} IA transforme santé (leviers + risques).
        \item \textbf{Arg 1 :} Diagnostic (DAO/IA) $\rightarrow$ détection précoce cancer + précision.
        \item \textbf{Arg 2 :} Chirurgie robotisée $\rightarrow$ précision, moins séquelles.
        \item \textbf{Arg 3 :} Big Data $\rightarrow$ médecine personnalisée (génétique).
        \item \textbf{Arg 4 (Cependant) :} Enjeux éthiques : données, responsabilité, fracture numérique (pays riches/pauvres).
        \item \textbf{Conclusion :} Levier puissant SI encadrement juridique + coopération.
    \end{itemize}
    \textbf{Temps 3 : Rédaction brouillon.}
    \emph{Dans son article, l'auteur montre que l'IA révolutionne la santé par le diagnostic précoce (cancers), la chirurgie robotique de précision et la médecine personnalisée via le Big Data. Toutefois, des défis éthiques persistent : protection des données, responsabilité algorithmique et inégalité d'accès Nord/Sud. L'IA reste un atout majeur pour la santé mondiale, nécessitant un cadre légal strict et une coopération équitable.}
    \textbf{Temps 4 : Comptage.} 62 mots. \textbf{Légèrement au-dessus (60 max). Ajustement :} Supprimer « de précision », « via le Big Data », « algorithmique », « mondiale ».
    \textbf{Version finale (58 mots) :}
    \begin{aretenir}[Résumé final - 58 mots]
        \emph{Dans son texte, l'auteur démontre que l'intelligence artificielle révolutionne la santé grâce au diagnostic précoce des cancers, à la chirurgie robotisée et à la médecine personnalisée. Cependant, elle pose des problèmes éthiques : protection des données, responsabilité en cas d'erreur et fracture d'accès entre pays riches et pauvres. L'IA est un levier puissant pour la santé, conditionné à un encadrement juridique strict et une coopération internationale équitable.}
    \end{aretenir}
    \textbf{Temps 5 : Vérification.} Fidèle ? Oui. Neutre ? Oui. Cohérent ? Oui. Longueur ? 58 $\in$ [54, 66]. \textbf{Validé.}
\end{exoresolu}

\courssub{Le Vieux nègre et la médaille : contexte, exposé et bilan}

\begin{methode}[Traiter l'œuvre au programme : Contexte, Exposé, Bilan]
    \textbf{Objectif :} Maîtriser les 3 temps de l'étude de \emph{Le Vieux nègre et la médaille} (Ferdinand Oyono) pour l'oral/écrit du Bac.
    \begin{enumerate}
        \item \textbf{Contexte (Inscription) :}
            \begin{itemize}
                \item \textbf{Auteur :} Ferdinand \cle{Oyono} (Cameroun, 1929-2010). Diplomate, homme politique. Figure majeure de la \cle{littérature négro-africaine d'expression française}.
                \item \textbf{Époque :} Publié \textbf{1956} (fin période coloniale, veille des indépendances). Climat : nationalisme, critique du colonialisme, \cle{négritude} (Senghor, Césaire) mais Oyono = \cle{réalisme critique / satire}.
                \item \textbf{Genre :} \cle{Roman court} / \cle{Nouvelle longue}. Ton : \cle{Satirique}, \cle{Polémique}, \cle{Réaliste}.
            \end{itemize}
        \item \textbf{Exposé n°3 (Structure narrative) :} Le roman suit \cle{Meka}, vieux chef traditionnel loyal à l'administration coloniale.
            \begin{itemize}
                \item \textbf{Situation initiale :} Meka attend la médaille (Mérite civique) pour services rendus (recrutement porteurs, travail forcé, impôts).
                \item \textbf{Élément déclencheur :} Cérémonie de remise. Discours du Commandant / Administrateur.
                \item \textbf{Péripéties / Ironie :} Meka doit payer la médaille, le drapeau, la fête. Humiliation : on lui tend la médaille comme à un enfant / un chien. Discours creux du blanc.
                \item \textbf{Dénouement :} Retour au village. Meka réalise l'absurdité. Il enterre la médaille (refus symbolique) ou la jette. Prise de conscience lucide (\cle{anagnorisis}).
                \item \textbf{Situation finale :} Rupture avec l'illusion coloniale. Début de la conscience politique.
            \end{itemize}
        \item \textbf{Bilan (Thèmes \& Portée) :}
            \begin{itemize}
                \item \textbf{Critique du colonialisme :} Aliénation, exploitation (travail forcé, impôts), hypocrisie de la « mission civilisatrice », mépris racial.
                \item \textbf{La médaille :} \cle{Symbole fort} = Leurres, chaîne dorée, prix de la dignité vendue, objet vide de sens.
                \item \textbf{Le « Vieux Nègre » :} Figure du \cle{collabo} lucide trop tard. Représente la génération des « intermédiaires » (chefs, interprètes) piégés.
                \item \textbf{Style :} \cle{Ironie narrative} (focalisation interne sur Meka $\rightarrow$ décalage avec réalité), \cle{oralité} (proverbes, rythme), \cle{langue française} appropriée, détournée.
            \end{itemize}
    \end{enumerate}
    \textbf{Reflexe Bac :} Citer des passages précis (cérémonie, enterrement, proverbes). Lier forme (satire) et fond (libération).
\end{methode}

\begin{exoresolu}[Question type Bac : La médaille, symbole de l'aliénation]
    \textbf{Énoncé :} « La médaille n'est pas une récompense, c'est un piège. » À partir de l'étude de \emph{Le Vieux nègre et la médaille}, commentez cette affirmation en montrant comment l'œuvre articule satire et prise de conscience.
    \tcblower
    \textbf{Solution détaillée :}
    \textbf{Introduction :} \emph{Le Vieux nègre et la médaille} (1956) de Ferdinand Oyono met en scène Meka, chef traditionnel décoré pour sa « loyauté ». L'affirmation résume le basculement du roman : l'objet honorifique devient le révélateur de l'oppression.
    \textbf{Développement :}
    \begin{enumerate}
        \item \textbf{La médaille, leurre de la « mission civilisatrice » (Satire) :}
            \begin{itemize}
                \item \textbf{Inversion des valeurs :} Meka a « servi » : recrutement porteurs (corvée), perception impôts (spoliation). La médaille récompense \textbf{sa propre soumission}.
                \item \textbf{La cérémonie grotesque :} Meka paie la médaille, le drapeau, le vin. Le discours du Commandant est vide (\emph{« fils dévoué »}). Ironie : le blanc parle de « mérite » pour de l'exploitation. Tonalité \cle{satirique} : décalage entre solennité officielle et réalité sordide (achat de la médaille).
            \end{itemize}
        \item \textbf{La médaille, chaîne dorée (Aliénation économique et symbolique) :}
            \begin{itemize}
                \item \textbf{Coût financier :} Meka s'endette (vente bétail) pour payer la fête imposée. L'honneur ruine le « honoré ». Allégorie de l'économie coloniale : le colonisé paie sa propre domination.
                \item \textbf{Coût symbolique :} Portée au cou comme un collier de chien (\emph{« on lui passait la médaille au cou »}). Perte de la verticalité (homme debout $\rightarrow$ bête courbée).
            \end{itemize}
        \item \textbf{Le refus final : Prise de conscience lucide (Anagnorisis) :}
            \begin{itemize}
                \item \textbf{Retour au village :} Silence de la foule. Regard des jeunes (mépris/compréhension). Meka « voit » enfin.
                \item \textbf{L'enterrement :} \emph{« Il enterra la médaille au pied du bananier. »} Acte fondateur. Rejet du symbole. Passage du statut de « Vieux nègre » (objet colonial) à \textbf{Sujet historique} (homme libre dans sa tête).
                \item \textbf{Proverbe final :} \emph{« Quand le vieux singe se regarde dans l'eau, il voit sa laideur. »} Lucidité tardive mais totale.
            \end{itemize}
    \end{enumerate}
    \textbf{Conclusion :} La médaille structure le roman : \textbf{objet de la satire} (dénonciation du système) $\rightarrow$ \textbf{objet du réveil} (déclencheur de la conscience). Oyono montre que la libération commence par le refus des symboles de l'aliénation. Œuvre charnière : de la plainte à la révolte consciente.
\end{exoresolu}

\courssub{Synthèse transversale : De la réception à la production}

\begin{methode}[Articuler lecture (Le Vieux Nègre) et écriture (Contraction)]
    \textbf{Objectif :} Transférer les compétences d'analyse (réception) vers la contraction (production).
    \begin{enumerate}
        \item \textbf{Repérer la progression narrative} (linéaire chronologique) $\rightarrow$ \textbf{Respecter l'ordre chronologique} dans le résumé.
        \item \textbf{Identifier la tonalité satirique} $\rightarrow$ \textbf{La restituer par le vocabulaire} (ironie $\rightarrow$ termes distanciés : \emph{« prétendue récompense », « soi-disant honneur »}) sans l'imiter servilement.
        \item \textbf{Extraire les thèmes majeurs} (Aliénation, Prise de conscience, Symbole) $\rightarrow$ \textbf{Les transformer en macro-propositions} du résumé.
        \item \textbf{Utiliser les sigles/emprunts du domaine littéraire} : \cle{Anagnorisis}, \cle{Focalisation}, \cle{Diégèse}, \cle{Intertextualité} (négritude) pour condenser l'analyse.
    \end{enumerate}
    \textbf{Exemple de phrase contractée :} \emph{« L'anagnorisis de Meka, déclenchée par l'absurdité satyrique de la cérémonie, signe le passage de l'aliénation consentie à la lucidité révoltée. »} (22 mots au lieu de 3 phrases).
\end{methode}

\begin{exoresolu}[Contraction d'un paragraphe critique sur l'œuvre]
    \textbf{Énoncé :} Contractez le paragraphe suivant en \textbf{35 mots} en conservant l'analyse littéraire.
    \medskip
    \emph{« Dans Le Vieux nègre et la médaille, Ferdinand Oyono utilise la focalisation interne sur Meka pour créer un effet d'ironie dramatique. Le lecteur comprend avant le personnage que la médaille est un leurre. Cette ironie prépare l'anagnorisis finale où Meka, lucide, enterre la médaille, refusant ainsi l'aliénation coloniale. Le style oralisé ancre le récit dans le réel camerounais. »}
    \tcblower
    \textbf{Solution détaillée :}
    \textbf{Idées forces :}
    1. Focalisation interne $\rightarrow$ Ironie dramatique (lecteur > personnage).
    2. Médaille = Leurre.
    3. Anagnorisis (enterrement) = Refus aliénation.
    4. Style oralisé = Ancrage réel Cameroun.
    \textbf{Reformulation / Fusion :}
    \emph{« Par la focalisation interne sur Meka, Oyono instaure une ironie dramatique dévoilant le leurre de la médaille ; l'anagnorisis finale (enterrement) scelle le refus de l'aliénation, dans un style oralisé ancrant le récit au Cameroun. »}
    \textbf{Comptage :} 33 mots. \textbf{Validé.}
    \textbf{Techniques utilisées :} Fusion phrases (participes \emph{instaure, dévoilant, scellant}), nominalisation (\emph{refus} au lieu de \emph{refuse}), vocabulaire précis (\emph{ironie dramatique, anagnorisis, focalisation}), suppression \emph{« Le lecteur comprend... »} (implicite dans ironie dramatique).
\end{exoresolu}

\begin{attention}[Les 5 erreurs qui coûtent des points au Bac]
    \begin{enumerate}
        \item \textbf{Le « Copier-Coller » déguisé (Plagiat / Absence de reformulation) :} Reproduire des phrases entières du texte source en changeant juste 2 mots. \textbf{\cle{Sanction :} Note 0/5 ou 1/5 en production.} $\rightarrow$ \textbf{Reflexe :} Fermer le texte source, écrire à partir de vos mots-clés.
        \item \textbf{Non-respect de la contrainte de longueur :} Dépasser de +10\% ou être trop court (-20\%). \textbf{\cle{Sanction :} Pénalité lourde (souvent 2 à 3 points sur la note globale).} $\rightarrow$ \textbf{Reflexe :} Compter \textbf{avant} la copie propre. Ajuster sur brouillon (fusion / généralisation).
        \item \textbf{Rupture de neutralité / Subjectivité :} Écrire \emph{« Je trouve que... », « C'est scandaleux que... », « L'auteur a raison... »}. Le résumé est un \cle{texte tiers}. \textbf{\cle{Sanction :} Perte points « Respect consignes / Registre ».} $\rightarrow$ \textbf{Reflexe :} Présent de narration, passif, impersonnel (\emph{« Il est montré que... »}).
        \item \textbf{Oubli de l'amorce (Cadre de référence) :} Ne pas citer \textbf{Auteur, Titre (italiques), Date/Genre, Thèse principale} en ouverture. \textbf{\cle{Sanction :} 1 à 2 points en moins (cohérence/exigence).} $\rightarrow$ \textbf{Modèle :} \emph{« Dans [Titre] (Auteur, Date), [Sujet] montre que [Thèse]... »}
        \item \textbf{Inversion / Désordre des idées (Non-respect de la progression) :} Mélanger conclusion et arguments, ou ordre chronologique. \textbf{\cle{Sanction :} Perte points « Cohérence / Fidélité ».} $\rightarrow$ \textbf{Reflexe :} Suivre scrupuleusement le plan du texte original (linéaire ou thématique).
    \end{enumerate}
\end{attention}

\newpage

\coursec{Exercices}

\courssub{Je vérifie mes connaissances}

\begin{exercice}[QCM : Types de progression thématique]
Pour chacun des paragraphes ci-dessous, cochez la case correspondant au type de progression thématique identifié (constante, linéaire, dérivée).

\begin{enumerate}
\item \textbf{Paragraphe A :} « Le téléphone mobile est devenu indispensable. \emph{Il} permet de communiquer. \emph{Cet appareil} offre aussi des jeux. \emph{Le smartphone} remplace l'ordinateur. »
\begin{itemize}
\item[$\square$] Progression constante \quad $\square$ Progression linéaire \quad $\square$ Progression dérivée
\end{itemize}

\item \textbf{Paragraphe B :} « La jeunesse camerounaise utilise le mobile. \emph{Cette utilisation} favorise l'accès à l'information. \emph{L'information} ainsi acquise développe l'esprit critique. \emph{L'esprit critique} renforce la citoyenneté. »
\begin{itemize}
\item[$\square$] Progression constante \quad $\square$ Progression linéaire \quad $\square$ Progression dérivée
\end{itemize}

\item \textbf{Paragraphe C :} « Les réseaux sociaux attirent les jeunes. \emph{Facebook} permet le partage. \emph{WhatsApp} facilite la messagerie. \emph{TikTok} valorise la vidéo courte. »
\begin{itemize}
\item[$\square$] Progression constante \quad $\square$ Progression linéaire \quad $\square$ Progression dérivée
\end{itemize}
\end{enumerate}
\end{exercice}

\begin{exercice}[Vrai / Faux justifié : Cohésion et cohérence]
Indiquez si les affirmations suivantes sont vraies (V) ou fausses (F) et justifiez brièvement votre réponse.

\begin{enumerate}
\item La cohésion textuelle repose uniquement sur l'enchaînement logique des idées. \hfill V / F
\item L'utilisation de connecteurs logiques (donc, cependant, par conséquent) relève de la cohérence. \hfill V / F
\item La reprise pronominale (il, ce, cela) et la chaîne nominale (le téléphone $\to$ cet appareil $\to$ le smartphone) sont des marques de cohésion. \hfill V / F
\item Un texte peut être cohésif sans être cohérent. \hfill V / F
\end{enumerate}
\end{exercice}

\begin{exercice}[Définitions et notions lexicales : Sigles et emprunts]
\begin{enumerate}
\item Donnez la définition précise du terme \cle{sigle} et citez deux exemples de sigles d'organismes publics camerounais (autres que ceux donnés en exemple dans le cours).
\item Expliquez la différence entre un \cle{sigle} et un \cle{acronyme}. Donnez un exemple d'acronyme courant.
\item Classez les emprunts suivants selon leur langue d'origine probable (anglais, allemand, arabe, italien, langues africaines) : \emph{leader, weekend, kindergarten, algorithme, safari, jazz, pizza, robot, toubab}.
\end{enumerate}
\end{exercice}

\begin{exercice}[Application directe : Reformulation par condensation]
Reformulez chacune des phrases suivantes en une seule proposition nominale ou verbale condensée, sans en changer le sens (nominalisation, suppression du superflu, synonymie).

\begin{enumerate}
\item Le fait que les jeunes utilisent massivement le téléphone mobile est un phénomène que l'on observe partout.
\item Il est nécessaire que les parents surveillent l'usage que font leurs enfants des réseaux sociaux.
\item Les autorités ont pris la décision d'interdire les téléphones dans les salles de classe.
\item C'est une évidence que l'accès à internet facilite la recherche documentaire pour les élèves.
\item Le coût élevé des forfaits data constitue un frein majeur pour la connexion des populations rurales.
\end{enumerate}
\end{exercice}

\courssub{Je m'entraîne}

\begin{exercice}[Repérage et hiérarchisation des idées essentielles]
Voici un texte explicatif d'environ 200 mots sur \emph{l'usage du téléphone mobile chez les jeunes Camerounais}. Lisez-le attentivement.

\textit{« Le téléphone mobile s'est imposé comme le compagnon inséparable de la jeunesse camerounaise. Dans les rues de Yaoundé comme dans les villages de l'Extrême-Nord, l'écran lumineux capte l'attention. Cet outil, jadis réservé aux élites, est devenu un bien de consommation courante grâce à la baisse des coûts des terminaux et à la concurrence entre les opérateurs. Les usages sont multiples. La communication vocale et la messagerie instantanée restent dominantes pour maintenir le lien familial et amical. Mais le mobile est aussi une porte ouverte sur le monde : réseaux sociaux, information en ligne, divertissement. Pour les élèves et étudiants, c'est un auxiliaire d'étude : dictionnaires, cours en ligne, révisions en groupe. L'envers du décor existe pourtant. La dépendance guette : anxiété, perturbation du sommeil, décrochage scolaire. Le cyberharcèlement et les arnaques en ligne sont des risques réels. Le coût des données, bien qu'en baisse, pèse encore sur le budget des familles modestes. Face à ces défis, l'éducation aux médias et à l'information s'impose. Parents, école et État doivent collaborer pour un usage responsable et citoyen du numérique. »}

\begin{enumerate}
\item Segmentez le texte en macro-structures (parties logiques) et donnez un titre à chaque partie.
\item Repérez les \textbf{idées essentielles} (maximum 5) en les numérotant. Justifiez chaque choix par un critère (thème principal, récurrence, place dans la structure).
\item Identifiez deux idées \textbf{secondaires} (exemples, détails, illustrations) qu'il faudra écarter du résumé.
\end{enumerate}
\end{exercice}

\begin{exercice}[Rédaction du résumé : Contrainte de forme]
À partir du travail de repérage effectué à l'exercice précédent, rédigez le \textbf{résumé} de ce texte en respectant scrupuleusement les consignes suivantes :
\begin{itemize}
\item Longueur : \textbf{100 mots $\pm$ 10 \%} (soit entre 90 et 110 mots).
\item Style : impersonnel, neutre, pas de « je », pas d'opinion personnelle.
\item Structure : phrases complètes, connecteurs logiques assurant la cohésion.
\item Fidélité : pas de contresens, pas d'ajout d'information extérieure.
\end{itemize}
Comptez vos mots et indiquez le total en bas de votre copie.
\end{exercice}

\begin{exercice}[Tonalités satirique et polémique : Le Vieux nègre et la médaille]
Relisez l'extrait de \emph{Le Vieux nègre et la médaille} (Ferdinand Oyono) correspondant à la \textbf{cérémonie de remise de la médaille} (chapitre consacré à la fête du 14 juillet : de l'arrivée du commandant au retour de Meka).

\begin{enumerate}
\item Relevez \textbf{quatre procédés} relevant de la \cle{tonalité satirique} (ironie, antiphrase, hyperbole, décalage, caricature, etc.) et citez le passage concerné.
\item Pour chaque procédé, expliquez la \textbf{portée critique} : quelle réalité coloniale est-elle dénoncée ? (ex. : la collaboration, l'assimilation illusoire, la violence symbolique, le mépris raciste).
\item En deux phrases, définissez la différence entre \cle{tonalité satirique} et \cle{tonalité polémique}. Le texte relève-t-il aussi de la polémique ? Justifiez.
\end{enumerate}
\end{exercice}

\begin{exercice}[Exposé : Introduction et annonce de plan]
Sujet : \emph{« L'inscription du Vieux nègre et la médaille dans le contexte de la négritude et de la littérature anticoloniale des années 1950. »}

Rédigez l'\textbf{introduction complète} d'un exposé oral de 10 minutes (environ 15 à 20 lignes écrites). Elle doit contenir :
\begin{enumerate}
\item Une \textbf{amorce} (accroche) contextualisant l'œuvre (auteur, date, genre, réception).
\item La \textbf{définition des termes clés} du sujet (négritude, littérature anticoloniale, engagement).
\item La \textbf{problématique} (question centrale à laquelle l'exposé répondra).
\item L'\textbf{annonce de plan} en trois parties équilibrées (ex. : Contexte historique et littéraire $\to$ Analyse thématique et stylistique $\to$ Portée et postérité).
\end{enumerate}
\end{exercice}

\courssub{Je me prépare au Bac}

\begin{exercice}[Sujet type Probatoire / Baccalauréat Technique : Contraction de texte (Coefficient 3)]
\textbf{Durée conseillée : 1 h 30}

\textbf{Texte support (environ 300 mots) :}

\begin{itshape}« L'intelligence artificielle (IA) transforme en profondeur les sociétés africaines. Au Cameroun, des jeunes entreprises innovantes développent des solutions locales : diagnostic agricole par analyse d'images, assistants conversationnels en langues nationales pour l'alphabétisation, prévision des crues des fleuves. L'État, à travers le ministère en charge des Postes et Télécommunications, impulse une stratégie nationale du numérique : formation aux métiers du secteur, cadre légal sur la protection des données, hébergement souverain des données nationales.

Pourtant, le fossé numérique persiste. Selon l'Union Internationale des Télécommunications (UIT), moins de la moitié de la population camerounaise avait accès à internet en 2023. L'électricité insuffisante, le coût du matériel et la faible littératie numérique excluent les zones rurales et les femmes. L'IA risque d'aggraver ces inégalités si elle reste l'apanage d'une élite urbaine connectée. De plus, l'importation massive de modèles entraînés sur des données étrangères menace l'identité culturelle et la souveraineté numérique du pays.

L'enjeu est donc double : innover en ancrant l'IA dans les réalités locales (données nationales, langues, savoirs traditionnels) et inclure par la formation massive, dans les écoles et les centres de formation professionnelle. L'éthique ne saurait être importée ; elle doit être pensée sur place, par des chercheurs, juristes et philosophes camerounais, pour que l'IA soit un levier de développement durable et non un nouvel instrument de dépendance. »\end{itshape}

\textbf{Consignes :}
\begin{enumerate}
\item \textbf{Analyse du texte (6 points) :}
\begin{enumerate}
\item Identifiez le \textbf{type de texte} et sa \textbf{nature} (argumentatif, explicatif, etc.). Justifiez par deux indices textuels.
\item Relevez \textbf{trois connecteurs logiques} différents et précisez pour chacun la relation logique qu'il marque (cause, conséquence, concession, opposition, addition, etc.).
\item Dégagez la \textbf{thèse principale} de l'auteur en une phrase.
\item Repérez \textbf{deux arguments principaux} développés pour soutenir cette thèse et, pour chacun, citez un exemple ou une illustration tirée du texte.
\end{enumerate}

\item \textbf{Contraction -- Résumé (10 points) :}
\begin{quote}
« Résumez le texte ci-dessus en \textbf{100 mots $\pm$ 10 \%} (90 à 110 mots). Votre résumé restituera fidèlement les idées essentielles, dans l'ordre du texte, sans opinion personnelle, en français correct, structuré par des connecteurs appropriés. Comptez vos mots et portez le total en marge. »
\end{quote}

\item \textbf{Langue et style (4 points) :}
\begin{enumerate}
\item Dans le texte, relevez \textbf{un sigle} (autre que IA) et \textbf{deux emprunts} à l'anglais éventuellement présents dans votre propre résumé ou dans le vocabulaire numérique courant. Pour le sigle, donnez son développement complet ; pour chaque emprunt, proposez un équivalent français.
\item Transformez la phrase suivante par \textbf{nominalisation} (tourner le verbe en nom) : « L'État impulse une stratégie nationale. » $\to$ « L'\emph{impulsion} par l'État d'une stratégie nationale... » (Reprenez le début et terminez la phrase nominalisée).
\item Réécrivez la phrase « L'IA risque d'aggraver ces inégalités si elle reste l'apanage d'une élite » en utilisant une \textbf{subordonnée concessive} (bien que / quoique + subjonctif).
\end{enumerate}
\end{enumerate}
\end{exercice}

\begin{exercice}[Sujet type Bac : Corpus et Question de littérature (Le Vieux nègre et la médaille)]
\textbf{Durée conseillée : 1 h}

\textbf{Corpus :} Extrait du chapitre de la cérémonie de la médaille — \textit{« Le commandant arriva... Meka tendit la main... La médaille brillait sur sa poitrine... Les discours furent longs... Le Vieux nègre rentra chez lui, la médaille au cou, mais le cœur vide. »} (Environ 25 lignes).

\textbf{Questions :}
\begin{enumerate}
\item \textbf{Contextualisation (3 points) :} Situez précisément cet extrait dans l'intrigue du roman (moment, lieu, personnages présents, enjeu immédiat pour Meka).
\item \textbf{Analyse stylistique (8 points) :} « La cérémonie de la médaille est le théâtre d'une satire féroce du système colonial. » À partir de l'extrait, \textbf{montrez comment l'écriture d'Oyono construit cette satire}. Vous organiserez votre réponse en \textbf{trois axes} distincts (ex. : La mise en scène grotesque du pouvoir / La déshumanisation du « bénéficiaire » / L'ironie dramatique et le décalage des discours). Pour chaque axe, citez le texte et commentez le procédé (vocabulaire, syntaxe, figures de style, tonalité).
\item \textbf{Portée critique (4 points) :} En quoi cette scène résume-t-elle la \textbf{critique de l'assimilation} et de la \textbf{collaboration} portées par le roman ? Répondez en 10 lignes maximum en articulant l'analyse du texte et votre connaissance de l'œuvre entière.
\item \textbf{Écriture d'invention (5 points) :} Imaginez le \textbf{monologue intérieur de Meka} sur le chemin du retour, cette nuit-là. Il rumine sa honte, sa colère sourde, sa lucidité nouvelle. Respectez le style d'Oyono (phrases courtes, images fortes, tonalité amère, vocabulaire du terroir). Longueur : 15 à 20 lignes.
\end{enumerate}
\end{exercice}

\begin{exercice}[Sujet type Bac : Production d'écrits -- Dissertation / Expression écrite]
\textbf{Durée conseillée : 1 h 30}

\textbf{Sujet :}
\textit{« Les nouvelles technologies (téléphone mobile, internet, intelligence artificielle) sont-elles un facteur de progrès ou un danger pour la jeunesse camerounaise ? »}

\textbf{Consignes :}
\begin{enumerate}
\item Vous traiterez ce sujet sous forme de \textbf{dissertation argumentée} (introduction, développement en deux ou trois parties équilibrées, conclusion).
\item Vous vous appuierez sur \textbf{vos connaissances} (textes étudiés en classe, actualité, observations personnelles) et sur \textbf{les documents du dossier} (textes de la leçon, statistiques, extraits d'articles de presse fournis ci-dessous).
\item La copie fera \textbf{300 mots minimum} (environ 1,5 à 2 pages). La qualité de la langue (syntaxe, vocabulaire, connecteurs, orthographe) sera notée sur 6 points.
\end{enumerate}

\textbf{Documents du dossier (extraits) :}

\textbf{Doc 1 (Statistiques indicatives, UIT 2023) :} Taux de pénétration d'internet au Cameroun : environ 40 \% de la population. Les jeunes de 15 à 24 ans sont nettement plus connectés que la moyenne ; les femmes et les zones rurales restent en retrait.

\textbf{Doc 2 (Témoignage d'élève, Lycée de Ngaoundéré) :} « Mon téléphone me permet de réviser mes cours sur WhatsApp avec mes camarades, de regarder des tutoriels de mathématiques sur YouTube. Mais je perds du temps sur TikTok, je dors moins. »

\textbf{Doc 3 (Extrait d'éditorial de presse, 2024) :} « L'IA offre une chance historique au Cameroun : saut technologique dans l'agriculture, la santé, l'éducation. À condition de former massivement, de protéger nos données et de refuser la dépendance algorithmique. »

\textbf{Doc 4 (Proverbe bamiléké) :} « L'outil ne fait pas l'artisan, c'est l'artisan qui fait l'outil. »
\end{exercice}



\newpage

\coursec{Corrigés des exercices}

\coursec{Corrigés des exercices — Leçon 04 : Contraction de texte}

\begin{corrige}[Exercice 1 — QCM : Types de progression thématique]
\begin{enumerate}
\item \textbf{Paragraphe A : progression constante.} Le même thème (le téléphone mobile) est repris d'une phrase à l'autre par un pronom (\emph{il}) et par des synonymes (\emph{cet appareil}, \emph{le smartphone}) : le thème reste identique, seule la formulation varie.
\item \textbf{Paragraphe B : progression linéaire.} Le rhème (l'information nouvelle) de chaque phrase devient le thème de la phrase suivante : \emph{utilisation} $\to$ \emph{information} $\to$ \emph{esprit critique}. Le texte avance « en chaîne ».
\item \textbf{Paragraphe C : progression dérivée.} Le thème général (\emph{les réseaux sociaux}) est éclaté en sous-thèmes particuliers (\emph{Facebook}, \emph{WhatsApp}, \emph{TikTok}) qui dérivent tous du thème initial.
\end{enumerate}
\textbf{Points de vigilance :} ne pas confondre progression constante (même thème repris) et progression dérivée (thème éclaté en sous-thèmes) ; la présence de synonymes ou de pronoms signe toujours la progression constante.
\end{corrige}

\begin{corrige}[Exercice 2 — Vrai / Faux justifié : Cohésion et cohérence]
\begin{enumerate}
\item \textbf{Faux.} L'enchaînement logique des idées relève de la \cle{cohérence} (le sens). La \cle{cohésion} repose sur les marques linguistiques de surface : reprises pronominales, chaînes nominales, connecteurs, substitutions lexicales.
\item \textbf{Faux.} Les connecteurs logiques sont des \emph{marques de cohésion} : ce sont des outils linguistiques visibles. C'est la relation logique qu'ils expriment (cause, opposition, conséquence) qui participe de la cohérence.
\item \textbf{Vrai.} Reprise pronominale et chaîne nominale sont les deux procédés majeurs de la cohésion référentielle : elles assurent la continuité du thème sans répétition lourde.
\item \textbf{Vrai.} Un texte peut aligner des phrases bien reliées grammaticalement (pronoms, connecteurs) mais dont les idées se contredisent ou n'ont aucun rapport logique : il est alors cohésif en surface mais incohérent en profondeur. Le résumé exige les deux.
\end{enumerate}
\textbf{Points de vigilance :} la justification est exigée ; une réponse « V » ou « F » seule ne rapporte aucun point. Retenir la formule : la cohésion se \emph{voit} (mots, pronoms, connecteurs), la cohérence se \emph{comprend} (logique des idées).
\end{corrige}

\begin{corrige}[Exercice 3 — Définitions et notions lexicales : Sigles et emprunts]
\begin{enumerate}
\item Un \cle{sigle} est une abréviation formée des initiales d'un groupe de mots et prononcée lettre par lettre. Exemples camerounais : \textbf{MINEDUB} (Ministère de l'Éducation de base), \textbf{ENAM} (École Nationale d'Administration et de Magistrature), \textbf{CRTV} (Cameroon Radio Television), \textbf{SONARA} (Société Nationale de Raffinage).
\item Le \cle{sigle} se prononce lettre par lettre (ex. : \textbf{C-R-T-V}) ; l'\cle{acronyme} se prononce comme un mot ordinaire. Exemples d'acronymes : \emph{SIDA} (Syndrome d'Immunodéficience Acquise), \emph{OVNI}.
\item Classement : \emph{leader, weekend, jazz} : anglais ; \emph{kindergarten} : allemand ; \emph{algorithme} : arabe (du nom du mathématicien Al-Khwarizmi) ; \emph{pizza} : italien ; \emph{robot} : tchèque (entré en français par l'anglais, pièce de Karel Čapek, 1920) ; \emph{safari} : swahili (langue africaine, via l'arabe \emph{safar}, voyage) ; \emph{toubab} : wolof (langue africaine).
\end{enumerate}
\textbf{Points de vigilance :} \emph{robot} est d'origine tchèque et non anglaise ; \emph{algorithme} vient de l'arabe et non du grec. En contraction, connaître sigles et emprunts permet de condenser (écrire « l'UIT » au lieu du nom complet) et de préférer l'équivalent français dans un résumé au style neutre.
\end{corrige}

\begin{corrige}[Exercice 4 — Application directe : Reformulation par condensation]
Propositions de reformulation (d'autres formulations correctes sont possibles) :
\begin{enumerate}
\item L'usage massif du téléphone mobile par les jeunes est un phénomène généralisé. (\emph{nominalisation} : « utilisent » $\to$ « l'usage » ; suppression de « le fait que ».)
\item Les parents doivent surveiller l'usage des réseaux sociaux par leurs enfants. (\emph{suppression de la tournure impersonnelle} « il est nécessaire que ».)
\item Les autorités ont décidé d'interdire les téléphones en classe. (\emph{condensation} : « ont pris la décision de » $\to$ « ont décidé de ».)
\item Internet facilite manifestement la recherche documentaire des élèves. (\emph{suppression de} « c'est une évidence que ».)
\item Le coût élevé des forfaits data freine la connexion des populations rurales. (\emph{verbalisation} : « constitue un frein majeur » $\to$ « freine ».)
\end{enumerate}
\textbf{Points de vigilance :} la reformulation ne doit jamais changer le sens (pas de contresens) ni ajouter d'idée ; supprimer le superflu (tournures impersonnelles, périphrases) sans supprimer l'information essentielle.
\end{corrige}

\begin{corrige}[Exercice 5 — Repérage et hiérarchisation des idées essentielles]
\begin{enumerate}
\item \textbf{Segmentation en quatre parties :}
\begin{itemize}
\item Partie 1 (phrases 1 à 3) : la généralisation du téléphone mobile chez les jeunes Camerounais.
\item Partie 2 (phrases 4 à 7) : la diversité des usages (communication, information, divertissement, études).
\item Partie 3 (phrases 8 à 11) : les dangers et limites (dépendance, risques en ligne, coût).
\item Partie 4 (phrases 12 et 13) : la solution préconisée : l'éducation aux médias et la collaboration des acteurs.
\end{itemize}
\item \textbf{Idées essentielles :}
\begin{enumerate}
\item Le mobile s'est généralisé chez les jeunes Camerounais (\emph{thème principal}, annoncé dès l'incipit).
\item Ses usages sont multiples : communication, information, divertissement, études (\emph{développement central}).
\item Il comporte des dangers : dépendance, cyberharcèlement, arnaques (\emph{retournement marqué par} « pourtant »).
\item Son coût reste un obstacle pour les familles modestes (\emph{nuance concessive}).
\item L'éducation aux médias et la collaboration parents-école-État s'imposent (\emph{conclusion du texte}, place finale stratégique).
\end{enumerate}
\item \textbf{Idées secondaires à écarter :} les exemples géographiques (« rues de Yaoundé », « villages de l'Extrême-Nord ») ; les illustrations concrètes des usages (dictionnaires, cours en ligne) ; les détails sur les symptômes de la dépendance (anxiété, sommeil).
\end{enumerate}
\textbf{Points de vigilance :} une idée est essentielle si sa suppression détruit le sens global ; un exemple, une répétition ou une précision chiffrée est presque toujours secondaire. Les marqueurs structurants (« pourtant », « face à ces défis ») signalent les articulations du texte : ils guident la segmentation.
\end{corrige}

\begin{corrige}[Exercice 6 — Rédaction du résumé : Contrainte de forme]
\textit{« Le téléphone mobile s'est généralisé chez les jeunes Camerounais, devenant un bien courant grâce à la baisse des prix. Ses usages sont variés : communication, information, divertissement et soutien scolaire. Cependant, cet outil présente des dangers réels : dépendance, cyberharcèlement et arnaques en ligne. En outre, le coût des données demeure lourd pour les familles modestes. C'est pourquoi l'éducation aux médias s'impose : parents, école et État doivent unir leurs efforts pour garantir un usage responsable et citoyen du numérique. »}

\textbf{Total : 78 mots.} Ce modèle est volontairement sous la barre pour montrer la marge de manœuvre ; une copie réelle devra atteindre 90 à 110 mots en développant légèrement chaque idée (par exemple en précisant les usages ou les risques). Vérifier : ordre du texte respecté, aucun « je », connecteurs variés (\emph{cependant}, \emph{en outre}, \emph{c'est pourquoi}), aucune information ajoutée.

\textbf{Barème indicatif (sur 10) :} respect du nombre de mots (2 pts) ; fidélité et exhaustivité des idées essentielles (4 pts) ; cohésion et connecteurs (2 pts) ; correction de la langue (2 pts).

\textbf{Points de vigilance :} tout dépassement ou insuffisance du nombre de mots est sanctionné ; le « je », les citations longues et les exemples du texte sont des fautes de méthode ; recopier des phrases entières du texte est assimilé à du plagiat, non à une contraction.
\end{corrige}

\begin{corrige}[Exercice 7 — Tonalités satirique et polémique : Le Vieux nègre et la médaille]
\begin{enumerate}
\item \textbf{Procédés satiriques attendus (exemples) :}
\begin{itemize}
\item \emph{L'hyperbole et la grandiloquence} : les discours officiels amplifient le « mérite » de Meka, présenté comme un modèle de fidélité, alors qu'il n'a fait que donner ses terres.
\item \emph{L'ironie et l'antiphrase} : la médaille, présentée comme une « récompense » glorieuse, est en réalité une babiole sans valeur qui paie la spoliation.
\item \emph{Le décalage} : le contraste entre la solennité de la cérémonie (uniformes, discours, fanfare) et la misère réelle de Meka (son pagne, sa vieillesse, sa pauvreté).
\item \emph{La caricature} : le portrait des notables blancs et des chefs africains ralliés, ridicules dans leur empressement à célébrer le colonisateur.
\end{itemize}
\item \textbf{Portée critique :} l'hyperbole des discours dénonce l'hypocrisie du discours colonial ; l'antiphrase de la « récompense » dénonce la spoliation foncière déguisée en bienfait ; le décalage dénonce le mépris raciste et l'assimilation illusoire ; la caricature dénonce la collaboration des élites africaines.
\item \textbf{Satire et polémique :} la \cle{satire} critique par le rire et la dérision (ironie, caricature) ; la \cle{polémique} attaque frontalement, par une argumentation ouverte et combative. Le roman relève d'abord de la satire, mais il touche à la polémique lorsque la lucidité de Meka, après son arrestation et son humiliation, se transforme en révolte explicite contre l'injustice coloniale.
\end{enumerate}
\textbf{Points de vigilance :} chaque procédé doit être appuyé par une citation courte et suivi de sa portée critique ; identifier un procédé sans l'expliquer ne vaut que la moitié des points.
\end{corrige}

\begin{corrige}[Exercice 8 — Exposé : Introduction et annonce de plan]
\textit{« En 1956, paraît à Paris un court roman qui dérange : Le Vieux nègre et la médaille, du Camerounais Ferdinand Oyono, déjà remarqué la même année pour Une vie de boy. L'œuvre est saluée par la critique et s'impose vite comme l'un des textes majeurs de la littérature africaine francophone de l'époque coloniale. Ce roman s'inscrit dans un double contexte. D'une part, la négritude, mouvement littéraire et culturel fondé par Senghor, Césaire et Damas, qui revendique la dignité et les valeurs de la civilisation noire face au racisme colonial. D'autre part, la littérature anticoloniale, ensemble d'œuvres engagées — c'est-à-dire qui prennent parti dans les débats de leur temps — dénonçant les méfaits de la colonisation : spoliation des terres, mépris raciste, assimilation illusoire. Dès lors, une question centrale se pose : en quoi Le Vieux nègre et la médaille illustre-t-il, par son intrigue et son écriture, les combats de la négritude et de la littérature anticoloniale des années 1950 ? Pour y répondre, nous situerons d'abord l'œuvre dans son contexte historique et littéraire ; nous analyserons ensuite ses thèmes et ses procédés satiriques ; nous mesurerons enfin sa portée critique et sa postérité. »}

\textbf{Barème indicatif :} amorce contextualisée (auteur, date, genre) : 1 pt ; définition des termes clés : 1 pt ; problématique claire : 1 pt ; annonce de plan en trois parties équilibrées : 1 pt ; qualité de la langue et ton soutenu adapté à l'oral : 1 pt.

\textbf{Points de vigilance :} vérifier la présence des quatre éléments obligatoires (amorce, définitions, problématique, annonce de plan) ; ne pas annoncer plus de trois parties ; à l'oral, employer le « nous » d'exposé et non le « je ».
\end{corrige}

\begin{corrige}[Exercice 9 — Sujet type Probatoire / Bac : Contraction de texte]
\textbf{1. Analyse du texte (6 points)}
\begin{enumerate}
\item \textbf{Type et nature (1,5 pt) :} texte \emph{explicatif} à visée \emph{argumentative}. Indices : présence d'une thèse (double enjeu : innover et inclure) ; connecteurs logiques organisant une démonstration ; exemples concrets étayant les idées.
\item \textbf{Connecteurs (1,5 pt) :} « Pourtant » (opposition) ; « De plus » (addition) ; « donc » (conséquence) ; « si » (condition/hypothèse). Trois connecteurs correctement identifiés avec leur relation logique suffisent.
\item \textbf{Thèse (1,5 pt) :} pour être un levier de développement durable, l'IA au Cameroun doit être ancrée dans les réalités locales et accessible à tous par la formation.
\item \textbf{Arguments (1,5 pt) :} (a) le fossé numérique exclut les zones rurales et les femmes (exemple : moins de la moitié de la population connectée en 2023) ; (b) l'importation de modèles étrangers menace l'identité culturelle et la souveraineté numérique (exemple : modèles entraînés sur des données étrangères).
\end{enumerate}

\textbf{2. Résumé (10 points) — modèle, 100 mots :} \textit{« L'intelligence artificielle transforme les sociétés africaines. Au Cameroun, des solutions locales apparaissent et l'État impulse une stratégie numérique nationale. Cependant, le fossé numérique persiste : une grande partie de la population, surtout rurale et féminine, reste exclue, et l'IA pourrait aggraver ces inégalités. En outre, l'importation de modèles étrangers menace l'identité culturelle et la souveraineté numérique. L'enjeu est donc double : innover en s'appuyant sur les réalités locales et inclure tous les citoyens grâce à la formation. Une éthique pensée sur place doit faire de l'IA un levier de développement durable, et non un instrument de dépendance. »} (100 mots.)

\textbf{Barème du résumé :} respect des 90 à 110 mots (2 pts) ; restitution des quatre idées essentielles dans l'ordre (4 pts) ; reformulation personnelle, sans recopie (2 pts) ; cohésion et correction de la langue (2 pts).

\textbf{3. Langue et style (4 points)}
\begin{enumerate}
\item Sigle : \textbf{UIT} = Union Internationale des Télécommunications. Emprunts à l'anglais du vocabulaire numérique : \emph{chatbot} $\to$ « assistant conversationnel » ; \emph{data} $\to$ « données » ; \emph{start-up} $\to$ « jeune entreprise innovante ».
\item Nominalisation : « L'impulsion par l'État d'une stratégie nationale du numérique favorise la formation et la protection des données. »
\item Subordonnée concessive : « Bien que l'IA soit un outil prometteur, elle risque d'aggraver ces inégalités si elle reste l'apanage d'une élite. » (subjonctif obligatoire après \emph{bien que} : \emph{soit}.)
\end{enumerate}

\textbf{Points de vigilance :} après \emph{bien que} et \emph{quoique}, le subjonctif est exigé (l'indicatif est une faute sanctionnée) ; la nominalisation doit conserver le sens exact de la phrase de départ ; dans le résumé, le total des mots doit figurer en marge sous peine de perdre le point de comptage.
\end{corrige}

\begin{corrige}[Exercice 10 — Sujet type Bac : Corpus et Question de littérature]
\begin{enumerate}
\item \textbf{Contextualisation (3 pts) :} l'extrait se situe au cœur du roman, lors de la fête du 14 juillet, dans la ville coloniale. Meka, vieux paysan qui a cédé ses terres à la mission, reçoit une médaille des mains du commandant, en présence des autorités coloniales, des notables et de la foule. Enjeu immédiat : Meka croit être honoré et intégré ; la scène prépare sa désillusion, qui éclatera après son arrestation nocturne. (Moment : 1 pt ; lieu et personnages : 1 pt ; enjeu pour Meka : 1 pt.)
\item \textbf{Analyse stylistique (8 pts) — attendus par axe :}
\begin{itemize}
\item \emph{Mise en scène grotesque du pouvoir :} accumulation des signes du faste (uniformes, discours interminables : « les discours furent longs ») ; l'hyperbole vide ridiculise l'autorité coloniale.
\item \emph{Déshumanisation du bénéficiaire :} Meka est réduit à un geste mécanique (« Meka tendit la main ») ; il est un objet de propagande, non une personne ; la médaille « brillait » : l'éclat de l'objet contraste avec l'effacement de l'homme.
\item \emph{Ironie dramatique et décalage :} le lecteur perçoit ce que Meka ignore encore ; la chute finale (« la médaille au cou, mais le cœur vide ») brise l'illusion par une antithèse cruelle.
\end{itemize}
(Barème : chaque axe = citation pertinente 1 pt + identification du procédé 0,5 pt + commentaire de la portée 1 pt ; organisation et langue 0,5 pt.)
\item \textbf{Portée critique (4 pts) :} la scène condense la dénonciation de l'assimilation : la médaille est le symbole dérisoire d'une intégration impossible, qui paie la spoliation des terres d'un hochet métallique. Elle dénonce aussi la collaboration : les notables africains qui applaudissent participent à leur propre aliénation. La suite du roman (arrestation, humiliation, rejet par les siens) révèle la vérité de cette cérémonie : le colonisateur ne reconnaît jamais le colonisé comme un égal.
\item \textbf{Écriture d'invention (5 pts) :} appréciée selon la fidélité au style d'Oyono (phrases courtes, images du terroir, amertume), la cohérence psychologique (honte, colère, lucidité) et la qualité de la langue.
\end{enumerate}
\textbf{Points de vigilance :} toute affirmation analytique doit être prouvée par une citation courte de l'extrait ; ne pas résumer l'intrigue à la place de l'analyse ; dans l'écriture d'invention, respecter la personne (« je » pour un monologue intérieur) et l'époque du récit.
\end{corrige}

\begin{corrige}[Exercice 11 — Sujet type Bac : Production d'écrits — Dissertation]
\textbf{Plan dialectique attendu (deux ou trois parties) :}
\begin{itemize}
\item \textbf{I. Les nouvelles technologies, facteur de progrès pour la jeunesse :} accès au savoir et aux études (Doc 2 : révisions, tutoriels) ; lien social et ouverture au monde ; opportunités économiques et innovation (Doc 3 : agriculture, santé, éducation).
\item \textbf{II. Mais un danger réel si l'usage est mal maîtrisé :} dépendance, perte de temps, troubles du sommeil (Doc 2) ; cyberharcèlement, arnaques, désinformation ; creusement des inégalités (Doc 1 : femmes et zones rurales en retrait).
\item \textbf{III. Synthèse : un progrès conditionné par l'éducation et la responsabilité :} éducation aux médias, formation massive, protection des données (Doc 3) ; l'outil dépend de l'usage qu'en fait l'utilisateur (Doc 4 : « L'outil ne fait pas l'artisan »).
\end{itemize}
\textbf{Introduction attendue :} amorce (généralisation du numérique chez les jeunes Camerounais), définition des termes (nouvelles technologies, progrès, danger), problématique (reformulée sous forme de question), annonce du plan.

\textbf{Barème indicatif (sur 20) :} introduction complète (3 pts) ; développement en parties équilibrées avec arguments et exemples (8 pts) ; mobilisation d'au moins trois documents du dossier (2 pts) ; conclusion (1 pt) ; qualité de la langue : syntaxe, vocabulaire, connecteurs, orthographe (6 pts).

\textbf{Points de vigilance :} respecter le minimum de 300 mots ; ne pas recopier les documents mais les exploiter (les citer brièvement ou les reformuler) ; éviter le plan « catalogue » sans articulation logique : chaque partie doit être reliée à la suivante par une transition ; la conclusion répond explicitement à la problématique sans introduire d'idée nouvelle.
\end{corrige}

\newpage

\coursec{Fiche bilan}

\begin{popfigure}
\begin{tikzpicture}[mindmap, grow cyclic, every node/.style={concept, text=white, font=\small},
  root concept/.append style={concept color=popblue, font=\bfseries\small, minimum size=3.2cm},
  level 1/.append style={level distance=3.6cm, sibling angle=60},
  level 1 concept/.append style={minimum size=2.6cm, font=\scriptsize}]
\node[root concept]{Contraction : rédiger un résumé}
  child[concept color=poporange]{node{Texte explicatif : neutralité, connecteurs, tonalités}}
  child[concept color=popgreen]{node{Cohésion et cohérence : reprises, connecteurs, thème unique}}
  child[concept color=poppurple]{node{Progression : thématique, analogique, en chaîne}}
  child[concept color=poppink]{node{Sigle et emprunt : économie lexicale}}
  child[concept color=popgold]{node{Reformuler : dégager, grouper, condenser}}
  child[concept color=popdark]{node{Rédiger et améliorer : brouillon, relecture, respect du quota}};
\end{tikzpicture}
\legende{Carte mentale de la leçon : les six compétences clés pour réussir la contraction de texte au Baccalauréat technique.}
\end{popfigure}

\begin{bilan}
\textbf{Définitions clés}
\begin{itemize}
  \item \cle{Résumé} : texte bref, rédigé avec ses propres mots, qui restitue fidèlement les idées essentielles d'un texte sans commentaire ni jugement personnel.
  \item \cle{Texte explicatif} : texte qui informe et fait comprendre ; il se caractérise par la neutralité, l'ordre logique, les connecteurs d'explication (\emph{car, en effet, c'est-à-dire}) et des définitions.
  \item \cle{Cohésion} : lien de surface entre les éléments du texte (reprises nominales et pronominales, connecteurs, substitutions).
  \item \cle{Cohérence} : unité de sens du texte (un seul thème, une progression logique, pas de contradiction).
  \item \cle{Sigle} : mot formé des initiales prononcées lettre par lettre (ONU, MINESEC) ; \cle{emprunt} : mot venant d'une autre langue (\emph{week-end}, \emph{business}).
\end{itemize}

\textbf{Les trois types de progression thématique}
\begin{center}
\begin{tabularx}{\linewidth}{|l|X|X|}
\hline
\rowcolor{popblueL} \textbf{Type} & \textbf{Principe} & \textbf{Exemple de reprise} \\
\hline
Linéaire (en chaîne) & le rème d'une phrase devient le thème de la suivante & \og la médaille\ldots{} Cette médaille\ldots\fg \\
\hline
À thème constant & le même thème est repris dans chaque phrase & \og Le vieux nègre\ldots{} Il\ldots{} Lui\ldots\fg \\
\hline
À thèmes dérivés & les thèmes dérivent d'un hyperthème général & Afrique $\rightarrow$ village, case, chef \\
\hline
\end{tabularx}
\end{center}

\textbf{Tonalités à reconnaître}
\begin{itemize}
  \item \cle{Tonalité satirique} : railler les défauts d'une personne ou d'une société (ironie, caricature, exagération) ; très présente dans \emph{Le Vieux nègre et la médaille} de Ferdinand Oyono.
  \item \cle{Tonalité polémique} : contester vivement une opinion ou un adversaire (questions rhétoriques, concession, dénonciation).
\end{itemize}

\textbf{Méthode express du résumé (4 étapes)}
\begin{enumerate}
  \item \textbf{Lire} deux fois le texte ; souligner les idées essentielles et repérer la structure.
  \item \textbf{Dégager} le plan : une idée principale par partie ; supprimer exemples, répétitions, citations, détails.
  \item \textbf{Reformuler} avec ses propres mots : condenser (nominalisation, sigles), grouper les idées voisines, relier par des connecteurs.
  \item \textbf{Rédiger} au présent de vérité générale, sans \og je \fg{}, puis compter les mots et corriger (orthographe, cohésion, quota imposé).
\end{enumerate}
\end{bilan}

\begin{aretenir}[Checklist avant le Bac]
\begin{itemize}
  \item $\square$ Je sais définir le résumé et le distinguer du compte rendu.
  \item $\square$ Je repère les caractéristiques d'un texte explicatif (neutralité, connecteurs, définitions).
  \item $\square$ Je distingue cohésion (liens de surface) et cohérence (unité de sens).
  \item $\square$ J'identifie les trois types de progression : linéaire, à thème constant, à thèmes dérivés.
  \item $\square$ Je reconnais les reprises nominales, pronominales et les connecteurs logiques.
  \item $\square$ Je sais ce qu'est un sigle et un emprunt, et je les utilise pour condenser.
  \item $\square$ Je repère les procédés de la tonalité satirique (ironie, caricature) et polémique.
  \item $\square$ Je dégage les idées essentielles et j'élimine exemples, répétitions et citations.
  \item $\square$ Je reformule avec mes propres mots, sans recopier de phrases du texte.
  \item $\square$ Je rédige sans avis personnel, sans \og je \fg{}, avec des connecteurs variés.
  \item $\square$ Je respecte le nombre de mots imposé et je relis pour corriger l'orthographe.
  \item $\square$ Je situe \emph{Le Vieux nègre et la médaille} dans son contexte (colonialisme, dérision de l'administration coloniale).
\end{itemize}
\end{aretenir}

\findecours

\end{document}
